% L'extraction des minerais — Allemand Tle A — Cours signé OnBuch+
% Programme officiel MINESEC. Compiler avec : tectonic AL12-l-extraction-des-minerais.tex
\newcommand{\DOCMATIERE}{Allemand}
\newcommand{\DOCNIVEAU}{Tle A}
\newcommand{\DOCMODULE}{Chapitre 3 : Environnement, santé et bien-être}
\newcommand{\DOCLECON}{AL12}
\newcommand{\DOCTITRE}{L'extraction des minerais}
% OnBuch+ — préambule « cours signés » ENRICHI (compilé avec tectonic / XeLaTeX)
% Système visuel « Pop » d'OnBuch+ : fond crème (jamais blanc), bordures encre
% épaisses, ombres « dures » décalées, titres Archivo Black en capitales.
% Version MATHÉMATIQUES : graphes (pgfplots/tikz), boîtes pédagogiques
% (définition, propriété, méthode, attention, le savais-tu, exercice résolu,
% exercice, TP, bilan), page de garde et signature OnBuch+ ancrée en bas.
%
% Macros attendues dans main.tex AVANT \input{preamble} :
%   \DOCMATIERE  \DOCNIVEAU  \DOCMODULE  \DOCLECON (numéro)  \DOCTITRE
\documentclass[11pt]{article}

\usepackage{fontspec}
\usepackage[ngerman,french]{babel} % français = langue principale ; l'allemand via \all{..} / textebox
\usepackage[a4paper,margin=2cm,top=3.4cm,bottom=2.8cm,headheight=1.4cm,footskip=1.3cm]{geometry}
\usepackage{amsmath,amssymb,amsfonts}
\usepackage{mathtools} % \xmapsto, \coloneqq... souvent utilisés par les modèles
\usepackage{cancel} % simplifications barrées (bilans de forces)
% \abs{z} (module, valeur absolue) et \norm{u} (norme), utilisés par les modèles
\providecommand{\abs}{}\let\abs\relax\DeclarePairedDelimiter{\abs}{\lvert}{\rvert}
\providecommand{\norm}{}\let\norm\relax\DeclarePairedDelimiter{\norm}{\lVert}{\rVert}
\usepackage[table]{xcolor} % [table] : fournit aussi \rowcolors (alternance de lignes)
\usepackage{pagecolor}
\usepackage{tikz}
\usetikzlibrary{arrows.meta,positioning,calc,shapes.geometric,shapes.misc,shapes.arrows,decorations.pathmorphing,decorations.pathreplacing,decorations.markings,patterns,shadows,mindmap,trees,fit,backgrounds,matrix,intersections,angles,quotes}
\usepackage{pgfplots}
\usepackage[version=4]{mhchem} % équations nucléaires \ce{^{235}_{92}U}
\usepackage[european,siunitx]{circuitikz} % schémas électriques
\pgfplotsset{compat=1.18}
\usepgfplotslibrary{fillbetween,groupplots} % aires entre courbes (calcul intégral)
\usepackage{enumitem}
\usepackage{fancyhdr}
% [most] charge toutes les bibliothèques tcolorbox (listings, minted, raster,
% documentation...) et gonfle inutilement la mémoire TeX sur un document long
% (~300 boîtes) : on ne charge que ce qui est réellement utilisé.
\usepackage[skins,breakable]{tcolorbox}
\usepackage{graphicx}
\usepackage{colortbl}
\usepackage{booktabs}
\usepackage{tabularx}
\usepackage{diagbox} % case d'en-tête diagonale des tableaux à double entrée
% Colonnes extensibles centrée / gauche / droite (C, L, R), souvent utilisées par les modèles
\newcolumntype{C}{>{\centering\arraybackslash}X}
\newcolumntype{L}{>{\raggedright\arraybackslash}X}
\newcolumntype{R}{>{\raggedleft\arraybackslash}X}
\usepackage{array}
\usepackage{multicol}
\usepackage{multirow}
\usepackage{siunitx}
% parse-numbers=false : accepte aussi \SI{2,5 \times 10^{-3}}{...}
\sisetup{locale=FR,output-decimal-marker={,},group-minimum-digits=4,parse-numbers=false}
\DeclareSIUnit\ohm{\ensuremath{\symup{\Omega}}} % U+2126 absent de la police
\DeclareSIUnit\division{div} % réglages d'oscilloscope (ms/div, V/div)
\DeclareSIUnit\tr{tr} % tours (vitesse de rotation, ex. \SI{1500}{\tr\per\minute})
\DeclareSIUnit\molar{M} % concentration molaire (ex. \SI{3}{\molar}, \SI{1}{\micro\molar})
\DeclareSIUnit\an{an} % année (le modèle confond parfois avec \year, un compteur TeX) — ex. \SI{5}{\kilo\meter\per\an}
\DeclareSIUnit\FCFA{FCFA} % franc CFA (ex. \SI{500}{\FCFA\per\kilo\gram})
\DeclareSIUnit\degreeBrix{\textdegree Brix} % teneur en sucre d'un sirop/jus (réfractométrie)
\DeclareSIUnit\month{mois} % le modèle utilise parfois \month, un compteur TeX, comme unité de durée
\usepackage{qrcode}
% \degree : commande fréquemment utilisée par les modèles pour un angle ou une
% température en dehors de \SI{}{} (non fournie par les paquets chargés).
\providecommand{\degree}{^{\circ}}
% \qed (fin de démonstration), utilisé par les modèles sans amsthm
\providecommand{\qed}{\hfill$\square$}
% \barre : double barre des valeurs interdites dans un tableau de variations (tkz-tab)
\providecommand{\barre}{\ensuremath{\|}}
% environnement proof (amsthm non chargé) utilisé par les modèles
\ifdefined\proof\else\newenvironment{proof}[1][Démonstration]{\par\noindent\textit{#1.}\ }{\qed\par}\fi
\usepackage{lastpage}
\usepackage{hyperref}
\hypersetup{hidelinks}

% ============================================================
% Palette « Pop » OnBuch+ — mêmes hex que la classe P (plus_ui.dart)
% ============================================================
\definecolor{poporange}{HTML}{F59321}
\definecolor{popdark}{HTML}{E07A0C}
\definecolor{popcream}{HTML}{FFF6E8}
\definecolor{popink}{HTML}{1C1714}
\definecolor{poppurple}{HTML}{7A5AE0}
\definecolor{popgreen}{HTML}{2E9E6A}
\definecolor{poppink}{HTML}{E0457B}
\definecolor{popblue}{HTML}{2F7DE1}
\definecolor{popgold}{HTML}{C98A12}
\definecolor{popmuted}{HTML}{8A7F76}
\definecolor{popline}{HTML}{EADFD2}
\definecolor{popgray}{HTML}{8A7F76}
\colorlet{popgrey}{popgray}
% Alias tolérants (le modèle nomme parfois les couleurs différemment)
\colorlet{popwhite}{white}
\colorlet{popblack}{popink}
\colorlet{popred}{poppink}
\colorlet{popyellow}{popgold}
% Teintes claires (fonds de tableaux/encadrés) — le modèle nomme parfois
% les couleurs avec un suffixe L (ex. popblueL) sans les avoir définies.
\colorlet{poporangeL}{poporange!20}
\colorlet{popdarkL}{popdark!20}
\colorlet{poppurpleL}{poppurple!20}
\colorlet{popgreenL}{popgreen!20}
\colorlet{poppinkL}{poppink!20}
\colorlet{popblueL}{popblue!20}
\colorlet{popgoldL}{popgold!20}
\colorlet{popredL}{popred!20}
\colorlet{popgrayL}{popgray!20}
% Teintes claires pour fonds de tableaux / zones
\colorlet{popblueL}{popblue!10!white}
\colorlet{poporangeL}{poporange!14!white}
\colorlet{popgreenL}{popgreen!12!white}
\colorlet{poppurpleL}{poppurple!12!white}
\colorlet{poppinkL}{poppink!12!white}
\colorlet{popgoldL}{popgold!14!white}

\pagecolor{popcream}

% ============================================================
% Typographie
% ============================================================
\defaultfontfeatures{Ligatures=TeX}
\setmainfont{PlusJakartaSans-Medium.ttf}[Path=../../../fonts/,BoldFont=PlusJakartaSans-Bold.ttf]
% Maths en sans-serif (Fira Math) avec chiffres et lettres droites de Plus
% Jakarta, pour que les formules se fondent dans le texte.
\usepackage[math-style=ISO,bold-style=ISO]{unicode-math}
\setmathfont{FiraMath-Regular.otf}[Path=../../../fonts/]
\setmathfont{PlusJakartaSans-Medium.ttf}[Path=../../../fonts/,range={up/num,up/{Latin,latin},\mathup}]
\usepackage{icomma} % virgule décimale sans espace en mode math
% Caractères mathématiques Unicode absents des polices de texte (Plus Jakarta,
% Archivo Black) : rendus par leur équivalent mathématique, en texte comme en
% formule — sinon ils s'affichent en carré vide (ex. « Arithmétique dans ℤ »).
\usepackage{newunicodechar}
\newunicodechar{ℝ}{\ensuremath{\mathbb{R}}}
\newunicodechar{ℤ}{\ensuremath{\mathbb{Z}}}
\newunicodechar{ℂ}{\ensuremath{\mathbb{C}}}
\newunicodechar{ℕ}{\ensuremath{\mathbb{N}}}
\newunicodechar{ℚ}{\ensuremath{\mathbb{Q}}}
\newunicodechar{𝔻}{\ensuremath{\mathbb{D}}}
\newunicodechar{α}{\ensuremath{\alpha}}
\newunicodechar{β}{\ensuremath{\beta}}
\newunicodechar{γ}{\ensuremath{\gamma}}
\newunicodechar{δ}{\ensuremath{\delta}}
\newunicodechar{ε}{\ensuremath{\varepsilon}}
\newunicodechar{θ}{\ensuremath{\theta}}
\newunicodechar{λ}{\ensuremath{\lambda}}
\newunicodechar{μ}{\ensuremath{\mu}}
\newunicodechar{σ}{\ensuremath{\sigma}}
\newunicodechar{τ}{\ensuremath{\tau}}
\newunicodechar{φ}{\ensuremath{\varphi}}
\newunicodechar{ψ}{\ensuremath{\psi}}
\newunicodechar{ω}{\ensuremath{\omega}}
\newunicodechar{Σ}{\ensuremath{\Sigma}}
% majuscules produites par \MakeUppercase dans les titres (α-aminés -> Α)
\newunicodechar{Α}{\ensuremath{\alpha}}
\newunicodechar{Β}{\ensuremath{\beta}}
\newunicodechar{Γ}{\ensuremath{\Gamma}}
\newunicodechar{Δ}{\ensuremath{\Delta}}
\newunicodechar{Ε}{\ensuremath{\varepsilon}}
\newunicodechar{Θ}{\ensuremath{\Theta}}
\newunicodechar{Λ}{\ensuremath{\Lambda}}
\newunicodechar{Μ}{\ensuremath{\mu}}
\newunicodechar{Τ}{\ensuremath{\tau}}
\newunicodechar{Φ}{\ensuremath{\Phi}}
\newunicodechar{Ψ}{\ensuremath{\Psi}}
\newunicodechar{Ω}{\ensuremath{\Omega}}
\newunicodechar{↦}{\ensuremath{\mapsto}}
\newunicodechar{∈}{\ensuremath{\in}}
\newunicodechar{∉}{\ensuremath{\notin}}
\newunicodechar{∧}{\ensuremath{\wedge}}
\newunicodechar{⇔}{\ensuremath{\Leftrightarrow}}
\newunicodechar{⟺}{\ensuremath{\Longleftrightarrow}}
\newunicodechar{⇒}{\ensuremath{\Rightarrow}}
\newunicodechar{⟹}{\ensuremath{\Longrightarrow}}
\newunicodechar{∘}{\ensuremath{\circ}}
\newunicodechar{∪}{\ensuremath{\cup}}
\newunicodechar{∩}{\ensuremath{\cap}}
\newunicodechar{≡}{\ensuremath{\equiv}}
\newunicodechar{⊂}{\ensuremath{\subset}}
\newunicodechar{⊥}{\ensuremath{\perp}}
\newunicodechar{∃}{\ensuremath{\exists}}
\newunicodechar{∀}{\ensuremath{\forall}}
\newunicodechar{∛}{\ensuremath{\sqrt[3]{\,}}}
\newunicodechar{∜}{\ensuremath{\sqrt[4]{\,}}}
\newunicodechar{⃗}{\ensuremath{{}^{\to}}}
\newunicodechar{ⁿ}{\ifmmode^{n}\else\textsuperscript{\ensuremath{n}}\fi}
\newunicodechar{ˣ}{\ifmmode^{x}\else\textsuperscript{\ensuremath{x}}\fi}
\newunicodechar{ᵇ}{\ifmmode^{b}\else\textsuperscript{\ensuremath{b}}\fi}
\newunicodechar{ʳ}{\ifmmode^{r}\else\textsuperscript{\ensuremath{r}}\fi}
\newunicodechar{ᵘ}{\ifmmode^{u}\else\textsuperscript{\ensuremath{u}}\fi}
\newunicodechar{ʸ}{\ifmmode^{y}\else\textsuperscript{\ensuremath{y}}\fi}
\newunicodechar{ᵃ}{\ifmmode^{a}\else\textsuperscript{\ensuremath{a}}\fi}
\newunicodechar{ᵉ}{\ifmmode^{e}\else\textsuperscript{\ensuremath{e}}\fi}
\newunicodechar{ᵏ}{\ifmmode^{k}\else\textsuperscript{\ensuremath{k}}\fi}
\newunicodechar{ᵖ}{\ifmmode^{p}\else\textsuperscript{\ensuremath{p}}\fi}
\newunicodechar{ᴺ}{\ifmmode^{N}\else\textsuperscript{\ensuremath{N}}\fi}
\newunicodechar{ᶜ}{\ifmmode^{c}\else\textsuperscript{\ensuremath{c}}\fi}
\newunicodechar{ʰ}{\ifmmode^{h}\else\textsuperscript{\ensuremath{h}}\fi}
\newunicodechar{ᵠ}{\ifmmode^{q}\else\textsuperscript{\ensuremath{q}}\fi}
\newunicodechar{ₙ}{\ifmmode_{n}\else\textsubscript{\ensuremath{n}}\fi}
\newunicodechar{ₐ}{\ifmmode_{a}\else\textsubscript{\ensuremath{a}}\fi}
\newunicodechar{ᵢ}{\ifmmode_{i}\else\textsubscript{\ensuremath{i}}\fi}
\newunicodechar{ᵣ}{\ifmmode_{r}\else\textsubscript{\ensuremath{r}}\fi}
\newunicodechar{ₖ}{\ifmmode_{k}\else\textsubscript{\ensuremath{k}}\fi}
\newunicodechar{ᵦ}{\ifmmode_{\beta}\else\textsubscript{\ensuremath{\beta}}\fi}
\newunicodechar{ₓ}{\ifmmode_{x}\else\textsubscript{\ensuremath{x}}\fi}
\newfontfamily\popdisplay{ArchivoBlack-Regular.ttf}[Path=../../../fonts/]
\newfontfamily\poplogo{SpaceGrotesk-Bold.ttf}[Path=../../../fonts/]
\newfontfamily\popsemi{PlusJakartaSans-SemiBold.ttf}[Path=../../../fonts/]
\newfontfamily\popxbold{PlusJakartaSans-ExtraBold.ttf}[Path=../../../fonts/]
\newfontfamily\popxboldls{PlusJakartaSans-ExtraBold.ttf}[Path=../../../fonts/,LetterSpace=3.0]

\setlist[enumerate]{leftmargin=1.7em,itemsep=5pt,topsep=4pt}
\setlist[itemize]{leftmargin=1.7em,itemsep=4pt,topsep=4pt,label={\color{poporange}\textbullet}}
\setlength{\parindent}{0pt}
\setlength{\parskip}{5pt}
\renewcommand{\arraystretch}{1.25}

% pgfplots : style Pop par défaut
\pgfplotsset{
  every axis/.append style={axis line style={popink,line width=1pt},tick style={popink},
    grid style={popline,line width=0.5pt},label style={font=\small},tick label style={font=\footnotesize}},
  popcurve/.style={poporange,line width=1.8pt,no markers,smooth},
}
% Même style disponible dans un \draw TikZ simple (hors pgfplots)
\tikzset{popcurve/.style={poporange,line width=1.8pt}}
\tikzset{popgrid/.style={popline,very thin}}
\colorlet{popcurve}{poporange} % le nom du style est parfois utilisé comme couleur

% ============================================================
% Pill / squiggle
% ============================================================
\newcommand{\poppill}[3]{%
  \tikz[baseline=-0.55ex]\node[draw=popink,line width=1.2pt,rounded corners=2.6pt,%
    fill=#2,inner xsep=6.5pt,inner ysep=3pt]{%
    \popxboldls\fontsize{7}{7}\selectfont\color{#3}\MakeUppercase{#1}};%
}
\newcommand{\popsquiggle}[1]{%
  \tikz[baseline=-0.4ex]{
    \draw[#1,line width=2.6pt,line cap=round]
      plot [smooth] coordinates {(0,0.09) (0.34,0.01) (0.68,0.21) (1.02,0.01) (1.36,0.10)};
  }%
}
% Mot-clé surligné (marqueur orange sous le texte)
\newcommand{\cle}[1]{{\popxbold\color{popdark}#1}}

% ============================================================
% En-tête / pied
% ============================================================
\pagestyle{fancy}
\fancyhf{}
\renewcommand{\headrulewidth}{0pt}
\renewcommand{\footrulewidth}{0pt}
\lhead{\raisebox{-0.15ex}{\poplogo\fontsize{13}{13}\selectfont\color{popink}OnBuch\color{poporange}+}}
\rhead{\poppill{\DOCMATIERE\ \textbullet\ \DOCNIVEAU\ \textbullet\ Leçon \DOCLECON}{white}{popink}}
\lfoot{\popxbold\fontsize{7.6}{7.6}\selectfont\color{popmuted}\MakeUppercase{Cours signé OnBuch+}\ \textbullet\ {\popsemi\DOCTITRE}}
\rfoot{\tikz[baseline=-0.6ex]\node[rounded corners=8.5pt,fill=popink,minimum height=17pt,minimum width=17pt,inner xsep=5pt,inner sep=0]{\popxbold\fontsize{8}{8}\selectfont\color{popcream}\thepage\,/\,\pageref*{LastPage}};}

% ============================================================
% Titres
% ============================================================
\newcommand{\chaptitre}[2]{%
  \begin{tcolorbox}[enhanced,colback=poporange,colframe=popink,boxrule=2.6pt,arc=11pt,
      drop shadow=popink,left=15pt,right=15pt,top=12pt,bottom=11pt,before skip=3pt,after skip=18pt]
    \poppill{#2}{popink}{popcream}\par\vspace{9pt}
    {\raggedright\hyphenpenalty=10000\exhyphenpenalty=10000\popdisplay\fontsize{22}{24}\selectfont\color{popcream}\MakeUppercase{#1}\par}
  \end{tcolorbox}
}
% Section numérotée : pastille orange avec le numéro + titre Archivo
\newcounter{popsec}
\newcommand{\coursec}[1]{%
  \stepcounter{popsec}\setcounter{popsub}{0}%
  \par\vspace{16pt}\noindent
  \begin{minipage}{\linewidth}
  \tikz[baseline=-0.7ex]\node[draw=popink,line width=1.6pt,fill=poporange,rounded corners=4pt,minimum size=22pt,inner sep=0]{\popdisplay\fontsize{12}{12}\selectfont\color{popcream}\thepopsec};%
  \hspace{8pt}{\popdisplay\fontsize{15}{17}\selectfont\color{popink}\MakeUppercase{#1}}\par
  \vspace{4pt}\noindent{\color{poporange}\rule{36pt}{3.6pt}}
  \end{minipage}\par\nopagebreak\vspace{8pt}
}
\newcounter{popsub}
\newcommand{\courssub}[1]{%
  \stepcounter{popsub}%
  \par\vspace{9pt}\noindent{\popxbold\fontsize{11.5}{13}\selectfont\color{popink}{\color{poporange}\thepopsec.\thepopsub}\hspace{6pt}#1}\par\nopagebreak\vspace{4pt}
}

% ============================================================
% Boîtes pédagogiques — même recette : fond blanc, bordure épaisse colorée,
% ombre dure encre, badge plein. Titre optionnel [#1].
% ============================================================
\tcbset{popbase/.style={breakable,enhanced,colback=white,boxrule=2.3pt,arc=9pt,
  shadow={3pt}{-3pt}{0pt}{popink},left=14pt,right=14pt,top=11pt,bottom=11pt,before skip=11pt,after skip=15pt,
  fontupper=\popsemi\fontsize{10.6}{14.5}\selectfont\color{popink},
  fontlower=\popsemi\fontsize{10.6}{14.5}\selectfont\color{popink}}}
% Coche dessinée (le glyphe ✓ n'existe pas dans les polices du document)
\newcommand{\popcheck}[1]{\tikz[baseline=-0.5ex]\draw[#1,line width=1.6pt,line cap=round,line join=round](-0.13,0.0)--(-0.04,-0.09)--(0.13,0.1);}
\providecommand{\coche}{\popcheck{popgreen}}
\providecommand{\resultat}[1]{\par\smallskip\noindent\fcolorbox{popgreen}{popgreenL}{\parbox{\dimexpr\linewidth-2\fboxsep-2\fboxrule\relax}{\textbf{Résultat.} #1}}\par\smallskip}
\newcommand{\popboxhead}[3]{\poppill{#1}{#2}{white}\ifx\relax#3\relax\else\hspace{6pt}{\popxbold\fontsize{10.6}{12}\selectfont\color{popink}#3}\fi\par\vspace{7pt}}

\newtcolorbox{aretenir}[1][]{popbase,colframe=popgreen,before upper={\popboxhead{À retenir}{popgreen}{#1}}}
% Texte en allemand (document de lecture, dialogue, extrait) : cadre
% d'étude. Un titre optionnel (source, auteur) s'affiche dans l'en-tête.
\newtcolorbox{textebox}[1][]{popbase,colframe=popblue,before upper={\popboxhead{Text}{popblue}{#1}\begin{otherlanguage}{ngerman}},after upper={\end{otherlanguage}}}
% Marques ✓ / ✗ (forme correcte / fautive) souvent utilisées par les modèles
\providecommand{\cross}{\textcolor{poppink}{\ensuremath{\times}}}
\providecommand{\xmark}{\cross}\providecommand{\crossmark}{\cross}\providecommand{\cmark}{\ensuremath{\checkmark}}
% Ligne de réponse / justification à compléter (inventée par les modèles)
\providecommand{\justification}[1]{\par\noindent\textit{Justification~:} \dotfill\par}
% \textipa (package tipa non chargé) : la police ne couvre pas l'API -> simple passage
\providecommand{\textipa}[1]{#1}
% Soulignement (ulem non chargé) : \uline, \ul
\providecommand{\uline}[1]{\underline{#1}}\providecommand{\ul}[1]{\underline{#1}}
% \alert (beamer) inventée par les modèles : texte mis en évidence
\providecommand{\alert}[1]{\textbf{\textcolor{poppink}{#1}}}
% variantes ASCII d'ordinaux écrites en macro
\providecommand{\iere}{\textsuperscript{re}}\providecommand{\ieme}{\textsuperscript{e}}\providecommand{\ere}{\textsuperscript{re}}\providecommand{\eme}{\textsuperscript{e}}\providecommand{\er}{\textsuperscript{er}}\providecommand{\ie}{\textsuperscript{e}}
% Ordinaux français écrits en macro par les modèles : 1\ière, 3\ième
\newcommand{\ière}{\textsuperscript{re}}\newcommand{\ième}{\textsuperscript{e}}
% \exemple (inventée par les modèles) : étiquette « Exemple : »
\providecommand{\exemple}{\textbf{Exemple~:}\ }
% \square utilisable aussi hors mode math (cases à cocher)
\AtBeginDocument{\let\squareMath\square\renewcommand{\square}{\ensuremath{\squareMath}}}
% Mot ou phrase en allemand à l'intérieur d'un paragraphe français
\newcommand{\all}[1]{\ifmmode\text{\textit{#1}}\else\begingroup\selectlanguage{ngerman}\itshape #1\endgroup\fi}
\let\esp\all % alias toléré
\newtcolorbox{exemplebox}[1][]{popbase,colframe=poporange,before upper={\popboxhead{Exemple}{poporange}{#1}}}
\newtcolorbox{definition}[1][]{popbase,colframe=popblue,before upper={\popboxhead{Définition}{popblue}{#1}}}
\newtcolorbox{propriete}[1][]{popbase,colframe=poppurple,before upper={\popboxhead{Propriété}{poppurple}{#1}}}
\newtcolorbox{methode}[1][]{popbase,colframe=popgold,before upper={\popboxhead{Méthode}{popgold}{#1}}}
\newtcolorbox{attention}[1][]{popbase,colframe=poppink,before upper={\popboxhead{Attention}{poppink}{#1}}}
\newtcolorbox{experience}[1][]{popbase,colframe=popblue,colback=popblueL,before upper={\popboxhead{Expérience}{popblue}{#1}}}
% « Le savais-tu ? » : carte violette pleine (contexte, culture, Cameroun)
\newtcolorbox{savaistu}[1][]{popbase,colback=poppurple,colframe=popink,
  fontupper=\popsemi\fontsize{10.4}{14.2}\selectfont\color{white},
  before upper={\poppill{Le savais-tu\,?}{white}{poppurple}\ifx\relax#1\relax\else\hspace{6pt}{\popxbold\color{white}#1}\fi\par\vspace{7pt}}}
% Exercice résolu : énoncé en haut, \tcblower puis la solution sur fond crème
\newcounter{popexo}\newcounter{popexr}
\newtcolorbox{exoresolu}[1][]{popbase,colframe=poporange,colbacklower=popcream,
  segmentation style={popink,line width=1.2pt,dashed},
  before upper={\stepcounter{popexr}\popboxhead{Exercice résolu \thepopexr}{poporange}{#1}},
  before lower={\poppill{Solution}{popink}{popcream}\par\vspace{6pt}}}
% Exercice (non résolu en place) — la correction est dans la partie Corrigés
\newtcolorbox{exercice}[1][]{popbase,colframe=popink,
  before upper={\stepcounter{popexo}\popboxhead{Exercice \thepopexo}{popink}{#1}}}
% Corrigé
\newtcolorbox{corrige}[1][]{popbase,colframe=popgreen,colback=popgreenL,
  before upper={\popboxhead{Corrigé}{popgreen}{#1}}}
% Objectifs / prérequis / bilan
\newtcolorbox{objectifs}{popbase,colframe=popink,colback=poporangeL,
  before upper={\popboxhead{Objectifs de la leçon}{popink}{}}}
\newtcolorbox{prerequis}{popbase,colframe=popink,colback=popblueL,
  before upper={\popboxhead{Prérequis}{popblue}{}}}
\newtcolorbox{bilan}{popbase,colframe=popink,colback=white,boxrule=2.8pt,
  before upper={\popboxhead{Fiche bilan}{poporange}{L'essentiel à connaître pour l'examen}}}
% Figure encadrée avec légende
\newtcolorbox{popfigure}[1][]{enhanced,breakable=false,colback=white,colframe=popline,boxrule=1.4pt,arc=8pt,
  left=10pt,right=10pt,top=10pt,bottom=8pt,before skip=10pt,after skip=12pt,halign=center,
  lower separated=false,halign lower=center,
  fontlower=\popsemi\fontsize{9}{11.5}\selectfont\color{popmuted},
  #1}
\newcommand{\legende}[1]{\par\vspace{6pt}{\popsemi\fontsize{9}{11.5}\selectfont\color{popmuted}#1}}

% ============================================================
% Signature OnBuch+
% ============================================================
\newcommand{\poplogoinline}[1]{{\poplogo\fontsize{#1}{#1}\selectfont\color{popink}OnBuch\color{poporange}+}}
\newcommand{\signatureonbuch}{%
  \begin{tcolorbox}[enhanced,colback=popink,colframe=popink,boxrule=2.2pt,arc=10pt,
      drop shadow=poporange,left=14pt,right=14pt,top=10pt,bottom=10pt,before skip=0pt,after skip=0pt]
    \tikz[baseline=-0.7ex]\node[circle,fill=popgreen,draw=popcream,line width=1.3pt,minimum size=22pt,inner sep=0]{\popcheck{white}};%
    \hspace{9pt}\begin{minipage}[c]{0.62\linewidth}
      {\popxbold\fontsize{12}{13}\selectfont\color{popcream}Signé\ \textbullet\ {\poplogo OnBuch\color{poporange}+}}\par\vspace{2pt}
      {\raggedright\popsemi\fontsize{8}{10.5}\selectfont\color{popline}\MakeUppercase{Équipe pédagogique OnBuch+}\\\MakeUppercase{Conforme au programme officiel MINESEC}\par}
    \end{minipage}\hfill
    {\poplogo\fontsize{22}{22}\selectfont\color{popcream}OnBuch\color{poporange}+}
  \end{tcolorbox}%
}

% ============================================================
% Page de garde
% #1 titre · #2 matière · #3 niveau · #4 module · #5 n° leçon · #6 durée ·
% #7 description / objectifs (texte ou itemize)
% ============================================================
\newcommand{\pagegarde}[7]{%
  \thispagestyle{empty}
  \newgeometry{a4paper,margin=1.7cm,top=1.6cm,bottom=1.4cm}%
  \begin{tikzpicture}[remember picture,overlay]
    \fill[poporange!18] ([xshift=-3.2cm,yshift=-3.6cm]current page.north east) circle (4.6cm);
    \fill[poppurple!14] ([xshift=2.4cm,yshift=7.6cm]current page.south west) circle (3.8cm);
  \end{tikzpicture}%
  {\poplogo\fontsize{30}{30}\selectfont\color{popink}OnBuch\color{poporange}+}\hfill
  \raisebox{0.6ex}{\poppill{Cours signé \textbullet\ Premium}{popink}{popcream}}\par
  \vspace{1pt}\popsquiggle{poppurple}\par
  \vspace{8pt}
  \begin{tcolorbox}[enhanced,colback=poporange,colframe=popink,boxrule=2.8pt,arc=13pt,
      drop shadow=popink,left=18pt,right=18pt,top=12pt,bottom=13pt,after skip=10pt]
    \poppill{#2 \textbullet\ #3}{popink}{popcream}\hspace{5pt}\poppill{#4}{white}{popink}\par\vspace{8pt}
    {\popdisplay\fontsize{14}{14}\selectfont\color{popink}LEÇON #5}\par\vspace{4pt}
    {\raggedright\hyphenpenalty=10000\exhyphenpenalty=10000\popdisplay\fontsize{25}{27}\selectfont\color{popcream}\MakeUppercase{#1}\par}
    \vspace{8pt}
    {\popxbold\fontsize{9.5}{9.5}\selectfont\color{popink}\MakeUppercase{Durée conseillée : #6}}
  \end{tcolorbox}
  \begin{tcolorbox}[enhanced,colback=white,colframe=popink,boxrule=2.2pt,arc=10pt,
      drop shadow=popink,left=16pt,right=16pt,top=10pt,bottom=10pt,after skip=8pt]
    \poppill{Dans cette leçon}{poppurple}{white}\par\vspace{5pt}
    \popsemi\fontsize{9.3}{12.6}\selectfont\color{popink}#7
  \end{tcolorbox}
  \noindent\begin{minipage}[t]{0.487\linewidth}
    \begin{tcolorbox}[enhanced,colback=popgreen,colframe=popink,boxrule=2.2pt,arc=10pt,
        drop shadow=popink,left=11pt,right=11pt,top=10pt,bottom=10pt,height=3.1cm,valign=center]
      \tcbox[colback=white,colframe=popink,boxrule=1.3pt,arc=5pt,boxsep=0pt,nobeforeafter,
        left=3pt,right=3pt,top=3pt,bottom=3pt,valign=center]{\qrcode[height=1.9cm]{https://play.google.com/store/apps/details?id=com.luvvix.onbuch&hl=fr}}%
      \hspace{8pt}\begin{minipage}[c]{0.5\linewidth}
        {\popdisplay\fontsize{11}{12.5}\selectfont\color{white}\MakeUppercase{Télécharge OnBuch}}\par\vspace{4pt}
        {\raggedright\popsemi\fontsize{8.4}{11}\selectfont\color{white}Gratuit \textbullet\ Google Play\\Tuteur IA, annales, concours, résultats\par}
      \end{minipage}
    \end{tcolorbox}
  \end{minipage}\hfill
  \begin{minipage}[t]{0.487\linewidth}
    \begin{tcolorbox}[enhanced,colback=poppurple,colframe=popink,boxrule=2.2pt,arc=10pt,
        drop shadow=popink,left=11pt,right=11pt,top=10pt,bottom=10pt,height=3.1cm,valign=center]
      \tikz[baseline=-0.7ex]\node[circle,fill=white,draw=popink,line width=1.3pt,minimum size=1.6cm,inner sep=0]{\popdisplay\fontsize{24}{24}\selectfont\color{poppurple}+};%
      \hspace{8pt}\begin{minipage}[c]{0.6\linewidth}
        {\popdisplay\fontsize{11}{12.5}\selectfont\color{white}\MakeUppercase{Passe à OnBuch+}}\par\vspace{4pt}
        {\raggedright\popsemi\fontsize{8.4}{11}\selectfont\color{white}Tous les cours signés, Tuteur IA illimité, fiches PDF — onglet OnBuch+ de l'app\par}
      \end{minipage}
    \end{tcolorbox}
  \end{minipage}
  \vspace{8pt}
  \popcontenu
  \vfill
  \signatureonbuch
  \restoregeometry
  \newpage
}

% Rangée « Ce cours contient » (page de garde)
\newcommand{\poptile}[3]{% #1 couleur #2 gros texte #3 légende
  \begin{tcolorbox}[enhanced,colback=#1,colframe=popink,boxrule=2pt,arc=9pt,shadow={2.5pt}{-2.5pt}{0pt}{popink},
      left=6pt,right=6pt,top=9pt,bottom=9pt,halign=center,valign=center,height=1.75cm,width=\linewidth,nobeforeafter]
    {\popdisplay\fontsize{12.5}{13}\selectfont\color{white}\MakeUppercase{#2}}\par\vspace{3pt}
    {\centering\popsemi\fontsize{7.8}{9.5}\selectfont\color{white}#3\par}
  \end{tcolorbox}}
\newcommand{\popcontenu}{%
  \par\noindent{\popxboldls\fontsize{7.5}{7.5}\selectfont\color{popmuted}CE COURS SIGNÉ CONTIENT}\par\vspace{7pt}
  \noindent\begin{minipage}[t]{0.235\linewidth}\poptile{popblue}{Cours}{complet, illustré, pas à pas}\end{minipage}\hfill
  \begin{minipage}[t]{0.235\linewidth}\poptile{poporange}{Méthodes}{et exercices résolus}\end{minipage}\hfill
  \begin{minipage}[t]{0.235\linewidth}\poptile{poppink}{Exercices}{type examen + corrigés}\end{minipage}\hfill
  \begin{minipage}[t]{0.235\linewidth}\poptile{popgreen}{Fiche bilan}{l'essentiel à retenir}\end{minipage}\par}

% Sommaire
\newtcolorbox{sommaire}{enhanced,colback=white,colframe=popink,boxrule=2.2pt,arc=10pt,
  drop shadow=popink,left=18pt,right=18pt,top=13pt,bottom=13pt,before skip=6pt,after skip=20pt,
  before upper={\poppill{Sommaire}{poppurple}{white}\par\vspace{9pt}\popsemi\fontsize{10.6}{14.5}\selectfont\color{popink}}}

% Signature de fin de cours — poussée tout en bas de la dernière page
\newcommand{\findecours}{%
  \par\vfill\nopagebreak
  \begin{center}\popsquiggle{poporange}\end{center}\vspace{4pt}
  \signatureonbuch
}
\AtBeginDocument{\def\llbracket{\mathopen{[\mkern-3.5mu[}}\def\rrbracket{\mathclose{]\mkern-3.5mu]}}} % intervalles d'entiers (glyphe ⟦ absent de la police)
% Glyphes absents de Fira Math : redéfinis à partir de glyphes disponibles
\AtBeginDocument{%
  \def\setminus{\mathbin{\backslash}}\let\smallsetminus\setminus
  \def\vdots{\mathord{\vbox{\baselineskip3.2pt\lineskiplimit0pt\kern4pt\hbox{.}\hbox{.}\hbox{.}}}}%
  \def\ddots{\mathinner{\mkern1mu\raise6.5pt\hbox{.}\mkern2mu\raise3.6pt\hbox{.}\mkern2mu\raise0.7pt\hbox{.}\mkern1mu}}%
  \def\triangle{\mathord{\scalebox{0.9}{$\symup{\Delta}$}}}%
  \def\nabla{\mathord{\raisebox{\depth}{\scalebox{1}[-1]{$\symup{\Delta}$}}}}%
  \def\checkmark{\popcheck{popdark}}%
  \def\star{\mathbin{\ast}}\def\bigstar{\mathord{\ast}}%
  \def\diamond{\mathbin{\scalebox{0.6}{$\Diamond$}}}%
  \def\bigcup{\mathop{\mathchoice{\raisebox{-0.3em}{\scalebox{1.7}{$\cup$}}}{\scalebox{1.25}{$\cup$}}{\cup}{\cup}}\displaylimits}%
  \def\bigcap{\mathop{\mathchoice{\raisebox{-0.3em}{\scalebox{1.7}{$\cap$}}}{\scalebox{1.25}{$\cap$}}{\cap}{\cap}}\displaylimits}%
  \def\lesseqqgtr{\mathrel{\lessgtr}}\def\bigoplus{\mathop{\oplus}\displaylimits}%
}
\providecommand{\ssi}{si et seulement si} % abréviation fréquente
\AtBeginDocument{\providecommand{\wideparen}{\overparen}} % arc de cercle

%% FIGFIX-BEGIN v5 — correctifs globaux des figures (docs/onbuch-generation/tools/figfix)
\usepackage{adjustbox}\usepackage{etoolbox}\tcbuselibrary{raster}
% toute figure TikZ/pgfplots plus large que la ligne est réduite à la largeur disponible
\BeforeBeginEnvironment{tikzpicture}{\begin{adjustbox}{max width=\linewidth}}
\AfterEndEnvironment{tikzpicture}{\end{adjustbox}}
% légendes pgfplots lisibles par-dessus les courbes
\pgfplotsset{every axis/.append style={legend style={fill=white,fill opacity=0.92,text opacity=1,draw=black!15,inner sep=3pt}}}
% étiquettes d'arêtes (syntaxe edge["..."]) avec fond blanc
\tikzset{every edge quotes/.append style={fill=white,inner sep=1.2pt}}
%% FIGFIX-END
\begin{document}

\pagegarde{L'extraction des minerais}{Allemand}{Tle A}{Chapitre 3 : Environnement, santé et bien-être}{AL12}{2 h}{Cette leçon propose un texte authentique sur l'extraction des minerais, son lexique spécialisé, les structures du passif au prétérit et les relatives avec préposition, suivis d'activités de compréhension, de commentaire et de production guidée.
\vspace{4pt}
\begin{multicols}{2}\small
\begin{itemize}
\item Comprendre globalement et en détail un texte allemand sur l'extraction des minerais.
\item Maîtriser le vocabulaire thématique (Erz, Bergbau, Mine, Gold, Rohstoff, Arbeiter) avec articles et pluriels.
\item Former et reconnaître le passif au prétérit (wurde + Partizip II) dans des phrases authentiques.
\item Construire des relatives avec préposition (z. B. der Bergbau, in dem ... gearbeitet wird).
\item Répondre par écrit à des questions de compréhension en allemand.
\item Rédiger un résumé structuré du texte support (environ 80 mots).
\end{itemize}\end{multicols}}

\begin{sommaire}
\begin{multicols}{2}
\begin{enumerate}
\item 1. Mise en situation \& texte support
\item 2. Compréhension globale et détaillée
\item 3. Lexique thématique par champs
\item 4. Grammaire : Le passif au prétérit
\item 5. Grammaire : Relatives avec préposition
\item 6. Commentaire guidé \& production écrite
\item 7. Production orale guidée \& réinvestissement
\item Activité d'intégration
\item Méthodes et exercices résolus
\item Exercices
\item Corrigés des exercices
\item Fiche bilan
\end{enumerate}
\end{multicols}
\end{sommaire}

\begin{objectifs}
\begin{itemize}
\item Comprendre globalement et en détail un texte allemand sur l'extraction des minerais.
\item Maîtriser le vocabulaire thématique (Erz, Bergbau, Mine, Gold, Rohstoff, Arbeiter) avec articles et pluriels.
\item Former et reconnaître le passif au prétérit (wurde + Partizip II) dans des phrases authentiques.
\item Construire des relatives avec préposition (z. B. der Bergbau, in dem ... gearbeitet wird).
\item Répondre par écrit à des questions de compréhension en allemand.
\item Rédiger un résumé structuré du texte support (environ 80 mots).
\item Produire un court paragraphe d'opinion sur les impacts écologiques du Bergbau.
\item Appliquer les règles de conjugaison du prétérit passif à de nouveaux verbes du champ lexical.
\end{itemize}
\end{objectifs}

\begin{prerequis}
\begin{itemize}
\item Connaître le présent et le parfait du passif (wird + Partizip II / ist + Partizip II).
\item Savoir former le prétérit actif des verbes réguliers et irréguliers courants.
\item Maîtriser les pronoms relatifs simples (der, die, das) au nominatif et à l'accusatif.
\item Disposer d'un vocabulaire de base sur l'environnement et l'économie (Umwelt, Wirtschaft, Rohstoff).
\item Être capable de rédiger un paragraphe cohérent en allemand (niveau B1).
\end{itemize}
\end{prerequis}

\newpage

\coursec{Mise en situation \& texte support}

\courssub{Pourquoi parler des minerais en cours d'allemand ?}

Imaginez un reportage à la radio : à \textbf{Minim-Martap}, dans l'Adamaoua, on prépare l'exploitation de la bauxite ; à \textbf{Bétaré-Oya}, dans l'Est, l'orpaillage fait vivre des milliers de familles. En Allemagne, la région de la \textbf{Ruhr} a été pendant deux siècles le cœur de l'extraction du charbon, avant la fermeture de la dernière mine en 2018. En \textbf{Autriche}, on exploite encore le sel à Hallstatt, un site connu depuis la préhistoire. La \textbf{Suisse}, elle, est pauvre en minerais : elle importe des matières premières pour son industrie de haute technologie (montres, machines de précision).

Ces réalités camerounaises et germanophones posent notre \textbf{problématique} : comment un texte allemand raconte-t-il l'histoire, les techniques et les conséquences de l'extraction minière, et quels outils de langue (vocabulaire, passif au prétérit, relatives avec préposition) faut-il maîtriser pour le comprendre et en parler au Bac ?

\begin{savaistu}[Le Cameroun et l'Allemagne : une histoire minière]
Sous le protectorat allemand (1884-1916), les premières prospections géologiques du Cameroun ont été menées par des ingénieurs allemands. Aujourd'hui encore, des entreprises germanophones vendent des machines d'extraction (\all{Bergbaumaschinen}) en Afrique centrale. Le mot allemand \all{der Bergbau} vient de \all{der Berg} (la montagne) et du verbe \all{bauen} (construire, ici : exploiter) : littéralement « travailler la montagne ». Le mineur se dit \all{der Bergmann} (pluriel \all{die Bergleute}) : « l'homme de la montagne ».
\end{savaistu}

\courssub{Carte mentale : le vocabulaire minier}

Avant de lire, activons les six mots-clés du thème. Observez la carte mentale : chaque mot est donné avec son \textbf{article} et son \textbf{pluriel}, à apprendre toujours ensemble.

\begin{popfigure}
\begin{center}\tcbox[colback=popblue,colframe=popblue,coltext=white,arc=9pt,boxsep=4pt,left=6pt,right=6pt]{\bfseries\small der Bergbau l'exploitation minière}\end{center}
\vspace{-2pt}
\begin{tcbraster}[raster columns=3,raster equal height=rows,raster column skip=6pt,raster row skip=6pt,enhanced,blankest]
\begin{tcolorbox}[enhanced,colback=poporange,colframe=poporange,coltext=white,boxrule=0.9pt,arc=3pt,left=4pt,right=4pt,top=3pt,bottom=3pt,boxsep=1pt,halign=left]\bfseries\scriptsize das Erz, -e le minerai\end{tcolorbox}
\begin{tcolorbox}[enhanced,colback=popgreen,colframe=popgreen,coltext=white,boxrule=0.9pt,arc=3pt,left=4pt,right=4pt,top=3pt,bottom=3pt,boxsep=1pt,halign=left]\bfseries\scriptsize die Mine, -n la mine\end{tcolorbox}
\begin{tcolorbox}[enhanced,colback=popgold,colframe=popgold,coltext=white,boxrule=0.9pt,arc=3pt,left=4pt,right=4pt,top=3pt,bottom=3pt,boxsep=1pt,halign=left]\bfseries\scriptsize das Gold l'or\end{tcolorbox}
\begin{tcolorbox}[enhanced,colback=poppurple,colframe=poppurple,coltext=white,boxrule=0.9pt,arc=3pt,left=4pt,right=4pt,top=3pt,bottom=3pt,boxsep=1pt,halign=left]\bfseries\scriptsize der Rohstoff, -e la matière première\end{tcolorbox}
\begin{tcolorbox}[enhanced,colback=poppink,colframe=poppink,coltext=white,boxrule=0.9pt,arc=3pt,left=4pt,right=4pt,top=3pt,bottom=3pt,boxsep=1pt,halign=left]\bfseries\scriptsize der Arbeiter, - l'ouvrier\end{tcolorbox}
\begin{tcolorbox}[enhanced,colback=popblue!60,colframe=popblue!60,coltext=white,boxrule=0.9pt,arc=3pt,left=4pt,right=4pt,top=3pt,bottom=3pt,boxsep=1pt,halign=left]\bfseries\scriptsize die Kohle, -n le charbon\end{tcolorbox}
\end{tcbraster}

\legende{Carte mentale du vocabulaire minier : article + pluriel à mémoriser ensemble.}
\end{popfigure}

\begin{definition}[Les six mots-clés du thème]
\begin{itemize}
\item \cle{das Erz}, -e : le minerai (roche contenant un métal). \all{das Eisenerz} = le minerai de fer.
\item \cle{der Bergbau} (sans pluriel) : l'exploitation minière, l'industrie minière.
\item \cle{die Mine}, -n : la mine (le lieu d'extraction souterrain).
\item \cle{das Gold} (sans pluriel) : l'or. Attention au genre : \all{das}, pas \all{der} !
\item \cle{der Rohstoff}, -e : la matière première. Famille : \all{roh} (brut) + \all{der Stoff} (la matière).
\item \cle{der Arbeiter}, - : l'ouvrier ; au féminin \all{die Arbeiterin}, -nen. Famille : \all{die Arbeit} (le travail), \all{arbeiten} (travailler).
\end{itemize}
\end{definition}

\courssub{Lecture du texte support}

Lisons le texte à haute voix, en respectant la mélodie de la phrase allemande. Les lignes sont numérotées pour faciliter le repérage. En lisant, soulignez mentalement toutes les formes en \all{wurde / wurden + participe II} : ce sera notre point de grammaire.

\begin{textebox}[Der Bergbau in Deutschland – Geschichte und Gegenwart]
\textbf{1} Der Bergbau hat in Deutschland eine lange Geschichte. Schon in der Antike\\
\textbf{2} wurde in den Bergen Erz abgebaut. Im Mittelalter wurden viele Minen\\
\textbf{3} eröffnet, besonders im Harz und im Erzgebirge. Dort wurde Silber gefunden,\\
\textbf{4} und die Städte in der Nähe wurden reich.\\
\textbf{5} Im 19. Jahrhundert begann die Industrialisierung. In der Ruhrregion\\
\textbf{6} wurde Kohle gefördert, und überall wurden neue Zechen gebaut. Tausende\\
\textbf{7} Arbeiter kamen aus dem Osten, um in den Minen zu arbeiten. Die Arbeit\\
\textbf{8} war hart und gefährlich: Viele Bergleute wurden bei Unfällen verletzt\\
\textbf{9} oder getötet. Trotzdem war der Bergbau die Basis der deutschen Industrie.\\
\textbf{10} Aber der Bergbau hatte auch Folgen für die Umwelt. Flüsse wurden\\
\textbf{11} verschmutzt, Wälder wurden abgeholzt, und die Luft in den Städten\\
\textbf{12} wurde schlecht. Noch heute sieht man in der Ruhrregion die Spuren\\
\textbf{13} dieser Zeit.\\
\textbf{14} Im Jahr 2018 wurde die letzte Zeche in Deutschland geschlossen.\\
\textbf{15} Viele alte Minen wurden in Museen umgewandelt; die Zeche Zollverein\\
\textbf{16} in Essen wurde sogar von der UNESCO zum Weltkulturerbe erklärt.\\
\textbf{17} Heute importiert Deutschland die meisten Rohstoffe, zum Beispiel aus\\
\textbf{18} Afrika. Auch in Kamerun werden Gold und Bauxit abgebaut. Die Frage\\
\textbf{19} bleibt: Wie kann man Rohstoffe gewinnen, ohne die Natur zu zerstören?\\
\textbf{20} Diese Frage wurde schon oft gestellt – eine Antwort wurde noch nicht\\
\textbf{21} gefunden.
\end{textebox}

\textbf{Glossaire} (mot — genre et pluriel — traduction) :

\begin{center}
\begin{tabularx}{\linewidth}{|l|l|X|}
\hline
\rowcolor{popblueL} Mot allemand & Genre / pluriel & Traduction \\
\hline
abbauen (trennbar) & verbe & extraire, exploiter \\
\hline
das Erzgebirge & nom propre & les Monts Métallifères \\
\hline
fördern & verbe & extraire (le charbon), produire \\
\hline
die Zeche, -n & fém. & la mine de charbon (mot régional, Ruhr) \\
\hline
der Bergmann, -leute & masc. & le mineur \\
\hline
der Unfall, -fälle & masc. & l'accident \\
\hline
die Umwelt & fém. (sans pluriel) & l'environnement \\
\hline
verschmutzen & verbe & polluer \\
\hline
abholzen (trennbar) & verbe & déboiser \\
\hline
die Spur, -en & fém. & la trace \\
\hline
umwandeln (trennbar) & verbe & transformer, convertir \\
\hline
das Weltkulturerbe & neutre (sans pluriel) & le patrimoine mondial de l'humanité \\
\hline
gewinnen (gewann, gewonnen) & verbe fort & obtenir, extraire, gagner \\
\hline
zerstören & verbe & détruire \\
\hline
\end{tabularx}
\end{center}

\begin{methode}[Lire un texte allemand pour la première fois]
\begin{enumerate}
\item Lire le titre et regarder la frise : quel est le thème, quelle est la période ?
\item Première lecture silencieuse : repérer les mots connus et les mots transparents (\all{Industrialisierung, Museum, Import}).
\item Deuxième lecture à haute voix : souligner les formes \all{wurde(n) + Partizip II} et entourer les mots inconnus.
\item Vérifier les mots inconnus dans le glossaire avant de répondre aux questions.
\item Résumer chaque paragraphe en une phrase française : histoire — techniques — impacts — avenir.
\end{enumerate}
\end{methode}

\coursec{Compréhension globale et détaillée}

\courssub{Compréhension globale : la structure du texte}

Le texte suit un plan chronologique clair. Repérons les quatre grandes étapes :

\begin{itemize}
\item \textbf{Lignes 1-4 :} l'histoire ancienne (Antiquité, Moyen Âge) — \all{Antike, Mittelalter}.
\item \textbf{Lignes 5-9 :} l'industrialisation au XIX\textsuperscript{e} siècle — \all{Industrialisierung, Ruhrregion}.
\item \textbf{Lignes 10-13 :} les impacts écologiques — \all{Umwelt, verschmutzt, abgeholzt}.
\item \textbf{Lignes 14-21 :} la situation actuelle et l'avenir — \all{2018, heute, die Frage bleibt}.
\end{itemize}

\begin{popfigure}
\begin{tikzpicture}
\draw[-{Stealth}, thick, popdark] (0,0) -- (14.5,0);
\foreach \x/\t/\d in {0.8/{Antiquité}/{\all{Erz wurde}\\\all{abgebaut}}, 4.6/{Moyen Âge}/{\all{Minen wurden}\\\all{eröffnet}}, 8.4/{XIX\textsuperscript{e} s.}/{\all{Kohle wurde}\\\all{gefördert}}, 12.2/{Aujourd'hui}/{\all{Zeche wurde}\\\all{geschlossen}}}{
  \fill[poporange] (\x,0) circle (0.12);
  \node[above, align=center, font=\small\bfseries, text=popblue] at (\x,0.25) {\t};
  \node[below, align=center, font=\footnotesize, text=popdark] at (\x,-0.25) {\d};
}
\end{tikzpicture}
\legende{Frise chronologique du texte : chaque étape est marquée par une forme passive au prétérit.}
\end{popfigure}

\courssub{Compréhension détaillée : questions sur le texte}

\begin{exercice}[Questions de compréhension]
Répondez en allemand, par des phrases complètes, en vous appuyant sur les lignes indiquées.
\begin{enumerate}
\item Wo wurde schon in der Antike Erz abgebaut? (l. 1-2)
\item Warum wurden die Städte im Harz und im Erzgebirge reich? (l. 3-4)
\item Warum kamen tausende Arbeiter in die Ruhrregion? (l. 6-7)
\item Welche Folgen hatte der Bergbau für die Umwelt? (l. 10-12)
\item Was geschah im Jahr 2018? (l. 14)
\item Welche Frage bleibt am Ende des Textes offen? (l. 19)
\end{enumerate}
\end{exercice}

\begin{corrige}[Questions de compréhension]
\begin{enumerate}
\item \all{Schon in der Antike wurde in den Bergen Erz abgebaut.} « Dès l'Antiquité, on extrayait du minerai dans les montagnes. »
\item \all{Die Städte wurden reich, weil dort Silber gefunden wurde.} « Les villes sont devenues riches parce qu'on y a trouvé de l'argent. » (subordonnée avec \all{weil} : verbe à la fin)
\item \all{Sie kamen, um in den Minen zu arbeiten.} « Ils sont venus pour travailler dans les mines. »
\item \all{Flüsse wurden verschmutzt, Wälder wurden abgeholzt und die Luft wurde schlecht.} « Les rivières ont été polluées, les forêts déboisées et l'air est devenu mauvais. »
\item \all{Im Jahr 2018 wurde die letzte Zeche in Deutschland geschlossen.} « En 2018, la dernière mine de charbon d'Allemagne a été fermée. »
\item \all{Es bleibt die Frage offen, wie man Rohstoffe gewinnen kann, ohne die Natur zu zerstören.} « Reste ouverte la question de savoir comment extraire des matières premières sans détruire la nature. »
\end{enumerate}
\end{corrige}

\begin{exemplebox}[Repérage visuel des formes passives dans le texte]
Relevons les formes en \all{wurde / wurden + Partizip II} :
\begin{itemize}
\item \all{wurde ... abgebaut} (l. 2) : était extrait / a été extrait
\item \all{wurden ... eröffnet} (l. 2-3) : ont été ouvertes
\item \all{wurde Silber gefunden} (l. 3) : on a trouvé de l'argent
\item \all{wurde Kohle gefördert} (l. 6) : on extrayait du charbon
\item \all{wurden neue Zechen gebaut} (l. 6) : de nouvelles mines ont été construites
\item \all{wurden ... verletzt oder getötet} (l. 8-9) : ont été blessés ou tués
\item \all{wurden verschmutzt / abgeholzt} (l. 11) : ont été polluées / déboisées
\item \all{wurde ... geschlossen} (l. 14) : a été fermée
\item \all{wurde ... zum Weltkulturerbe erklärt} (l. 16) : a été classée au patrimoine mondial
\end{itemize}
Observation : la forme est toujours \cle{wurde(n)} (auxiliaire conjugué au prétérit, en 2\textsuperscript{e} position) + \cle{participe II} (renvoyé à la fin de la phrase).
\end{exemplebox}

\coursec{Lexique thématique par champs}

\courssub{Champ 1 : les matières et les lieux}

\begin{center}
\begin{tabularx}{\linewidth}{|l|l|X|}
\hline
\rowcolor{popblueL} Allemand & Genre / pluriel & Français \\
\hline
das Erz & -e & le minerai \\
\hline
das Gold & sans pluriel & l'or \\
\hline
das Silber & sans pluriel & l'argent (métal) \\
\hline
die Kohle & -n & le charbon \\
\hline
das Bauxit & sans pluriel & la bauxite \\
\hline
der Rohstoff & -e & la matière première \\
\hline
die Mine & -n & la mine \\
\hline
die Zeche & -n & la mine de charbon \\
\hline
der Bergbau & sans pluriel & l'exploitation minière \\
\hline
\end{tabularx}
\end{center}

\courssub{Champ 2 : les personnes et le travail}

\begin{center}
\begin{tabularx}{\linewidth}{|l|l|X|}
\hline
\rowcolor{popgreenL} Allemand & Genre / pluriel & Français \\
\hline
der Arbeiter / die Arbeiterin & - / -nen & l'ouvrier / l'ouvrière \\
\hline
der Bergmann & -leute & le mineur \\
\hline
die Arbeit & -en & le travail \\
\hline
der Unfall & -fälle & l'accident \\
\hline
gefährlich & adjectif & dangereux \\
\hline
hart & adjectif & dur \\
\hline
\end{tabularx}
\end{center}

\courssub{Champ 3 : les verbes de l'extraction et de l'environnement}

\begin{center}
\begin{tabularx}{\linewidth}{|l|l|X|}
\hline
\rowcolor{poporangeL} Verbe & Prétérit / Partizip II & Français \\
\hline
abbauen (trennbar) & baute ab / abgebaut & extraire, exploiter \\
\hline
fördern & förderte / gefördert & extraire, produire \\
\hline
gewinnen & gewann / gewonnen & obtenir, extraire \\
\hline
eröffnen & eröffnete / eröffnet & ouvrir \\
\hline
schließen & schloss / geschlossen & fermer \\
\hline
verschmutzen & verschmutzte / verschmutzt & polluer \\
\hline
abholzen (trennbar) & holzte ab / abgeholzt & déboiser \\
\hline
zerstören & zerstörte / zerstört & détruire \\
\hline
umwandeln (trennbar) & wandelte um / umgewandelt & transformer \\
\hline
importieren & importierte / importiert & importer \\
\hline
\end{tabularx}
\end{center}

\begin{attention}[Faux amis et pièges lexicaux]
\begin{itemize}
\item \all{die Mine} n'est pas « la mine » au sens de « mine de crayon » (\all{die Mine} convient aussi, mais au Bac c'est toujours la mine d'extraction) ; ne pas confondre avec \all{das Minenfeld} (champ de mines).
\item \all{fördern} = extraire, produire ; \all{fördern} avec Umlaut signifie « encourager, soutenir » : \all{Der Staat fördert junge Unternehmen.} « L'État soutient les jeunes entreprises. » Même orthographe, sens différents selon le contexte !
\item \all{gewinnen} = extraire, mais aussi « gagner » : \all{eine Medaille gewinnen} « gagner une médaille ».
\item \all{das Silber} = l'argent métal ; « l'argent » (monnaie) se dit \all{das Geld}. Ne jamais traduire « les mineurs gagnent de l'argent » par \all{Silber} !
\item Genre : \all{das Gold}, \all{das Erz}, \all{das Silber} (les métaux sont neutres), mais \all{die Kohle} et \all{die Mine} sont féminins.
\end{itemize}
\end{attention}

\coursec{Grammaire : le passif au prétérit}

\courssub{Formation et emploi}

\begin{propriete}[Formation du passif au prétérit (Präteritum Passiv)]
\textbf{Forme :} \cle{wurde} (singulier) ou \cle{wurden} (pluriel) + \cle{Partizip II} placé en fin de phrase.

\begin{center}
\begin{tabular}{|l|l|l|}
\hline
\rowcolor{popblueL} Personne & Singulier & Pluriel \\
\hline
1\textsuperscript{re} pers. & ich wurde & wir wurden \\
\hline
2\textsuperscript{e} pers. & du wurdest & ihr wurdet \\
\hline
3\textsuperscript{e} pers. & er/sie/es wurde & sie/Sie wurden \\
\hline
\end{tabular}
\end{center}

\textbf{Emploi :} on l'utilise quand l'\emph{action} compte plus que l'\emph{agent} (qui est souvent inconnu ou sans importance), exactement comme le passif français ou le « on » français. L'agent, s'il est exprimé, est introduit par \all{von} + datif : \all{von der UNESCO} (l. 16).

\textbf{Comparaison avec le français :} le français dit « la mine \emph{a été} ouverte » ; l'allemand dit \all{die Mine \emph{wurde} eröffnet}. L'auxiliaire passif allemand est \all{werden} (jamais \all{sein} !), dont le prétérit est \all{wurde(n)}.

\textbf{Attention aux verbes à particule :} le participe garde la particule insérée : \all{abbauen} $\rightarrow$ \all{abgebaut} ; \all{umwandeln} $\rightarrow$ \all{umgewandelt}.
\end{propriete}

\begin{popfigure}
\begin{tikzpicture}[node distance=0.6cm]
\node[draw=popblue, rounded corners, fill=popblueL, align=center, font=\small] (aktiv) {Actif : \all{Man eröffnete die Mine.}\\ « On a ouvert la mine. »};
\node[draw=poporange, rounded corners, fill=poporangeL, align=center, font=\small, below=of aktiv] (passiv) {Passif : \all{Die Mine \textbf{wurde} eröffnet.}\\ « La mine a été ouverte. »};
\draw[-{Stealth}, very thick, popdark] (aktiv) -- (passiv) node[midway, right, align=left, font=\footnotesize, text=popdark]{COD \all{die Mine} $\rightarrow$ sujet\\\all{man} disparaît\\verbe $\rightarrow$ \all{wurde + Partizip II}};
\end{tikzpicture}
\legende{Transformation actif $\rightarrow$ passif : le COD devient sujet, l'auxiliaire \all{werden} se met au prétérit.}
\end{popfigure}

\begin{exemplebox}[Exemples traduits]
\all{Die Mine wurde 1900 eröffnet.} « La mine a été ouverte en 1900. »\\
\all{Viele Minen wurden im Mittelalter eröffnet.} « Beaucoup de mines ont été ouvertes au Moyen Âge. »\\
\all{In Kamerun wurde Gold gefunden.} « Au Cameroun, on a trouvé de l'or. »\\
\all{Die Flüsse wurden verschmutzt.} « Les rivières ont été polluées. »\\
\all{Die Zeche wurde von der UNESCO zum Weltkulturerbe erklärt.} « La mine a été classée par l'UNESCO au patrimoine mondial. »\\
\all{Die Arbeiter wurden bei einem Unfall verletzt.} « Les ouvriers ont été blessés dans un accident. »
\end{exemplebox}

\begin{attention}[Erreur fréquente : wurde ou wurden ?]
Les francophones confondent souvent \all{wurde} et \all{wurden} devant les sujets collectifs ou les noms de matière. Règle simple : le verbe s'accorde avec le \textbf{sujet grammatical}, pas avec le sens.
\begin{itemize}
\item \all{Die Kohle wurde gefördert.} (singulier : \all{die Kohle}) « Le charbon a été extrait. »
\item \all{Tausende Arbeiter wurden verletzt.} (pluriel : \all{Arbeiter}) « Des milliers d'ouvriers ont été blessés. »
\item Faux : \all{Gold und Bauxit wurde abgebaut.} — Correct : \all{Gold und Bauxit \textbf{wurden} abgebaut} (deux éléments = pluriel).
\end{itemize}
Autre piège : ne jamais écrire \all{war ... eröffnet} pour raconter une action passée ; \all{war + Partizip II} (Zustandspassiv) décrit un \emph{état résultant} : \all{Die Mine war geschlossen.} « La mine était fermée » (état), mais \all{Die Mine wurde 2018 geschlossen.} « La mine a été fermée en 2018 » (action).
\end{attention}

\begin{exoresolu}[Mise en application immédiate]
Énoncé : mettez les phrases suivantes au passif au prétérit, comme dans le texte.
\begin{enumerate}
\item Man eröffnete 1900 eine Mine in der Ruhrregion.
\item Man förderte viele Jahre lang Kohle.
\item Man verschmutzte die Flüsse.
\end{enumerate}
\tcblower
Solution détaillée :
\begin{enumerate}
\item Le sujet passif est \all{eine Mine} (singulier) $\rightarrow$ \all{wurde} ; participe \all{eröffnet} en fin : \all{1900 wurde eine Mine in der Ruhrregion eröffnet.} (ou \all{Eine Mine wurde 1900 in der Ruhrregion eröffnet.}) « Une mine a été ouverte en 1900 dans la Ruhr. »
\item Sujet \all{Kohle} (singulier, nom de matière) $\rightarrow$ \all{wurde} : \all{Viele Jahre lang wurde Kohle gefördert.} « Pendant de nombreuses années, on a extrait du charbon. »
\item Sujet \all{die Flüsse} (pluriel) $\rightarrow$ \all{wurden} : \all{Die Flüsse wurden verschmutzt.} « Les rivières ont été polluées. »
\end{enumerate}
Méthode : (1) identifier le complément d'objet direct qui devient sujet ; (2) choisir \all{wurde} ou \all{wurden} selon ce sujet ; (3) placer le participe II en fin de phrase ; (4) supprimer \all{man} ou le remplacer par \all{von} + datif si l'agent est utile.
\end{exoresolu}

\coursec{Grammaire : les relatives avec préposition}

\courssub{Observation}

Dans le texte, la ligne 7 dit : \all{Tausende Arbeiter kamen, um in den Minen zu arbeiten.} On pourrait dire autrement : \all{Die Minen, \textbf{in denen} die Arbeiter arbeiteten, waren gefährlich.} « Les mines \emph{dans lesquelles} les ouvriers travaillaient étaient dangereuses. » C'est la \textbf{relative avec préposition}.

\begin{propriete}[La relative avec préposition]
Quand l'antécédent est le complément d'une \textbf{préposition} (dans, avec, pour, sur...), la préposition se place \textbf{devant le pronom relatif}, qui se met au \textbf{cas exigé par la préposition} (accusatif ou datif) et s'accorde en genre et en nombre avec l'antécédent.

\begin{center}
\begin{tabular}{|l|l|l|}
\hline
\rowcolor{popblueL} Antécédent & Datif (in, mit, von, zu, bei...) & Accusatif (für, durch, gegen...) \\
\hline
masculin & mit \textbf{dem} & für \textbf{den} \\
\hline
féminin & mit \textbf{der} & für \textbf{die} \\
\hline
neutre & mit \textbf{dem} & für \textbf{das} \\
\hline
pluriel & mit \textbf{denen} & für \textbf{die} \\
\hline
\end{tabular}
\end{center}

\textbf{Ordre des mots :} la relative est une subordonnée : le verbe conjugué va \textbf{à la fin}.\\
\all{Die Mine, in der er arbeitete, war alt.} « La mine dans laquelle il travaillait était vieille. »

\textbf{Cas particulier :} si l'antécédent est une \textbf{chose} (pas une personne), on peut remplacer \emph{préposition + relatif} par \cle{wo(r)- + préposition} : \all{womit} (avec quoi), \all{worüber} (sur quoi), \all{wovon} (de quoi), \all{worauf} (sur quoi, à quoi). Le \all{-r-} s'ajoute devant une voyelle.\\
\all{Das Werkzeug, womit ( = mit dem) die Bergleute arbeiteten, war schwer.} « L'outil avec lequel les mineurs travaillaient était lourd. »
\end{propriete}

\begin{exemplebox}[Exemples traduits]
\all{Die Mine, in der die Arbeiter arbeiteten, war tief.} « La mine dans laquelle les ouvriers travaillaient était profonde. »\\
\all{Der Rohstoff, für den man viel Geld bezahlt, ist Gold.} « La matière première pour laquelle on paie beaucoup d'argent, c'est l'or. »\\
\all{Die Bergleute, von denen wir sprechen, kamen aus dem Osten.} « Les mineurs dont nous parlons venaient de l'Est. »\\
\all{Die Frage, worüber alle diskutieren, bleibt offen.} « La question dont tout le monde discute reste ouverte. »\\
\all{Kamerun ist ein Land, in dem Gold und Bauxit abgebaut werden.} « Le Cameroun est un pays où l'on extrait de l'or et de la bauxite. »
\end{exemplebox}

\begin{attention}[Pièges pour les francophones]
\begin{itemize}
\item En français, « dont » cache souvent une préposition : « les mineurs \emph{dont} nous parlons » = \all{die Bergleute, \textbf{von denen} wir sprechen} (car \all{sprechen von} + datif). Il faut connaître la construction du verbe allemand : \all{sprechen von}, \all{warten auf} (+ acc.), \all{sich freuen über} (+ acc.), \all{arbeiten in} (+ dat.).
\item Ne jamais séparer la préposition du relatif : faux \all{die Mine, die er in arbeitete} ; correct : \all{die Mine, \textbf{in der} er arbeitete}.
\item \all{wo(r)-} ne s'emploie que pour les \textbf{choses}, jamais pour les personnes : \all{der Arbeiter, mit dem ich sprach} (jamais \all{womit} pour une personne).
\end{itemize}
\end{attention}

\begin{exoresolu}[Relatives avec préposition]
Énoncé : combinez les deux phrases avec une relative à préposition.
\begin{enumerate}
\item Die Mine war gefährlich. Die Arbeiter arbeiteten in der Mine.
\item Der Rohstoff ist teuer. Man bezahlt viel Geld für den Rohstoff.
\end{enumerate}
\tcblower
Solution détaillée :
\begin{enumerate}
\item La préposition est \all{in} (+ datif) ; antécédent \all{die Mine} (fém. sing.) $\rightarrow$ \all{in der} : \all{Die Mine, in der die Arbeiter arbeiteten, war gefährlich.} « La mine dans laquelle travaillaient les ouvriers était dangereuse. »
\item La préposition est \all{für} (+ accusatif) ; antécédent \all{der Rohstoff} (masc. sing.) $\rightarrow$ \all{für den} : \all{Der Rohstoff, für den man viel Geld bezahlt, ist teuer.} « La matière première pour laquelle on paie beaucoup d'argent est chère. »
\end{enumerate}
Méthode : (1) repérer la préposition et son cas ; (2) accorder le relatif avec l'antécédent ; (3) envoyer le verbe conjugué à la fin de la relative.
\end{exoresolu}

\coursec{Commentaire guidé \& production écrite}

\courssub{Méthode du commentaire de texte au Bac}

\begin{methode}[Commenter un texte documentaire en quatre étapes]
\begin{enumerate}
\item \textbf{Présenter} : type de texte, thème, structure (ici : texte documentaire chronologique sur l'histoire du charbon en Allemagne).
\item \textbf{Dégager l'idée principale} : le texte montre que le progrès industriel a un prix humain et écologique.
\item \textbf{Analyser les procédés} : chronologie (\all{in der Antike, im Mittelalter, im 19. Jahrhundert, heute}), passif au prétérit pour insister sur les actions subies, chiffres et exemples concrets (2018, Zeche Zollverein).
\item \textbf{Confronter au contexte camerounais} : quelles leçons pour l'exploitation de la bauxite de Minim-Martap ou de l'or de Bétaré-Oya ?
\end{enumerate}
\end{methode}

\courssub{Production écrite guidée}

\begin{exercice}[Résumé puis courte production]
\begin{enumerate}
\item Résumez le texte en 4 ou 5 phrases allemandes, en utilisant au moins trois formes de passif au prétérit.
\item Rédigez un court paragraphe (6-8 phrases) : \all{Der Bergbau in Kamerun – Chancen und Gefahren} (« L'exploitation minière au Cameroun — chances et dangers »). Utilisez : deux passifs au prétérit, une relative avec préposition, et le vocabulaire du thème (\all{das Gold, der Rohstoff, die Umwelt, der Arbeiter}).
\end{enumerate}
\end{exercice}

\begin{corrige}[Production écrite guidée — proposition de résumé]
\all{Schon in der Antike wurde in Deutschland Erz abgebaut, und im Mittelalter wurden viele Minen eröffnet. Im 19. Jahrhundert wurde in der Ruhrregion Kohle gefördert; die Arbeit der Bergleute war hart und gefährlich. Die Umwelt wurde stark verschmutzt: Flüsse und Wälder wurden zerstört. Im Jahr 2018 wurde die letzte Zeche geschlossen, und die Zeche Zollverein wurde zum Weltkulturerbe erklärt. Heute stellt sich die Frage, wie man Rohstoffe gewinnen kann, ohne die Natur zu zerstören.}

Proposition de paragraphe : \all{Auch in Kamerun werden Rohstoffe abgebaut: In Bétaré-Oya wurde Gold gefunden, und in Minim-Martap wird Bauxit gefördert. Der Bergbau, von dem viele Familien leben, gibt den Arbeitern Arbeit und Geld. Aber die Natur, in der die Menschen leben, wird oft zerstört: Flüsse wurden verschmutzt und Wälder wurden abgeholzt. Die Regierung muss die Umwelt schützen. Man muss Rohstoffe gewinnen, ohne die Natur zu zerstören.}
\end{corrige}

\coursec{Production orale guidée \& réinvestissement}

\begin{experience}[Débat : Faut-il exploiter les minerais du Cameroun ?]
La classe se divise en deux groupes.
\begin{itemize}
\item \textbf{Groupe A} (les « pour ») : le ministre des Mines. Arguments : emplois, argent pour l'État, développement. Expressions utiles : \all{Ich bin der Meinung, dass...} « Je suis d'avis que... » ; \all{Der Bergbau schafft Arbeitsplätze.} « L'exploitation minière crée des emplois. »
\item \textbf{Groupe B} (les « contre ») : les habitants du village et les écologistes. Arguments : pollution, destruction des forêts, accidents. Expressions utiles : \all{Ich bin dagegen, weil...} « Je suis contre, parce que... » ; \all{Die Umwelt wurde schon oft zerstört.} « L'environnement a déjà souvent été détruit. »
\end{itemize}
Chaque élève parle au moins une fois, avec une phrase au passif au prétérit ou une relative avec préposition. Un élève-secrétaire note les arguments au tableau ; la classe vote ensuite.
\end{experience}

\begin{exercice}[Pour aller plus loin]
Répondez en allemand, avec une phrase complète au passif au prétérit :
\begin{enumerate}
\item Wann wurde die letzte Zeche in Deutschland geschlossen?
\item Was wurde in der Ruhrregion gefördert?
\item Was wurde von der UNESCO zum Weltkulturerbe erklärt?
\end{enumerate}
\end{exercice}

\begin{corrige}[Exercice « Pour aller plus loin »]
\begin{enumerate}
\item \all{Die letzte Zeche wurde im Jahr 2018 geschlossen.} « La dernière mine a été fermée en 2018. »
\item \all{In der Ruhrregion wurde Kohle gefördert.} « Dans la Ruhr, on extrayait du charbon. »
\item \all{Die Zeche Zollverein in Essen wurde von der UNESCO zum Weltkulturerbe erklärt.} « La mine Zollverein à Essen a été classée par l'UNESCO au patrimoine mondial. »
\end{enumerate}
\end{corrige}

\begin{aretenir}[L'essentiel de la leçon]
\begin{itemize}
\item Le vocabulaire s'apprend avec article et pluriel : \all{das Erz (-e), der Bergbau, die Mine (-n), das Gold, der Rohstoff (-e), der Arbeiter (-)}.
\item Le passif au prétérit = \cle{wurde/wurden + Partizip II en fin de phrase} ; il exprime une action subie, comme « a été + participe » en français.
\item \all{wurde} = singulier, \all{wurden} = pluriel ; l'accord se fait avec le sujet grammatical ; l'agent est introduit par \all{von} + datif.
\item La relative avec préposition : préposition + relatif au cas exigé (\all{in der, für den, von denen}), verbe à la fin ; pour les choses, forme possible en \all{wo(r)-} (\all{womit, worüber}).
\item Le texte « Der Bergbau in Deutschland » suit un plan chronologique : Antiquité $\rightarrow$ Moyen Âge $\rightarrow$ industrialisation $\rightarrow$ aujourd'hui, et invite à comparer avec la situation minière du Cameroun.
\end{itemize}
\end{aretenir}

\coursec{Compréhension globale et détaillée}

Dans la section précédente, nous avons découvert le thème de l'extraction des minerais à travers une mise en situation et un premier contact avec le texte support. Nous allons maintenant apprendre à \cle{comprendre} ce texte de manière méthodique : d'abord une lecture globale (l'idée générale), puis une lecture détaillée (dates, lieux, chiffres, chronologie), et enfin l'identification du point de vue de l'auteur. C'est exactement la démarche exigée à l'épreuve de compréhension de l'écrit du baccalauréat.

\courssub{La méthode de lecture : du général au particulier}

Avant de répondre à la moindre question, un bon lecteur suit toujours le même chemin. Observez d'abord la méthode, puis appliquez-la au texte.

\begin{methode}[Comment aborder un texte de compréhension écrite]
\begin{enumerate}
\item \textbf{Lire le titre et les sous-titres.} Ils annoncent le thème et le plan du texte. Exemple : \all{Der Bergbau in Deutschland — Gestern und heute} annonce un texte historique et comparatif.
\item \textbf{Survoler le texte} (\all{der Text im Schnelldurchgang}) : lire la première et la dernière phrase de chaque paragraphe pour saisir l'idée générale, sans s'arrêter sur les mots inconnus.
\item \textbf{Lire les questions} avant la lecture détaillée : elles guident l'œil vers les informations à repérer.
\item \textbf{Lire attentivement} en soulignant les dates, les lieux, les chiffres et les connecteurs logiques (\all{zuerst}, \all{dann}, \all{schließlich}).
\item \textbf{Répondre en phrases complètes}, en reformulant avec ses propres mots, jamais en recopiant de longs passages.
\end{enumerate}
\end{methode}

\courssub{Le texte support : lecture guidée}

Relisons le texte support de la leçon, cette fois avec un glossaire complet pour préparer la lecture détaillée.

\begin{textebox}[Der Bergbau im Ruhrgebiet — ein Kapitel deutscher Geschichte]
Der Bergbau hat das Ruhrgebiet über 150 Jahre lang geprägt. Im 19. Jahrhundert wurden dort riesige Kohlevorkommen entdeckt, und bald wurden die ersten Zechen eröffnet. Kohle, Eisenerz und andere Rohstoffe wurden aus der Erde gefördert und in die Fabriken transportiert, wo sie für die Stahlproduktion gebraucht wurden.

Um 1950 arbeiteten fast 500\,000 Bergleute im Ruhrgebiet. Die Kohle wurde damals im Untertagebau gewonnen: Die Arbeiter stiegen jeden Tag Hunderte Meter tief in die Schächte hinab. Die Arbeit war hart und gefährlich. Viele Bergleute wurden bei Grubenunglücken verletzt oder getötet, und viele litten später an Staublunge, einer schweren Krankheit.

Seit den 1960er Jahren wurde der Bergbau immer weniger wichtig, weil billigere Kohle aus dem Ausland importiert wurde. Eine Zeche nach der anderen wurde geschlossen. Im Dezember 2018 wurde die letzte Zeche Deutschlands, die Zeche Prosper-Haniel in Bottrop, endgültig stillgelegt. Damit wurde ein ganzes Kapitel der Industriegeschichte beendet.

Heute werden viele alte Zechen als Museen genutzt. Die Zeche Zollverein in Essen zum Beispiel wurde von der UNESCO zum Weltkulturerbe erklärt und wird jedes Jahr von mehr als einer Million Touristen besucht. Der Bergbau wird zwar nicht mehr betrieben, aber er wird nicht vergessen: Er bleibt ein wichtiger Teil der Identität der Region.

(Environ 220 mots)
\end{textebox}

\textbf{Glossaire :}

\begin{center}
\begin{tabularx}{\linewidth}{|X|X|}
\hline
\rowcolor{popblueL} \textbf{Deutsch} & \textbf{Français} \\
\hline
der Bergbau (Sg.) & l'exploitation minière \\
\hline
das Ruhrgebiet, -e & le bassin de la Ruhr \\
\hline
prägen (hat geprägt) & marquer, façonner \\
\hline
das Kohlevorkommen, - & le gisement de charbon \\
\hline
die Zeche, -n & la mine de charbon, la fosse \\
\hline
das Eisenerz, -e & le minerai de fer \\
\hline
der Rohstoff, -e & la matière première \\
\hline
fördern (hat gefördert) & extraire \\
\hline
der Bergmann, -leute / die Bergleute (Pl.) & le mineur / les mineurs \\
\hline
der Untertagebau (Sg.) & l'exploitation souterraine \\
\hline
der Schacht, -e & le puits (de mine) \\
\hline
das Grubenunglück, -e & l'accident minier \\
\hline
die Staublunge, -n & la silicose (maladie des poumons) \\
\hline
schließen (wurde geschlossen) & fermer \\
\hline
stilllegen (wurde stillgelegt) & fermer définitivement \\
\hline
das Weltkulturerbe (Sg.) & le patrimoine mondial de l'UNESCO \\
\hline
\end{tabularx}
\end{center}

\begin{attention}[Orthographe : \all{stilllegen} avec deux « l »]
Le verbe \all{stilllegen} (mettre hors service, fermer définitivement) s'écrit avec \textbf{deux l} : il est composé de \all{still} (immobile) + \all{legen} (poser). Participe II : \all{stillgelegt}. Ne pas confondre avec \all{schließen} (fermer, de façon générale) : une boutique \all{schließt} le soir, mais une mine est \all{stillgelegt} pour toujours.
\end{attention}

\courssub{Compréhension globale : le QCM}

La compréhension globale vise l'\cle{idée générale} du texte. Au baccalauréat, elle est souvent testée par des questions à choix multiples (QCM, \all{Multiple-Choice-Fragen}). Voici le modèle de l'exercice, avec son corrigé détaillé.

\begin{exoresolu}[QCM — L'idée générale du texte]
\textbf{Énoncé.} Cochez la bonne réponse (\all{Kreuzen Sie die richtige Antwort an}).

\textbf{1.} Worum geht es in diesem Text?
\begin{itemize}
\item[a)] Um die Entdeckung von Gold in Deutschland.
\item[b)] Um die Geschichte und das Ende des Bergbaus im Ruhrgebiet.
\item[c)] Um den Tourismus in der Stadt Essen.
\end{itemize}

\textbf{2.} Welche Rohstoffe wurden gefördert?
\begin{itemize}
\item[a)] Nur Gold und Silber.
\item[b)] Kohle und Eisenerz.
\item[c)] Nur Diamanten.
\end{itemize}

\textbf{3.} Wann wurde die letzte Zeche geschlossen?
\begin{itemize}
\item[a)] 1950.
\item[b)] In den 1960er Jahren.
\item[c)] Im Dezember 2018.
\end{itemize}

\tcblower
\textbf{Solution détaillée.}

\textbf{1. b)} — Le texte retrace toute l'histoire du bassin de la Ruhr, de la découverte du charbon au XIX\textsuperscript{e} siècle jusqu'à la fermeture de la dernière fosse. Les options a) et c) ne correspondent qu'à des détails ou à des thèmes absents.

\textbf{2. b)} — Premier paragraphe : \all{Kohle, Eisenerz und andere Rohstoffe wurden aus der Erde gefördert.} « Le charbon, le minerai de fer et d'autres matières premières étaient extraits de la terre. »

\textbf{3. c)} — Troisième paragraphe : \all{Im Dezember 2018 wurde die letzte Zeche Deutschlands ... endgültig stillgelegt.} « En décembre 2018, la dernière mine d'Allemagne a été fermée définitivement. »

\textbf{Méthode :} pour chaque item, repérez dans le texte la phrase qui justifie la réponse. Une réponse de QCM sans preuve textuelle n'est qu'une intuition.
\end{exoresolu}

\begin{exemplebox}[Formuler les questions de compréhension]
Observez la structure des questions typiques et leur traduction :

\all{Welche Rohstoffe wurden gefördert?} → « Quelles matières premières ont été extraites ? »

\all{Wann wurde die erste Zeche eröffnet?} → « Quand la première mine a-t-elle été ouverte ? »

\all{Wie viele Bergleute arbeiteten um 1950 im Ruhrgebiet?} → « Combien de mineurs travaillaient dans la Ruhr vers 1950 ? »

\all{Warum wurde der Bergbau immer weniger wichtig?} → « Pourquoi l'exploitation minière est-elle devenue de moins en moins importante ? »

Remarquez l'ordre des mots : mot interrogatif + verbe conjugué en 2\textsuperscript{e} position + sujet. Au passif prétérit : \all{wurden} + sujet + participe II en fin de phrase.
\end{exemplebox}

\courssub{Compréhension détaillée : questions ouvertes}

Les questions ouvertes exigent une réponse rédigée en phrase complète. Prenons la question centrale de la leçon :

\textbf{\all{Warum ist der Bergbau für die Industrie wichtig?}} « Pourquoi l'exploitation minière est-elle importante pour l'industrie ? »

\begin{methode}[Répondre à une question ouverte en trois étapes]
\begin{enumerate}
\item \textbf{Repérer} le passage du texte qui contient la réponse (ici : le premier paragraphe).
\item \textbf{Reformuler} sans recopier mot à mot : remplacer les mots du texte par des synonymes (\all{fördern} → \all{aus der Erde holen}).
\item \textbf{Rédiger} une phrase complète avec un connecteur de cause : \all{weil} + verbe en fin de phrase.
\end{enumerate}
\textbf{Réponse modèle :} \all{Der Bergbau ist für die Industrie wichtig, weil er die Rohstoffe liefert, die man für die Produktion braucht, zum Beispiel Kohle und Eisenerz für die Stahlproduktion.} « L'exploitation minière est importante pour l'industrie parce qu'elle fournit les matières premières nécessaires à la production, par exemple le charbon et le minerai de fer pour la production d'acier. »
\end{methode}

\begin{exercice}[À vous de répondre]
Répondez par des phrases complètes (\all{Antworten Sie mit vollständigen Sätzen}) :

\begin{enumerate}
\item \all{Wohin wurden die Rohstoffe transportiert?}
\item \all{Warum war die Arbeit der Bergleute gefährlich?}
\item \all{Was passiert heute mit den alten Zechen?}
\end{enumerate}
\end{exercice}

\begin{corrige}[Exercice — questions ouvertes]
\begin{enumerate}
\item \all{Die Rohstoffe wurden in die Fabriken transportiert.} « Les matières premières étaient transportées vers les usines. » (On préfère \all{wohin} « vers où » à \all{wo} « où », car il s'agit d'un mouvement.)
\item \all{Die Arbeit war gefährlich, weil viele Bergleute bei Grubenunglücken verletzt oder getötet wurden.} « Le travail était dangereux parce que beaucoup de mineurs étaient blessés ou tués dans des accidents miniers. »
\item \all{Heute werden viele alte Zechen als Museen genutzt.} « Aujourd'hui, beaucoup d'anciennes mines sont utilisées comme musées. »
\end{enumerate}
\end{corrige}

\courssub{Recherche d'informations précises : dates, lieux, chiffres}

La lecture détaillée consiste aussi à extraire des données exactes. Entraînons-nous à remplir une fiche d'informations à partir du texte.

\begin{center}
\begin{tabularx}{\linewidth}{|X|X|}
\hline
\rowcolor{popblueL} \textbf{Information recherchée} & \textbf{Réponse dans le texte} \\
\hline
Durée de l'influence du Bergbau & über 150 Jahre \\
\hline
Siècle de la découverte du charbon & im 19. Jahrhundert \\
\hline
Nombre de mineurs vers 1950 & fast 500\,000 \\
\hline
Type d'exploitation du charbon & der Untertagebau \\
\hline
Début du déclin & seit den 1960er Jahren \\
\hline
Fermeture de la dernière mine & Dezember 2018, Zeche Prosper-Haniel in Bottrop \\
\hline
Musée célèbre & die Zeche Zollverein in Essen \\
\hline
Visiteurs par an & mehr als eine Million Touristen \\
\hline
\end{tabularx}
\end{center}

\begin{attention}[Le piège de « seit »]
En allemand, \all{seit} (depuis) s'emploie avec le \textbf{présent} quand l'action continue encore, alors que le français hésite entre présent et passé composé :

\all{Seit 2018 wird dort keine Kohle mehr gefördert.} → « Depuis 2018, on n'y extrait plus de charbon. » (l'allemand garde le présent)

\all{Seit den 1960er Jahren wurde der Bergbau immer weniger wichtig.} → « Depuis les années 1960, l'exploitation minière devenait de moins en moins importante. » (ici le prétérit est possible car le processus est décrit comme révolu)

\textbf{Erreur typique du francophone :} traduire \all{seit} par « pendant » ou le confondre avec \all{vor} (il y a). Retenez : \all{seit 1950} = « depuis 1950 » ; \all{vor 70 Jahren} = « il y a 70 ans ». Autre piège : \all{seit} exige le \textbf{datif} : \all{seit den 1960er Jahren}, \all{seit dem 19. Jahrhundert}.
\end{attention}

\begin{definition}[Deux types d'exploitation minière]
\textbf{der Tagebau} (pl. \all{die Tagebaue}) : la mine à ciel ouvert. On enlève la terre en surface pour atteindre le minerai. \all{Braunkohle wird oft im Tagebau gefördert.} « Le lignite est souvent extrait à ciel ouvert. »

\textbf{der Untertagebau} (pl. \all{die Untertagebaue}) : la mine souterraine. Les mineurs descendent dans des puits (\all{die Schächte}) pour extraire le minerai en profondeur. \all{Die Steinkohle im Ruhrgebiet wurde im Untertagebau gewonnen.} « La houille de la Ruhr était extraite en souterrain. »

Au Cameroun, on parlerait plutôt d'exploitation artisanale à ciel ouvert pour l'or de la région de l'Est ; en Allemagne, le charbon de la Ruhr était un cas typique d'\all{Untertagebau}.
\end{definition}

\courssub{Mise en ordre chronologique des événements}

Reconstituer la chronologie (\all{die Chronologie}) est un exercice classique : il prouve que vous avez compris l'enchaînement des faits. Les connecteurs à repérer sont : \all{zuerst} (d'abord), \all{dann / danach} (ensuite), \all{seit} (depuis), \all{schließlich} (finalement), \all{heute} (aujourd'hui).

\begin{exoresolu}[Remettre les événements dans l'ordre]
\textbf{Énoncé.} \all{Bringen Sie die Ereignisse in die richtige Reihenfolge (1–5).}

\begin{itemize}
\item[a)] Die letzte Zeche wird stillgelegt.
\item[b)] Riesige Kohlevorkommen werden im Ruhrgebiet entdeckt.
\item[c)] Die Zeche Zollverein wird zum Weltkulturerbe erklärt.
\item[d)] Fast 500\,000 Bergleute arbeiten im Ruhrgebiet.
\item[e)] Der Bergbau wird immer weniger wichtig, weil billige Kohle importiert wird.
\end{itemize}

\tcblower
\textbf{Solution détaillée.}

\begin{enumerate}
\item \textbf{b)} — La découverte des gisements au XIX\textsuperscript{e} siècle ouvre le récit (paragraphe 1).
\item \textbf{d)} — Vers 1950, l'apogée : près de 500\,000 mineurs (paragraphe 2).
\item \textbf{e)} — Depuis les années 1960, le déclin à cause des importations (paragraphe 3).
\item \textbf{a)} — En décembre 2018, la fermeture de la dernière fosse (paragraphe 3).
\item \textbf{c)} — Aujourd'hui, la reconversion : la Zeche Zollverein classée par l'UNESCO (paragraphe 4).
\end{enumerate}

Ordre correct : \textbf{b – d – e – a – c}. Astuce : suivez les marqueurs temporels du texte (\all{im 19. Jahrhundert} → \all{um 1950} → \all{seit den 1960er Jahren} → \all{2018} → \all{heute}).
\end{exoresolu}

\courssub{Le point de vue de l'auteur : critique, neutre ou optimiste ?}

Dernière étape de la lecture détaillée : dégager l'\cle{attitude de l'auteur} (\all{die Haltung des Autors}). Trois indices à observer :

\begin{itemize}
\item \textbf{Le lexique} : mots négatifs (\all{hart}, \all{gefährlich}, \all{getötet}, \all{Krankheit}) ou positifs (\all{Weltkulturerbe}, \all{besucht}).
\item \textbf{Les faits rapportés} : l'auteur cite-t-il les deux faces du sujet (richesse industrielle \emph{et} souffrance des mineurs) ?
\item \textbf{La conclusion} : le texte finit-il sur une note de regret, d'espoir ou de constat ?
\end{itemize}

Dans notre texte, l'auteur mentionne objectivement les chiffres et les dates (ton \textbf{neutre} et informatif), mais il souligne aussi les dangers du métier et la maladie des mineurs (ton \textbf{critique}), avant de terminer sur la mémoire préservée et le tourisme culturel (note plutôt \textbf{optimiste}). On dira donc : \all{Der Autor ist sachlich, aber auch kritisch; am Ende klingt der Text versöhnlich und optimistisch.} « L'auteur est factuel, mais aussi critique ; à la fin, le texte sonne de façon apaisée et optimiste. »

\begin{center}
\begin{tabularx}{\linewidth}{|X|X|X|}
\hline
\rowcolor{popblueL} \textbf{Point de vue} & \textbf{Indice dans le texte} & \textbf{Français} \\
\hline
neutral / sachlich & chiffres et dates (500\,000 Bergleute, 2018) & neutre / factuel \\
\hline
kritisch & \all{hart und gefährlich}, \all{Staublunge} & critique \\
\hline
optimistisch & \all{Weltkulturerbe}, \all{wird nicht vergessen} & optimiste \\
\hline
\end{tabularx}
\end{center}

\courssub{La structure du texte : un organigramme}

Pour bien mémoriser le plan du texte, visualisons sa structure en trois parties : introduction, corps, conclusion.

\begin{popfigure}
\begin{center}
\begin{tikzpicture}[
  box/.style={draw=popblue, fill=popblueL, rounded corners=4pt, align=center, text width=4.2cm, minimum height=1.1cm, font=\small},
  arr/.style={-{Stealth[length=3mm]}, thick, popdark}
]
\node[box, align=center] (intro) at (0,0){\textbf{Einleitung}\\Découverte du charbon au XIX\textsuperscript{e} s.};
\node[box, fill=poporangeL, draw=poporange, align=center] (corps) at (0,-2.2){\textbf{Hauptteil}\\Apogée (1950), dangers du métier,\\déclin et fermeture (2018)};
\node[box, fill=popgreenL, draw=popgreen, align=center] (concl) at (0,-4.6){\textbf{Schluss}\\Reconversion : musées et\\patrimoine mondial};
\draw[arr] (intro) -- (corps);
\draw[arr] (corps) -- (concl);
\end{tikzpicture}
\end{center}
\legende{Organigramme de la structure du texte : Einleitung – Hauptteil – Schluss.}
\end{popfigure}

\begin{savaistu}[D'où vient le mot « Bergbau » ?]
Le mot \all{der Bergbau} est un mot composé, très typique de l'allemand : \all{der Berg} (la montagne) + \all{der Bau} (la construction, l'exploitation). Littéralement, c'est donc « l'exploitation de la montagne » ! De la même famille : \all{der Bergmann} (litt. « l'homme de la montagne » = le mineur), \all{das Bergwerk} (la mine, la fosse), \all{bergen} (extraire, récupérer, sauver). Saviez-vous que les mineurs allemands se saluent traditionnellement par \all{Glückauf!} — littéralement « chance pour la remontée » — un vœu né de la peur des accidents ? Ce salut est encore utilisé aujourd'hui dans la Ruhr, par exemple par les supporters du club de football Schalke 04, le club des mineurs.
\end{savaistu}

\begin{aretenir}[Les cinq réflexes du bon lecteur]
\begin{enumerate}
\item Titre et sous-titres d'abord, survol du texte ensuite.
\item QCM : toujours justifier par une phrase du texte.
\item Questions ouvertes : phrase complète + connecteur (\all{weil} + verbe final).
\item Chronologie : suivre les marqueurs temporels (\all{im 19. Jahrhundert}, \all{um 1950}, \all{seit}, \all{heute}).
\item Point de vue de l'auteur : observer le lexique, les faits cités et la conclusion.
\end{enumerate}
\end{aretenir}

Cette double lecture — globale puis détaillée — vous a permis de construire une compréhension solide du texte. Nous allons maintenant, dans la section suivante, enrichir notre lexique thématique autour du monde minier, avant d'étudier les deux points de grammaire du chapitre : le passif au prétérit et les propositions relatives avec préposition.

\coursec{Lexique thématique par champs}

Après avoir compris le texte dans son ensemble puis dans le détail, nous allons maintenant organiser et approfondir le vocabulaire de l'extraction des minerais. Un bon vocabulaire ne s'apprend pas mot par mot : il s'apprend par \cle{champs lexicaux}, c'est-à-dire par familles de mots liés à un même thème. Nous allons donc classer le lexique du texte en cinq champs : les matières, les lieux, les acteurs, les verbes d'action et les impacts sur l'environnement. Pour chaque mot, retiens toujours trois informations : l'\cle{article} (le genre), le \cle{pluriel} et une \cle{phrase d'exemple}.

\begin{methode}[Créer une fiche de vocabulaire efficace]
Pour chaque mot nouveau, construis une fiche en cinq étapes :
\begin{enumerate}
\item Écris le mot avec son \textbf{article} : \all{das Erz} (jamais \all{Erz} seul).
\item Note le \textbf{pluriel} : \all{die Erze}.
\item Ajoute la \textbf{traduction} française : « le minerai ».
\item Rédige une \textbf{phrase d'exemple} tirée du texte ou inventée : \all{Das Erz wird in der Mine gefördert.}
\item Ajoute les \textbf{mots de la même famille} (dérivés, composés) : \all{der Bergbau}, \all{der Bergmann}, \all{das Bergwerk}, \all{bergmännisch}.
\end{enumerate}
Révise tes fiches en couvrant la colonne allemande et en retrouvant le mot à partir du français, puis l'inverse.
\end{methode}

\courssub{Champ 1 : Les matières — die Rohstoffe}

Observons d'abord quelques phrases du texte support :

\all{In der Mine wird Kupfer gefördert.} « Dans la mine, on extrait du cuivre. »

\all{Gold und Salz sind wertvolle Rohstoffe.} « L'or et le sel sont des matières premières précieuses. »

\all{Die Kohle wurde früher im Ruhrgebiet abgebaut.} « Le charbon était autrefois exploité dans la Ruhr. »

On remarque que les noms de matières sont souvent \textbf{neutres} (\all{das Gold}, \all{das Erz}, \all{das Salz}, \all{das Kupfer}, \all{das Eisen}) et qu'ils s'emploient généralement \textbf{sans article} quand on parle de la matière en général, comme en français : \all{Gold ist teuer.} « L'or est cher. » \all{Die Kohle} (féminin) est une exception notable.

\begin{definition}[Les matières et les minerais]
\begin{itemize}
\item \all{das Erz, die Erze} : le minerai
\item \all{das Gold} (pas de pluriel au sens matière) : l'or
\item \all{die Kohle} (plur. \all{die Kohlen} seulement pour « types de charbon » ou « morceaux ») : le charbon
\item \all{das Salz, die Salze} : le sel
\item \all{das Kupfer} (pas de pluriel au sens matière) : le cuivre
\item \all{das Eisen} (pas de pluriel au sens matière) : le fer
\item \all{der Rohstoff, die Rohstoffe} : la matière première
\item \all{das Metall, die Metalle} : le métal
\end{itemize}
\end{definition}

\begin{attention}[Mots sans pluriel (noms de matières)]
Les noms de matières comme \all{das Gold}, \all{das Kupfer}, \all{das Eisen}, \all{das Silber} n'ont en général \textbf{pas de pluriel}, car ce sont des noms non dénombrables (\emph{Massennamen}). On ne dit pas \all{viele Golde} ; on dit \all{viel Gold} (« beaucoup d'or »). De même en français on ne dit pas « des ors » dans ce sens. \all{Die Kohle} fonctionne comme \all{le charbon} (singulier collectif), le pluriel \all{die Kohlen} désigne des sortes ou des morceaux.
\end{attention}

\courssub{Champ 2 : Les lieux — die Orte des Bergbaus}

\begin{definition}[Les lieux de l'extraction]
\begin{itemize}
\item \all{die Mine, die Minen} : la mine (souterraine)
\item \all{der Stollen, die Stollen} : la galerie (horizontale)
\item \all{der Schacht, die Schächte} : le puits de mine (vertical)
\item \all{der Tagebau, die Tagebaue} : l'exploitation à ciel ouvert
\item \all{die Hütte, die Hütten} : l'usine métallurgique, la fonderie
\item \all{das Bergwerk, die Bergwerke} : la mine (l'ensemble de l'exploitation)
\item \all{der Bergbau} (pas de pluriel) : l'industrie minière, l'activité minière
\end{itemize}
\end{definition}

\begin{attention}[Faux amis : \all{der Schacht} et \all{die Hütte}]
\begin{itemize}
\item \all{Der Schacht} ne signifie pas « le chat » ! C'est le \textbf{puits de mine}, le conduit vertical par lequel on descend. Le chat se dit \all{die Katze}.
\item \all{Die Hütte} : dans le contexte minier/industriel, ce n'est pas une « hutte » (\all{die Hütte} au sens cabane existe aussi), mais une \textbf{usine de traitement} (fonderie, aciérie). \all{Die Hüttenwerke} = les usines sidérurgiques.
\end{itemize}
\end{attention}

\begin{exemplebox}[Les lieux dans des phrases]
\all{Die Arbeiter fahren mit dem Fahrstuhl in den Schacht.} « Les ouvriers descendent dans le puits avec l'ascenseur. »

\all{Im Stollen ist es dunkel und kalt.} « Dans la galerie, il fait sombre et froid. »

\all{Im Tagebau wird die Kohle an der Oberfläche abgebaut.} « Dans l'exploitation à ciel ouvert, le charbon est extrait en surface. »

\all{Das Erz wurde in die Hütte transportiert.} « Le minerai a été transporté à l'usine. »
\end{exemplebox}

\courssub{Champ 3 : Les acteurs — die Menschen im Bergbau}

\begin{definition}[Les personnes et les organisations]
\begin{itemize}
\item \all{der Arbeiter, die Arbeiter} : l'ouvrier
\item \all{die Arbeiterin, die Arbeiterinnen} : l'ouvrière
\item \all{der Bergmann, die Bergleute} (plur. principal) / \all{die Bergmänner} : le mineur
\item \all{der Ingenieur, die Ingenieure} : l'ingénieur
\item \all{die Ingenieurin, die Ingenieurinnen} : l'ingénieure
\item \all{die Gewerkschaft, die Gewerkschaften} : le syndicat
\item \all{der Unternehmer, die Unternehmer} : le chef d'entreprise, l'entrepreneur
\item \all{der Experte, die Experten} / \all{die Expertin, die Expertinnen} : l'expert(e)
\end{itemize}
\end{definition}

Remarque la formation du féminin régulier : on ajoute \all{-in} au masculin, et le pluriel devient \all{-innen} : \all{der Ingenieur} → \all{die Ingenieurin} → \all{die Ingenieurinnen}. \all{Der Bergmann} possède deux pluriels : \all{die Bergleute} (le plus courant, « les mineurs » comme collectif) et \all{die Bergmänner} (individus). Retiens surtout \all{die Bergleute}.

\courssub{Champ 4 : Les verbes d'action — die Verben}

\begin{definition}[Les verbes de l'extraction et du traitement]
\begin{itemize}
\item \all{fördern} (faible) : extraire, remonter (le minerai) — \all{er fördert, er förderte, er hat gefördert}
\item \all{abbauen} (séparable, faible) : exploiter, extraire (une ressource) — \all{Die Firma baut Kohle ab.}
\item \all{schmelzen} (fort) : fondre — \all{er schmilzt, er schmolz, er hat geschmolzen}
\item \all{transportieren} (faible, en \all{-ieren}) : transporter — \all{er hat transportiert} (pas de \all{ge-})
\item \all{verarbeiten} (inséparable, faible) : transformer, traiter — \all{er verarbeitet, er verarbeitete, er hat verarbeitet}
\item \all{exportieren} (faible, en \all{-ieren}) : exporter — \all{er hat exportiert}
\item \all{gewinnen} (fort) : extraire, obtenir, gagner — \all{er gewinnt, er gewann, er hat gewonnen}
\end{itemize}
\end{definition}

\begin{attention}[Particule séparable ou inséparable ? — Piège majeur]
Ne confonds pas \all{abbauen} (séparable : \all{Man baut das Erz ab.}, Partizip: \all{abgebaut}) et \all{verarbeiten} (inséparable : \all{Man verarbeitet das Erz.}, Partizip: \all{verarbeitet} — \textbf{pas de ge-}). Le préfixe \all{ver-} ne se sépare jamais et ne prend jamais de \all{ge-} au participe. Autres préfixes inséparables fréquents : \all{be-, ent-, er-, ge-, miss-, zer-}.
\end{attention}

\begin{exemplebox}[Exemples traduits — Passif et actif]
\all{Die Förderung von Kupfer ist teuer.} « L'extraction de cuivre est coûteuse. » (Nominalisation)

\all{Früher wurde hier Kohle gefördert.} « Autrefois, on extrayait du charbon ici. » (Passif Präteritum)

\all{Das Erz wird in der Hütte geschmolzen.} « Le minerai est fondu dans l'usine. » (Passif Présent)

\all{Die Rohstoffe wurden in den Hafen transportiert.} « Les matières premières ont été transportées au port. » (Passif Präteritum)

\all{In dieser Fabrik wird Kupfer verarbeitet.} « Dans cette usine, on transforme le cuivre. » (Passif Présent, verbe inséparable)
\end{exemplebox}

Remarque la famille de mots autour de \all{fördern} : le nom \all{die Förderung} (« l'extraction, la promotion »), \all{der Förderer} (« le convoyeur, le soutien »), \all{die Fördermenge} (« la quantité extraite »), \all{fordernd} (« encourageant »). Apprendre les familles de mots multiplie ton vocabulaire sans effort supplémentaire.

\courssub{Champ 5 : Les impacts — Umwelt und Zukunft}

\begin{definition}[Environnement et avenir]
\begin{itemize}
\item \all{die Umweltbelastung, die Umweltbelastungen} : la pollution, l'atteinte à l'environnement
\item \all{die Landschaftszerstörung} (gén. sg.) : la destruction du paysage
\item \all{das Recycling} (pas de pluriel) : le recyclage
\item \all{die Nachhaltigkeit} : la durabilité, le développement durable
\item \all{die Umwelt} : l'environnement
\item \all{der Abfall, die Abfälle} : le déchet
\item \all{die Verschmutzung, die Verschmutzungen} : la pollution
\item \all{schützen} (faible) : protéger ; \all{recyceln} (faible) : recycler ; \all{verursachen} (faible) : causer, provoquer
\end{itemize}
\end{definition}

\begin{exemplebox}[Exemples traduits]
\all{Der Tagebau führt zur Landschaftszerstörung.} « L'exploitation à ciel ouvert mène à la destruction du paysage. »

\all{Recycling ist wichtig für die Nachhaltigkeit.} « Le recyclage est important pour le développement durable. »

\all{Die Umweltbelastung durch den Bergbau ist groß.} « L'atteinte à l'environnement due à l'industrie minière est grande. »

\all{Alte Handys werden recycelt, um Gold zu gewinnen.} « Les vieux téléphones portables sont recyclés pour obtenir de l'or. » (Passif + finale)
\end{exemplebox}

\courssub{Texte d'application : le lexique en contexte}

\begin{textebox}[Rohstoffe für die Zukunft]
Der Bergbau hat eine lange Geschichte. Schon vor tausend Jahren wurde in Europa Salz und Kupfer abgebaut. Heute ist die Förderung von Rohstoffen wichtiger denn je: Ohne Metalle gibt es keine Handys, keine Autos und keine Windräder.

In einer modernen Mine arbeiten viele Menschen zusammen. Die Bergleute fahren früh am Morgen in den Schacht und fördern das Erz aus dem Stollen. Ingenieure planen die Arbeit und kontrollieren die Maschinen. Die Gewerkschaft kämpft für gute Löhne und für die Sicherheit der Arbeiter.

Aber der Bergbau hat auch negative Seiten. Der Tagebau zerstört die Landschaft, und die Umweltbelastung ist oft sehr hoch. Flüsse werden verschmutzt, Wälder verschwinden. Deshalb denken viele Experten über Alternativen nach. Eine Lösung ist das Recycling: Aus alten Computern und Handys kann man Gold, Kupfer und andere Metalle gewinnen. So wird weniger neues Erz abgebaut, und die Natur wird geschützt.

Auch in Afrika, zum Beispiel in Kamerun, spielen Rohstoffe eine große Rolle. Wenn die Rohstoffe im Land verarbeitet werden, entstehen Arbeitsplätze. Nachhaltigkeit bedeutet: Wir müssen heute so handeln, dass auch unsere Kinder noch genug Rohstoffe und eine gesunde Umwelt haben.
\end{textebox}

\textbf{Glossaire :} \all{der Bergbau} : l'industrie minière ; \all{der Schacht, die Schächte} : le puits ; \all{der Stollen, die Stollen} : la galerie ; \all{der Lohn, die Löhne} : le salaire ; \all{die Sicherheit} : la sécurité ; \all{verschmutzen} : polluer ; \all{der Wald, die Wälder} : la forêt ; \all{die Lösung, die Lösungen} : la solution ; \all{der Arbeitsplatz, die Arbeitsplätze} : l'emploi ; \all{handeln} : agir.

\textbf{Questions de compréhension :}
\begin{enumerate}
\item \all{Welche Rohstoffe wurden schon vor tausend Jahren abgebaut?}
\item \all{Was machen die Ingenieure in der Mine?}
\item \all{Welche negativen Folgen hat der Tagebau?}
\item \all{Warum ist Recycling eine Lösung?}
\item \all{Was bedeutet Nachhaltigkeit?}
\end{enumerate}

\courssub{Tableau récapitulatif du lexique}

\begin{center}
\begin{tabularx}{\linewidth}{|l|l|l|X|X|}
\hline
\rowcolor{popblueL} \textbf{Mot} & \textbf{Genre} & \textbf{Pluriel} & \textbf{Traduction} & \textbf{Exemple} \\
\hline
\all{Erz} & das & die Erze & le minerai & \all{Das Erz wird gefördert.} \\
\hline
\all{Gold} & das & --- & l'or & \all{Gold ist teuer.} \\
\hline
\all{Kohle} & die & --- (Sg. collectif) & le charbon & \all{Die Kohle wurde abgebaut.} \\
\hline
\all{Salz} & das & die Salze & le sel & \all{Salz wurde hier gewonnen.} \\
\hline
\all{Kupfer} & das & --- & le cuivre & \all{Kupfer wird verarbeitet.} \\
\hline
\all{Mine} & die & die Minen & la mine & \all{Die Mine ist sehr alt.} \\
\hline
\all{Stollen} & der & die Stollen & la galerie & \all{Der Stollen ist dunkel.} \\
\hline
\all{Schacht} & der & die Schächte & le puits & \all{Der Schacht ist 500 m tief.} \\
\hline
\all{Tagebau} & der & die Tagebaue & la mine à ciel ouvert & \all{Der Tagebau zerstört die Landschaft.} \\
\hline
\all{Hütte} & die & die Hütten & la fonderie & \all{Das Erz kommt in die Hütte.} \\
\hline
\all{Arbeiter} & der & die Arbeiter & l'ouvrier & \all{Der Arbeiter fördert Erz.} \\
\hline
\all{Bergmann} & der & die Bergleute & le mineur & \all{Der Bergmann arbeitet im Schacht.} \\
\hline
\all{Ingenieur} & der & die Ingenieure & l'ingénieur & \all{Der Ingenieur plant die Arbeit.} \\
\hline
\all{Gewerkschaft} & die & die Gewerkschaften & le syndicat & \all{Die Gewerkschaft kämpft für Löhne.} \\
\hline
\all{Umweltbelastung} & die & die Umweltbelastungen & la pollution & \all{Die Umweltbelastung ist hoch.} \\
\hline
\all{Recycling} & das & --- & le recyclage & \all{Recycling schützt die Natur.} \\
\hline
\end{tabularx}
\end{center}

\courssub{La formation du pluriel : deux modèles du texte}

Observons : \all{der Arbeiter} → \all{die Arbeiter} ; \all{der Ingenieur} → \all{die Ingenieure} ; \all{die Mine} → \all{die Minen} ; \all{die Hütte} → \all{die Hütten} ; \all{der Schacht} → \all{die Schächte}.

\begin{propriete}[Pluriel des noms en \all{-er} et noms féminins en \all{-e} / \all{-ung}]
\begin{enumerate}
\item \textbf{Noms masculins en \all{-er}} (souvent métiers, nationalités) : le pluriel est souvent \textbf{identique au singulier} (zéro marque) : \all{der Arbeiter} → \all{die Arbeiter}, \all{der Lehrer} → \all{die Lehrer}, \all{der Franzose} → \all{die Franzosen}. \textbf{Exception notable :} \all{der Ingenieur} → \all{die Ingenieure} (prend \all{-e}), \all{der Bauer} → \all{die Bauern} (Umlaut + \all{-n}).
\item \textbf{Noms féminins en \all{-e}} : forment leur pluriel en \all{-n} : \all{die Mine} → \all{die Minen}, \all{die Hütte} → \all{die Hütten}, \all{die Kohle} → \all{die Kohlen} (sortes), \all{die Lampe} → \all{die Lampen}.
\item \textbf{Noms féminins en \all{-ung}, \all{-schaft}, \all{-heit}, \all{-keit}, \all{-tion}} : prennent \all{-en} : \all{die Förderung} → \all{die Förderungen}, \all{die Gewerkschaft} → \all{die Gewerkschaften}, \all{die Sicherheit} → \all{die Sicherheiten}.
\item \textbf{Umlaut au pluriel} : nombreux monosyllabes (masc., neut., fém.) : \all{der Schacht} → \all{die Schächte}, \all{der Wald} → \all{die Wälder}, \all{das Buch} → \all{die Bücher}, \all{die Hand} → \all{die Hände}.
\end{enumerate}
\textbf{Comparaison avec le français :} le français ajoute simplement un « s » muet (la mine → les mines) ; l'allemand modifie la terminaison prononcée (Umlaut, \all{-e}, \all{-n}, \all{-en}, \all{-er}, \all{-s}), et l'article pluriel est toujours \all{die} quel que soit le genre du singulier.
\end{propriete}

\begin{attention}[Erreurs fréquentes des francophones]
\begin{itemize}
\item Ne dis jamais \all{die Arbeiters} : le pluriel de \all{der Arbeiter} est \all{die Arbeiter}, sans changement.
\item N'oublie pas l'Umlaut quand il existe : \all{der Schacht} → \all{die Sch\textbf{ä}chte}, \all{der Wald} → \all{die W\textbf{ä}lder}, \all{der Bergmann} → \all{die Bergm\textbf{ä}nner}.
\item L'article pluriel est toujours \all{die} : \all{das Erz} → \all{\textbf{die} Erze}, jamais \all{das Erze}.
\item \all{Die Kohle} (féminin) vs \all{das Gold} (neutre) : attention au genre, il conditionne l'article et les pronoms relatifs.
\end{itemize}
\end{attention}

\begin{popfigure}
\begin{center}
\begin{center}\tcbox[colback=popblue,colframe=popblue,coltext=white,arc=9pt,boxsep=4pt,left=6pt,right=6pt]{\bfseries\small der Bergbau}\end{center}
\vspace{-2pt}
\begin{tcbraster}[raster columns=3,raster equal height=rows,raster column skip=6pt,raster row skip=6pt,enhanced,blankest]
\begin{tcolorbox}[enhanced,colback=white,colframe=poporange,boxrule=0.9pt,arc=3pt,title={die Rohstoffe},fonttitle=\bfseries\scriptsize,coltitle=white,colbacktitle=poporange,left=4pt,right=4pt,top=2pt,bottom=2pt,boxsep=1pt,toptitle=1.5pt,bottomtitle=1.5pt,halign=left]
\begin{itemize}[nosep,leftmargin=*,itemsep=1pt,font=\scriptsize]
\item das Erz
\item das Gold
\item die Kohle
\item das Kupfer
\end{itemize}
\end{tcolorbox}
\begin{tcolorbox}[enhanced,colback=white,colframe=popgreen,boxrule=0.9pt,arc=3pt,title={die Orte},fonttitle=\bfseries\scriptsize,coltitle=white,colbacktitle=popgreen,left=4pt,right=4pt,top=2pt,bottom=2pt,boxsep=1pt,toptitle=1.5pt,bottomtitle=1.5pt,halign=left]
\begin{itemize}[nosep,leftmargin=*,itemsep=1pt,font=\scriptsize]
\item die Mine
\item der Schacht
\item der Stollen
\item der Tagebau
\item die Hütte
\end{itemize}
\end{tcolorbox}
\begin{tcolorbox}[enhanced,colback=white,colframe=poppurple,boxrule=0.9pt,arc=3pt,title={die Menschen},fonttitle=\bfseries\scriptsize,coltitle=white,colbacktitle=poppurple,left=4pt,right=4pt,top=2pt,bottom=2pt,boxsep=1pt,toptitle=1.5pt,bottomtitle=1.5pt,halign=left]
\begin{itemize}[nosep,leftmargin=*,itemsep=1pt,font=\scriptsize]
\item der Arbeiter
\item der Bergmann
\item der Ingenieur
\item die Gewerkschaft
\end{itemize}
\end{tcolorbox}
\begin{tcolorbox}[enhanced,colback=white,colframe=poppink,boxrule=0.9pt,arc=3pt,title={die Verben},fonttitle=\bfseries\scriptsize,coltitle=white,colbacktitle=poppink,left=4pt,right=4pt,top=2pt,bottom=2pt,boxsep=1pt,toptitle=1.5pt,bottomtitle=1.5pt,halign=left]
\begin{itemize}[nosep,leftmargin=*,itemsep=1pt,font=\scriptsize]
\item fördern
\item abbauen
\item schmelzen
\item verarbeiten
\end{itemize}
\end{tcolorbox}
\begin{tcolorbox}[enhanced,colback=white,colframe=popgold,boxrule=0.9pt,arc=3pt,title={die Umwelt},fonttitle=\bfseries\scriptsize,coltitle=white,colbacktitle=popgold,left=4pt,right=4pt,top=2pt,bottom=2pt,boxsep=1pt,toptitle=1.5pt,bottomtitle=1.5pt,halign=left]
\begin{itemize}[nosep,leftmargin=*,itemsep=1pt,font=\scriptsize]
\item das Recycling
\item die Nachhaltigkeit
\item die Verschmutzung
\end{itemize}
\end{tcolorbox}
\end{tcbraster}

\end{center}
\legende{Carte mentale des cinq champs lexicaux du thème « l'extraction des minerais ».}
\end{popfigure}

\begin{exoresolu}[Mise en relation : mot — article — pluriel — définition]
Énoncé : Relie chaque mot à son article, son pluriel et sa définition. Mots : \all{Mine}, \all{Schacht}, \all{Arbeiter}, \all{Erz}. Définitions : a) « Person, die in der Mine arbeitet » ; b) « Ort, wo man Rohstoffe fördert » ; c) « vertikaler Weg in die Mine » ; d) « Rohstoff, aus dem man Metall gewinnt ».
\tcblower
Solution détaillée :
\begin{itemize}
\item \all{\textbf{die} Mine, die Minen} → b) : la mine est le lieu où l'on extrait les matières premières.
\item \all{\textbf{der} Schacht, die Schächte} → c) : le puits est le chemin vertical vers la mine (attention au faux ami : ce n'est pas « le chat » !).
\item \all{\textbf{der} Arbeiter, die Arbeiter} → a) : l'ouvrier est la personne qui travaille dans la mine ; pluriel identique au singulier.
\item \all{\textbf{das} Erz, die Erze} → d) : le minerai est la matière première dont on tire le métal.
\end{itemize}
Méthode : commence toujours par vérifier le genre (les noms en \all{-er} désignant des personnes sont masculins ; les noms de matières sont souvent neutres), puis cherche la définition qui contient un mot de la même famille.
\end{exoresolu}

\begin{exercice}[À ton tour : fiches de vocabulaire]
\begin{enumerate}
\item Complète avec l'article et le pluriel : \all{\ldots{} Hütte (\ldots{})}, \all{\ldots{} Gewerkschaft (\ldots{})}, \all{\ldots{} Stollen (\ldots{})}, \all{\ldots{} Rohstoff (\ldots{})}.
\item Traduis en allemand : « Le mineur descend dans le puits. » ; « Le charbon était transporté à l'usine. » ; « Le recyclage protège l'environnement. »
\item Trouve le mot qui ne convient pas : \all{das Gold — die Kohle — das Kupfer — der Ingenieur}. Justifie ta réponse en allemand.
\item Forme le féminin et le pluriel : \all{der Ingenieur}, \all{der Arbeiter}, \all{der Experte}.
\end{enumerate}
\end{exercice}

\begin{corrige}[Exercice : fiches de vocabulaire]
\begin{enumerate}
\item \all{die Hütte, die Hütten} ; \all{die Gewerkschaft, die Gewerkschaften} ; \all{der Stollen, die Stollen} ; \all{der Rohstoff, die Rohstoffe}.
\item \all{Der Bergmann fährt in den Schacht.} (Accusatif directionnel) ; \all{Die Kohle wurde in die Hütte transportiert.} (Passif Präteritum) ; \all{Das Recycling schützt die Umwelt.} (Actif Présent).
\item \all{Der Ingenieur} : c'est une personne (Mensch), les trois autres sont des matières premières (Rohstoffe).
\item \all{die Ingenieurin, die Ingenieurinnen} ; \all{die Arbeiterin, die Arbeiterinnen} ; \all{die Expertin, die Expertinnen}.
\end{enumerate}
\end{corrige}

\begin{savaistu}[L'Allemagne et ses matières premières]
L'Allemagne possède encore du charbon (lignite) et du sel, mais elle importe aujourd'hui la quasi-totalité — on dit souvent \all{100 \%} — de son minerai de fer, principalement du Brésil, du Canada et d'Australie. La dernière mine de charbon allemand (\all{Steinkohle}) a fermé en 2018 dans la Ruhr (\all{Zeche Prosper-Haniel}). C'est pourquoi le recyclage (\all{das Recycling}, \all{die Kreislaufwirtschaft}) et la sécurisation des matières premières africaines (cobalt, lithium, terres rares) sont devenus des sujets centraux des discussions économiques entre l'Allemagne et des pays comme le Cameroun.
\end{savaistu}

\begin{aretenir}[L'essentiel du lexique]
\begin{itemize}
\item Cinq champs : les matières (\all{das Erz, das Gold, die Kohle, das Salz, das Kupfer}), les lieux (\all{die Mine, der Stollen, der Schacht, der Tagebau, die Hütte}), les acteurs (\all{der Arbeiter, der Bergmann, der Ingenieur, die Gewerkschaft}), les verbes (\all{fördern, abbauen, schmelzen, transportieren, verarbeiten, gewinnen}) et les impacts (\all{die Umweltbelastung, die Landschaftszerstörung, das Recycling, die Nachhaltigkeit}).
\item Chaque mot s'apprend avec son article, son pluriel et une phrase d'exemple.
\item Pièges : \all{der Schacht} n'est pas « le chat » ; \all{abbauen} est séparable, \all{verarbeiten} est inséparable ; les noms de matières comme \all{das Gold} n'ont pas de pluriel ; \all{der Ingenieur} fait exception à la règle du pluriel en \all{-er} invariable.
\end{itemize}
\end{aretenir}

\coursec{Grammaire : Le passif au prétérit}

\courssub{Du lexique à la grammaire : pourquoi le passif ?}

Dans la section précédente, nous avons construit le lexique du monde minier : \all{das Erz} (le minerai), \all{der Bergbau} (l'exploitation minière), \all{die Mine} (la mine), \all{das Gold} (l'or), \all{der Rohstoff} (la matière première), \all{der Arbeiter} (l'ouvrier). Or, quand on lit un article de presse ou un texte historique sur l'extraction des minerais, on s'aperçoit vite que l'on s'intéresse davantage au \cle{minerai extrait}, au \cle{charbon transporté}, à la \cle{mine fermée} qu'aux personnes qui font l'action. Pour cela, l'allemand — comme le français — utilise le \cle{passif}. Mais la presse et l'histoire racontent des faits passés : il nous faut donc le \textbf{passif au prétérit}, la forme reine du récit journalistique et historique en allemand.

\courssub{Observation : repérer la structure dans le texte}

Reprenons quelques phrases du texte support et observons ce qu'elles ont en commun.

\begin{exemplebox}[Phrases relevées dans le texte]
\begin{itemize}
\item \all{Im 19. Jahrhundert wurde in dieser Region viel Kohle gefördert.} — Au XIX\textsuperscript{e} siècle, on a extrait beaucoup de charbon dans cette région.
\item \all{Das Erz wurde mit schweren Maschinen aus der Mine geholt.} — Le minerai a été sorti de la mine avec des machines lourdes.
\item \all{Viele Arbeiter wurden von der Gesellschaft eingestellt.} — De nombreux ouvriers ont été embauchés par la société.
\item \all{Die gefährlichen Stollen wurden nach dem Unfall geschlossen.} — Les galeries dangereuses ont été fermées après l'accident.
\item \all{Das Gold wurde in die Hauptstadt transportiert.} — L'or a été transporté dans la capitale.
\end{itemize}
\end{exemplebox}

Observons : chaque phrase contient deux éléments verbaux. Le premier est une forme de \all{werden} au prétérit (\all{wurde} au singulier, \all{wurden} au pluriel) ; le second est un \cle{participe passé} (Partizip II) placé \textbf{à la fin de la phrase} : \all{gefördert}, \all{geholt}, \all{eingestellt}, \all{geschlossen}, \all{transportiert}. Le sujet grammatical (\all{Kohle}, \all{Erz}, \all{Arbeiter}, \all{Stollen}, \all{Gold}) ne fait pas l'action : il la \emph{subit}. C'est le \cle{sujet patient}.

\begin{propriete}[Règle de formation du passif au prétérit]
Le passif au prétérit se forme avec :
\begin{center}
\textbf{werden au Präteritum (wurde/wurden) + Partizip II du verbe principal (en fin de phrase)}
\end{center}
\begin{itemize}
\item \all{werden} est l'\cle{auxiliaire} : il porte la marque du temps (prétérit), de la personne et du nombre ; il se place à la \textbf{deuxième position} dans la phrase déclarative.
\item Le \cle{Partizip II} est \textbf{invariable} : il ne s'accorde jamais, contrairement au français (« la mine a été fermé\emph{e} » mais \all{die Mine wurde geschlossen}).
\item L'agent (celui qui fait l'action) est introduit par \all{von + Datif} : \all{Das Erz wurde \textbf{von den Arbeitern} gefördert.} (Le minerai a été extrait par les ouvriers.)
\item Si l'agent est un moyen ou un instrument, on utilise \all{durch + Akkusativ} ou \all{mit + Datif} : \all{Das Erz wurde \textbf{mit Maschinen} geholt.}
\end{itemize}
\end{propriete}

\courssub{Conjugaison complète : werden au prétérit + Partizip II}

Le verbe \all{werden} est irrégulier au prétérit : sa racine devient \all{wurd-}. Voici le paradigme complet avec le verbe \all{fördern} (extraire), verbe régulier dont le Partizip II est \all{gefördert}.

\begin{center}
\begin{tabular}{|l|l|l|}
\hline
\rowcolor{popblueL} \textbf{Personne} & \textbf{Passif prétérit} & \textbf{Français (écrit / oral)} \\
\hline
ich & wurde gefördert & je fus extrait / j'ai été extrait \\
\hline
du & wurdest gefördert & tu fus extrait / tu as été extrait \\
\hline
er / sie / es & wurde gefördert & il / elle fut extrait(e) / il / elle a été extrait(e) \\
\hline
wir & wurden gefördert & nous fûmes extraits / nous avons été extraits \\
\hline
ihr & wurdet gefördert & vous fûtes extraits / vous avez été extraits \\
\hline
sie / Sie & wurden gefördert & ils / elles furent extraits / ils / elles ont été extraits \\
\hline
\end{tabular}
\end{center}

\begin{attention}[wurde ou wurden ?]
Le choix entre \all{wurde} et \all{wurden} dépend uniquement du \textbf{sujet patient}, pas de l'agent. \all{Die Kohle wurde gefördert} (singulier) mais \all{Die Arbeiter wurden eingestellt} (pluriel). Attention aussi aux formes : \all{du wurdest}, \all{ihr wurdet} — n'oubliez pas le \all{-t} !
\end{attention}

\courssub{Le Partizip II : réguliers, irréguliers et verbes à particule}

Le Partizip II suit trois modèles que vous connaissez déjà du Perfekt :

\begin{center}
\begin{tabularx}{\linewidth}{|X|X|X|}
\hline
\rowcolor{popblueL} \textbf{Type de verbe} & \textbf{Formation} & \textbf{Exemple au passif prétérit} \\
\hline
Verbe faible (régulier) & ge + radical + t & \all{fördern → Das Erz wurde gefördert.} \\
\hline
Verbe fort (irrégulier) & ge + (racine changée) + en & \all{finden → Gold wurde gefunden.} \\
\hline
Verbe en -ieren & pas de ge- & \all{transportieren → Das Gold wurde transportiert.} \\
\hline
\end{tabularx}
\end{center}

\begin{propriete}[Verbes à particule séparable au passif prétérit]
Pour les verbes à \cle{particule séparable}, le \all{-ge-} du Partizip II s'insère \textbf{entre la particule et le radical} : \all{abbauen → abgebaut}, \all{ausbeuten → ausgebeutet}, \all{freilegen → freigelegt}. Au passif prétérit, la particule reste donc soudée au participe en fin de phrase :
\begin{itemize}
\item \all{Die Kohle wurde in dieser Mine \textbf{abgebaut}.} — Le charbon a été exploité dans cette mine.
\item \all{Die Rohstoffe wurden rücksichtslos \textbf{ausgebeutet}.} — Les matières premières ont été exploitées sans scrupule.
\end{itemize}
En revanche, les verbes à particule \textbf{inséparable} (\all{ver-}, \all{be-}, \all{er-}…) ne prennent pas de \all{ge-} : \all{verkaufen → Das Erz wurde verkauft} ; \all{zerstören → Die Landschaft wurde zerstört}.
\end{propriete}

\courssub{La transformation actif → passif}

Passer de la voix active à la voix passive est une opération mécanique en trois étapes. Observons d'abord la transformation :

\begin{center}
\begin{tabularx}{\linewidth}{|X|X|}
\hline
\rowcolor{popblueL} \textbf{Actif (Präteritum)} & \textbf{Passif (Präteritum)} \\
\hline
\all{Die Arbeiter förderten das Erz.} & \all{Das Erz wurde von den Arbeitern gefördert.} \\
\hline
\all{Die Firma baute die Mine ab.} & \all{Die Mine wurde von der Firma abgebaut.} \\
\hline
\all{Man transportierte das Gold.} & \all{Das Gold wurde transportiert.} \\
\hline
\end{tabularx}
\end{center}

\begin{methode}[Transformer une phrase active en phrase passive au prétérit]
\begin{enumerate}
\item \textbf{Identifier le sujet patient} : repérer le \emph{complément d'objet direct} (accusatif) de la phrase active. Il devient le sujet (nominatif) de la phrase passive. \all{das Erz} → \all{Das Erz wurde …}
\item \textbf{Mettre l'auxiliaire au prétérit} : conjuguer \all{werden} au prétérit en l'accordant avec le nouveau sujet : \all{wurde} (singulier) ou \all{wurden} (pluriel).
\item \textbf{Ajouter le Partizip II} du verbe principal \emph{à la fin} de la phrase : \all{gefördert}.
\item \textbf{Exprimer l'agent} (facultatif) avec \all{von + Datif} : \all{von den Arbeitern}. Si la phrase active a pour sujet \all{man} (on), l'agent disparaît simplement : \all{Man förderte Kohle. → Kohle wurde gefördert.}
\end{enumerate}
\end{methode}

\begin{attention}[Le piège n° 1 du francophone : « ist » au lieu de « wurde »]
En français, « le charbon \emph{a été} extrait » contient « est », ce qui pousse beaucoup d'élèves à écrire \all{*Die Kohle ist gefördert worden} quand ils racontent au passé. Retenez bien :
\begin{itemize}
\item \all{Die Kohle \textbf{wird} gefördert.} = le charbon \emph{est} extrait (présent).
\item \all{Die Kohle \textbf{wurde} gefördert.} = le charbon \emph{a été / fut} extrait (prétérit, écrit).
\item \all{Die Kohle \textbf{ist} gefördert \textbf{worden}.} = le charbon \emph{a été} extrait (Perfekt, oral).
\end{itemize}
Au prétérit passif, l'auxiliaire est \all{wurde/wurden}, jamais \all{ist}. Et notez la forme spéciale du participe de \all{werden} au Perfekt passif : \all{worden} (et non \all{geworden}) !
\end{attention}

\courssub{La frise des temps du passif}

La figure suivante situe le passif prétérit parmi les quatre temps du passif que vous rencontrerez au lycée.

\begin{popfigure}
\begin{center}
\begin{tikzpicture}[scale=1.0]
\draw[-{Stealth[length=3mm]}, thick, popdark] (0,0) -- (12.6,0);
\foreach \x/\t/\e/\c in {
  1.5/{Présent}/{wird gefördert}/popblue,
  4.8/{Perfekt}/{ist gefördert worden}/popgreen,
  8.1/{Prétérit}/{wurde gefördert}/poporange,
  11.4/{Plus-que-parfait}/{war gefördert worden}/poppurple}
{
  \fill[\c] (\x,0) circle (0.12);
  \node[above=0.15cm, align=center, font=\bfseries, text=\c] at (\x,0) {\t};
  \node[below=0.18cm, align=center, font=\small] at (\x,0) {\all{\e}};
}
\draw[thick, poporange, -{Stealth[length=2.5mm]}] (8.1,-1.1) -- (8.1,-1.7);
\node[below, align=center, text=poporange, font=\small\itshape] at (8.1,-1.7) {temps du récit écrit\\(presse, histoire)};
\end{tikzpicture}
\end{center}
\legende{Frise des temps du passif allemand : auxiliaire werden conjugué + Partizip II invariable.}
\end{popfigure}

\begin{savaistu}[Prétérit ou Perfekt : une question de canal]
En allemand, le choix entre passif prétérit et passif Perfekt dépend surtout du \textbf{registre}. À l'\emph{écrit} — articles de journaux, livres d'histoire, rapports officiels — c'est le prétérit qui domine : \all{Die Mine wurde 1960 geschlossen.} (La mine fut fermée en 1960.) À l'\emph{oral}, dans la conversation quotidienne, les Allemands préfèrent le Perfekt : \all{Die Mine ist 1960 geschlossen worden.} C'est exactement la même répartition que pour l'actif. Un bon article de presse sur l'extraction du cobalt ou de la bauxite sera donc presque entièrement écrit au passif prétérit.
\end{savaistu}

\courssub{Entraînement : un court texte au passif prétérit}

\begin{textebox}[Aus der Geschichte des Bergbaus]
\all{Schon vor hundert Jahren wurde in vielen Regionen Deutschlands Kohle gefördert. Zuerst wurde das Erz von Hand aus der Erde geholt, später wurden große Maschinen benutzt. Tausende Arbeiter wurden in den Minen beschäftigt, und ganze Städte wurden um die Zechen gebaut. Die Arbeit wurde von den Bergleuten oft als gefährlich beschrieben: Unfälle wurden gemeldet, kranke Arbeiter wurden ins Krankenhaus gebracht. Nach dem Zweiten Weltkrieg wurde der Bergbau modernisiert, aber viele Stollen wurden in den 1960er Jahren geschlossen, weil importierte Kohle billiger wurde. Heute werden alte Minen oft als Museen genutzt: Dort wird den Besuchern gezeigt, wie früher gearbeitet wurde.}\par
\emph{Glossaire :} die Zeche, -n (la mine de charbon) ; der Bergmann, die Bergleute (le mineur) ; beschäftigen (employer) ; der Stollen, - (la galerie) ; melden (signaler) ; importieren (importer) ; nutzen (utiliser) ; der Besucher, - (le visiteur).
\end{textebox}

Comptez les passifs prétérit dans ce texte : vous en trouverez plus de dix ! C'est bien la preuve que ce temps est la colonne vertébrale du récit historique en allemand.

\courssub{Exercices résolus et entraînement}

\begin{exoresolu}[Transformer à la voix passive]
Énoncé : mettez la phrase active suivante au passif prétérit : \all{Die Arbeiter förderten das Erz aus der Mine.}
\tcblower
Solution détaillée, étape par étape :
\begin{enumerate}
\item Sujet patient : le COD est \all{das Erz} (accusatif) → il devient sujet au nominatif : \all{Das Erz}.
\item Auxiliaire : \all{das Erz} est singulier → \all{wurde}, placé en deuxième position.
\item Partizip II de \all{fördern} (verbe faible) : \all{gefördert}, placé à la fin.
\item Agent : \all{die Arbeiter} → \all{von den Arbeitern} (datif pluriel).
\end{enumerate}
Résultat : \all{\textbf{Das Erz wurde von den Arbeitern aus der Mine gefördert.}} — Le minerai a été extrait de la mine par les ouvriers.
\end{exoresolu}

\begin{exoresolu}[Conjuguer au passif prétérit]
Énoncé : complétez avec la forme correcte : \all{Viele Minen … (schließen, Passiv Präteritum).}
\tcblower
Solution : le sujet \all{viele Minen} est pluriel → auxiliaire \all{wurden} ; \all{schließen} est un verbe fort → Partizip II \all{geschlossen}. Phrase complète : \all{\textbf{Viele Minen wurden geschlossen.}} — De nombreuses mines ont été fermées. (Attention : participe invariable, pas d'accord !)
\end{exoresolu}

\begin{exercice}[À vous : cinq transformations]
Mettez les phrases actives suivantes au passif prétérit :
\begin{enumerate}
\item \all{Die Firma baute die Kohle in dieser Region ab.}
\item \all{Man transportierte das Gold in die Hauptstadt.}
\item \all{Die Ingenieure modernisierten die Maschinen.}
\item \all{Die Gesellschaft stellte viele Arbeiter ein.}
\item \all{Die Bergleute entdeckten ein neues Erz.}
\end{enumerate}
\end{exercice}

\begin{corrige}[Exercice « À vous : cinq transformations »]
\begin{enumerate}
\item \all{Die Kohle wurde von der Firma in dieser Region abgebaut.} — Le charbon a été exploité par la firme dans cette région. (Verbe à particule : \all{abgebaut} !)
\item \all{Das Gold wurde in die Hauptstadt transportiert.} — L'or a été transporté dans la capitale. (\all{man} disparaît ; \all{transportieren} : pas de \all{ge-}.)
\item \all{Die Maschinen wurden von den Ingenieuren modernisiert.} — Les machines ont été modernisées par les ingénieurs. (Pluriel → \all{wurden}.)
\item \all{Viele Arbeiter wurden von der Gesellschaft eingestellt.} — De nombreux ouvriers ont été embauchés par la société. (Particule séparable : \all{eingestellt}.)
\item \all{Ein neues Erz wurde von den Bergleuten entdeckt.} — Un nouveau minerai a été découvert par les mineurs.
\end{enumerate}
\end{corrige}

\begin{aretenir}[L'essentiel du passif au prétérit]
\begin{itemize}
\item Forme : \cle{wurde/wurden + Partizip II} (participe à la fin, invariable).
\item \all{werden} au prétérit : \all{ich wurde, du wurdest, er wurde, wir wurden, ihr wurdet, sie wurden}.
\item Agent : \all{von + Datif} ; instrument : \all{mit + Datif} / \all{durch + Akkusativ}.
\item Verbes à particule séparable : \all{abbauen → wurde abgebaut} (ge- inséré).
\item Verbes à particule inséparable / en -ieren : pas de \all{ge-} (\all{verkauft, transportiert}).
\item Jamais \all{ist} au prétérit : \all{ist … worden} est le Perfekt passif, réservé à l'oral.
\item Emploi : récit écrit, presse, histoire — le temps idéal pour parler de l'extraction des minerais.
\end{itemize}
\end{aretenir}

Maintenant que vous savez raconter au passé ce qui \emph{a été extrait, transporté, fermé}, nous allons dans la section suivante enrichir nos phrases grâce aux \textbf{propositions relatives introduites par une préposition} (\all{die Mine, in der das Erz gefördert wurde…}), un outil précieux pour décrire précisément le monde minier.

\coursec{Grammaire : Relatives avec préposition}

Dans la section précédente, nous avons appris à raconter l'extraction des minerais au passif prétérit : \all{Das Erz wurde abgebaut, die Arbeiter wurden bezahlt.} (Le minerai a été extrait, les ouvriers ont été payés.) Mais un texte journalistique ou scientifique ne se contente pas de phrases courtes : il enchaîne les informations grâce aux \cle{propositions relatives}. Or, très souvent, le nom repris par le relatif est introduit par une \cle{préposition} : \all{die Mine, in der gearbeitet wurde} (la mine dans laquelle on travaillait). C'est cette construction, typique de l'allemand écrit soigné et redoutée des francophones, que nous allons maîtriser pas à pas.

\courssub{Observation : des exemples extraits du texte support}

Commençons par observer quatre phrases tirées de notre texte sur l'extraction des minerais. Soulignez mentalement la préposition et le pronom relatif.

\begin{exemplebox}[Phrases à observer]
\begin{itemize}
\item \all{Die Mine, in der gearbeitet wurde, ist geschlossen.} « La mine dans laquelle on travaillait est fermée. »
\item \all{Der Bergbau, von dem die Region lebt, ist riskant.} « L'exploitation minière, dont vit la région, est risquée. »
\item \all{Der Arbeiter, mit dem ich gesprochen habe, kennt die Mine gut.} « L'ouvrier avec qui j'ai parlé connaît bien la mine. »
\item \all{Das Gold, für das so viele Menschen starben, liegt tief unter der Erde.} « L'or pour lequel tant d'hommes sont morts repose profondément sous la terre. »
\end{itemize}
\end{exemplebox}

Que remarque-t-on ?

\begin{enumerate}
\item La proposition relative est toujours séparée par une \cle{virgule}.
\item Le pronom relatif (\all{der, dem, das}) est précédé d'une préposition (\all{in, von, mit, für}).
\item La préposition impose son \cle{cas} au pronom relatif : \all{in der} (datif), \all{von dem} (datif), \all{mit dem} (datif), \all{für das} (accusatif).
\item Le verbe conjugué de la relative se place, comme toujours, à la \cle{fin} de la proposition.
\end{enumerate}

\courssub{Rappel : les pronoms relatifs simples et leur accord}

Avant d'ajouter une préposition, rappelons le tableau des pronoms relatifs simples. Le pronom relatif s'accorde en \cle{genre} et en \cle{nombre} avec le nom auquel il se réfère (l'antécédent), mais son \cle{cas} dépend de sa fonction dans la proposition relative.

\begin{center}
\begin{tabular}{|l|l|l|l|}
\hline
\rowcolor{popblueL} Cas & Masculin & Féminin & Neutre / Pluriel \\
\hline
Nominatif & der & die & das / die \\
\hline
Accusatif & den & die & das / die \\
\hline
Datif & dem & der & dem / denen \\
\hline
Génitif & dessen & deren & dessen / deren \\
\hline
\end{tabular}
\end{center}

\begin{exemplebox}[Relatives simples sans préposition]
\begin{itemize}
\item \all{Das Erz, das man hier findet, ist sehr wertvoll.} « Le minerai qu'on trouve ici est très précieux. » (accusatif neutre)
\item \all{Die Arbeiter, die in der Mine arbeiten, sind mutig.} « Les ouvriers qui travaillent dans la mine sont courageux. » (nominatif pluriel)
\item \all{Der Ingenieur, dem wir das Gold zeigten, war überrascht.} « L'ingénieur à qui nous avons montré l'or était surpris. » (datif masculin)
\end{itemize}
\end{exemplebox}

\courssub{La construction : préposition + pronom relatif}

\begin{propriete}[Règle de formation]
Quand le nom repris par le relatif est le complément d'une \cle{préposition} (exigée par le verbe, l'adjectif ou le sens), la construction allemande est :

\textbf{antécédent + virgule + PRÉPOSITION + pronom relatif + ... + verbe conjugué en fin de proposition.}

Le pronom relatif prend le \cle{cas exigé par la préposition} (et non par le verbe) :
\begin{itemize}
\item \all{in} + datif (lieu) $\rightarrow$ \all{in dem / in der} ; \all{in} + accusatif (mouvement) $\rightarrow$ \all{in den / in das}
\item \all{mit, von, nach, zu, bei, aus} + datif $\rightarrow$ \all{mit dem, von der, nach dem...}
\item \all{für, gegen, ohne, durch, über} (sens abstrait) + accusatif $\rightarrow$ \all{für den, über das...}
\end{itemize}
\end{propriete}

\begin{propriete}[Tableau des prépositions fréquentes et de leur cas]
\begin{center}
\begin{tabularx}{\linewidth}{|X|X|X|}
\hline
\rowcolor{popblueL} Préposition & Cas & Exemple de relative \\
\hline
in (lieu) & Dativ & die Mine, in der... \\
\hline
an, auf (lieu) & Dativ & der Markt, auf dem... \\
\hline
in, auf (mouvement) & Akkusativ & das Land, in das... \\
\hline
mit, nach, von, zu, bei, aus & Dativ & der Arbeiter, mit dem... \\
\hline
für, gegen, ohne, durch, um & Akkusativ & das Gold, für das... \\
\hline
über (au sujet de) & Akkusativ & der Bericht, über den... \\
\hline
\end{tabularx}
\end{center}
\end{propriete}

\begin{popfigure}
\begin{tikzpicture}[node distance=0.6cm and 1.1cm,
box/.style={draw=popblue, fill=popblueL, rounded corners=2pt, align=center, minimum height=0.9cm, inner sep=4pt},
arr/.style={-{Stealth[length=3mm]}, thick, poporange}]
\node[box, align=center] (prep){\textbf{Préposition}\\ \all{in} (lieu)};
\node[box, right=of prep, align=center] (kasus){\textbf{Cas exigé}\\ Dativ};
\node[box, right=of kasus, align=center] (rel){\textbf{Pronom relatif}\\ \all{dem / der / dem}};
\node[box, right=of rel, align=center] (res){\textbf{Résultat}\\ \all{in dem, in der}};
\draw[arr] (prep) -- (kasus);
\draw[arr] (kasus) -- (rel);
\draw[arr] (rel) -- (res);
\node[box, below=0.9cm of prep, fill=popgreenL, draw=popgreen, align=center] (prep2){\textbf{Préposition}\\ \all{für}};
\node[box, right=of prep2, fill=popgreenL, draw=popgreen, align=center] (kasus2){\textbf{Cas exigé}\\ Akkusativ};
\node[box, right=of kasus2, fill=popgreenL, draw=popgreen, align=center] (rel2){\textbf{Pronom relatif}\\ \all{den / das}};
\node[box, right=of rel2, fill=popgreenL, draw=popgreen, align=center] (res2){\textbf{Résultat}\\ \all{für den, für das}};
\draw[arr] (prep2) -- (kasus2);
\draw[arr] (kasus2) -- (rel2);
\draw[arr] (rel2) -- (res2);
\end{tikzpicture}
\legende{De la préposition au pronom relatif : la préposition choisit le cas, l'antécédent choisit le genre.}
\end{popfigure}

\courssub{Comparaison avec le français}

Le français utilise souvent \emph{dont}, \emph{où}, \emph{auquel}, \emph{duquel}, \emph{pour lequel} ; l'allemand, lui, reste systématique : préposition + article relatif décliné.

\begin{center}
\begin{tabularx}{\linewidth}{|X|X|X|}
\hline
\rowcolor{popblueL} Français & Deutsch & Remarque \\
\hline
la mine dans laquelle & die Mine, in der & in + Dativ, féminin \\
\hline
l'ouvrier avec qui & der Arbeiter, mit dem & mit + Dativ, masculin \\
\hline
l'or pour lequel & das Gold, für das & für + Akkusativ, neutre \\
\hline
la région dont (de laquelle) & die Region, von der & von + Dativ, féminin \\
\hline
les ouvriers à qui / avec lesquels & die Arbeiter, mit denen & datif pluriel : denen \\
\hline
\end{tabularx}
\end{center}

\begin{attention}[Piège n° 1 : la virgule]
En allemand, la proposition relative est \textbf{toujours} précédée d'une virgule, même quand le français n'en met pas. On écrit : \all{Die Mine, in der gearbeitet wurde, ist geschlossen.} et jamais \all{Die Mine in der gearbeitet wurde...}. Oublier la virgule est la faute la plus fréquente des francophones à l'écrit.
\end{attention}

\begin{attention}[Piège n° 2 : le cas vient de la préposition, pas du français]
Ne traduisez pas mécaniquement « dont » par \all{von dem} : vérifiez la préposition exigée par le verbe allemand. « Le métal dont on a besoin » se dit \all{das Metall, das man braucht} (car \all{brauchen} se construit sans préposition), mais « l'aide dont on a besoin » se dit \all{die Hilfe, deren man bedarf} ou, plus simplement, \all{die Hilfe, die man braucht}. Interrogez toujours la construction du verbe allemand : \all{sprechen über + Akk.}, \all{träumen von + Dat.}, \all{warten auf + Akk.}, \all{leben von + Dat.}.
\end{attention}

\begin{savaistu}[À l'écrit et à l'oral]
En allemand standard et formel (presse, dissertation, administration), la préposition se place obligatoirement \textbf{devant} le pronom relatif : \all{der Stuhl, auf dem ich sitze}. À l'oral familier, on entend parfois la préposition reportée en fin de phrase avec \all{da(r)-} ou une reprise : \all{der Stuhl, wo ich drauf sitze} — tournure régionale et familière à éviter dans vos copies. 

Retenez aussi les formes contractées : \all{in dem $\rightarrow$ im}, \all{an dem $\rightarrow$ am}, \all{von dem $\rightarrow$ vom}. Elles sont correctes mais, dans les relatives formelles, on préfère souvent les formes pleines (\all{in dem}, \all{von dem}) pour lever toute ambiguïté avec le pronom démonstratif. Exemple : \all{der Ort, in dem (oder: im) ich wohne}.
\end{savaistu}

\courssub{Méthode : construire une relative avec préposition en quatre étapes}

\begin{methode}[Quatre étapes pour ne plus se tromper]
\begin{enumerate}
\item \textbf{Identifier l'antécédent} et noter son genre et son nombre : \all{die Mine} $\rightarrow$ féminin singulier.
\item \textbf{Identifier le verbe de la relative} et la préposition qu'il exige : \all{arbeiten in + Dativ} (lieu).
\item \textbf{Déterminer le cas} imposé par la préposition : \all{in} + lieu $\rightarrow$ datif.
\item \textbf{Choisir le pronom relatif} (genre de l'antécédent + cas de la préposition) : féminin datif $\rightarrow$ \all{der}. Résultat : \all{die Mine, in der...}. Ne pas oublier la virgule et le verbe conjugué en fin de proposition.
\end{enumerate}
\end{methode}

\begin{exoresolu}[Exercice de complétion résolu]
Complétez avec la préposition et le pronom relatif qui conviennent : \all{Der Rohstoff, \_\_\_\_ \_\_\_\_ die Wirtschaft abhängt, wird teurer.} (Le matériau brut dont dépend l'économie devient plus cher.)
\tcblower
\textbf{Solution détaillée.} Étape 1 : l'antécédent est \all{der Rohstoff}, masculin singulier. Étape 2 : le verbe est \all{abhängen}, qui se construit avec \all{von + Dativ} (dépendre de). Étape 3 : \all{von} exige le datif. Étape 4 : masculin datif $\rightarrow$ \all{dem}. Phrase correcte : \all{Der Rohstoff, von dem die Wirtschaft abhängt, wird teurer.} Pièges écartés : \all{von den} (accusatif, impossible après \all{von}) et \all{von der} (féminin, alors que \all{Rohstoff} est masculin).
\end{exoresolu}

\courssub{Texte d'application : lire et repérer les relatives}

\begin{textebox}[„Gold und Cobalt: Reichtum unter der Erde“]
In vielen Ländern Afrikas gibt es Bodenschätze, von denen die ganze Welt träumt. Das Gold, für das früher Tausende von Arbeitern in tiefen Minen schufteten, wird heute oft mit modernen Maschinen gefördert. Doch die Mine, in der die Menschen arbeiten, bleibt ein gefährlicher Ort. Der Arbeiter, mit dem unser Reporter gesprochen hat, erzählt: „Der Bergbau, von dem meine Familie lebt, hat mir alles gegeben, aber er hat auch meinen Vater krank gemacht.“ Die Region, in der das Erz abgebaut wird, verändert sich schnell: Straßen, auf denen früher nur Fußgänger gingen, fahren heute schwere Lastwagen. Es gibt Dörfer, in denen neue Schulen gebaut wurden, aber auch Flüsse, in denen das Wasser schmutzig wurde. Die Firma, für die die Menschen arbeiten, verspricht bessere Löhne und mehr Sicherheit. Viele Jugendliche denken über die Zukunft nach, über das Land, in dem ihre Kinder einmal leben werden. Experten warnen: Der Rohstoff, um den so hart gekämpft wird, ist nicht unendlich. Man muss über Methoden nachdenken, mit denen die Natur geschützt werden kann.
\end{textebox}

\textbf{Glossaire} : der Bodenschatz, die Bodenschätze (richesse du sous-sol) ; schuften (trimer) ; fördern (extraire) ; der Lastwagen, die Lastwagen (camion) ; der Lohn, die Löhne (salaire) ; unendlich (inépuisable) ; schützen (protéger).

\textbf{Questions de compréhension} : 1. \all{Wovon lebt die Familie des Arbeiters?} 2. \all{Was verspricht die Firma?} 3. \all{Wovor warnen die Experten?} 4. Relevez dans le texte quatre relatives avec préposition et indiquez pour chacune la préposition, le cas et l'antécédent.

\courssub{Entraînement et production}

\begin{exercice}[Choisissez la bonne préposition et le bon pronom relatif]
Complétez : 1. \all{Die Arbeiter, \_\_\_\_ \_\_\_\_ wir gesprochen haben, sind müde.} (mit) 2. \all{Das Dorf, \_\_\_\_ \_\_\_\_ die Mine liegt, heißt Bafia.} (in, lieu) 3. \all{Der Bericht, \_\_\_\_ \_\_\_\_ alle diskutieren, erscheint morgen.} (über) 4. \all{Die Maschinen, \_\_\_\_ \_\_\_\_ das Erz gefördert wurde, sind alt.} (mit) 5. \all{Der Freund, \_\_\_\_ \_\_\_\_ ich nach Yaoundé reise, ist Ingenieur.} (mit)
\end{exercice}

\begin{corrige}[Exercice « Choisissez la bonne préposition »]
1. \all{mit denen} (mit + Dativ, pluriel). 2. \all{in dem} (in + Dativ, neutre, lieu). 3. \all{über den} (über + Akkusativ, masculin : \all{diskutieren über + Akk.}). 4. \all{mit denen} (mit + Dativ, pluriel). 5. \all{mit dem} (mit + Dativ, masculin).
\end{corrige}

\begin{exercice}[Production écrite : trois phrases originales]
Rédigez trois phrases originales sur le thème minier (Cameroun ou pays germanophone), chacune contenant une relative avec préposition différente : une avec \all{in} + datif, une avec \all{mit} + datif, une avec \all{für} ou \all{von}. Modèle : \all{Der Markt, auf dem man Coltan verkauft, liegt im Zentrum der Stadt.}
\end{exercice}

\begin{experience}[Jeu de la phrase qui s'allonge]
Par groupes de quatre : le premier élève énonce un nom du champ lexical minier (\all{die Mine}), le deuxième ajoute une relative avec préposition (\all{die Mine, in der mein Onkel arbeitet}), le troisième transforme la phrase au passif prétérit (\all{die Mine, in der früher Gold gefördert wurde}), et le quatrième traduit en français. On change de rôle à chaque tour. L'enseignant note au tableau les prépositions utilisées et corrige les cas en direct.
\end{experience}

\begin{aretenir}[L'essentiel de la section]
\begin{itemize}
\item Structure : antécédent + \textbf{virgule} + préposition + pronom relatif + ... + verbe en fin.
\item Le \cle{genre} du pronom vient de l'antécédent ; le \cle{cas} vient de la préposition (\all{mit, von, in} (lieu) + datif ; \all{für, über, ohne} + accusatif).
\item Datif pluriel : \all{denen} ; génitif : \all{dessen / deren}.
\item Le français \emph{dont, où, auquel} se rend par préposition + relatif décliné, selon la construction du verbe allemand.
\item Jamais de relative sans virgule, jamais de préposition après le relatif à l'écrit formel.
\end{itemize}
\end{aretenir}

\coursec{Commentaire guidé \& production écrite}

Après avoir étudié dans la section précédente les propositions relatives introduites par une préposition (\all{die Mine, in der gearbeitet wird…}), nous disposons désormais de tous les outils nécessaires : le vocabulaire du thème minier, le passif au prétérit et les relatives complexes. Il est temps de les \cle{réinvestir} dans une production écrite complète : le \cle{commentaire de texte} (\all{der Kommentar}). C'est l'exercice roi de la classe de Terminale : il évalue à la fois la compréhension, la grammaire et l'expression personnelle, et il est au cœur des épreuves écrites du Baccalauréat.

\courssub{Qu'est-ce qu'un commentaire ?}

\begin{definition}[der Kommentar ou die Zusammenfassung ?]
\all{Die Zusammenfassung} (le résumé) est un texte qui reprend \textbf{uniquement} les idées du texte source, de manière objective, raccourcie, sans opinion personnelle.\par
\all{Der Kommentar} (le commentaire) va plus loin : il présente le texte, le résume, \textbf{analyse sa langue} (grammaire, vocabulaire, style) et se termine par une \textbf{opinion personnelle argumentée} (\all{die eigene Meinung}).\par
En Terminale A, le commentaire court (environ 120 mots) combine les deux : un résumé structuré \emph{plus} une analyse \emph{plus} un avis.
\end{definition}

\begin{attention}[Reformuler n'est pas copier]
L'erreur la plus grave dans un résumé ou un commentaire est de \textbf{recopier des phrases entières du texte}. Il faut reformuler avec ses propres mots : c'est ce qu'on appelle \all{mit eigenen Worten wiedergeben} (restituer avec ses propres mots). Recopier trois phrases du texte fait perdre la moitié des points de la partie « résumé ». Astuce : changez la structure de la phrase (actif/passif), remplacez les noms par des pronoms et les verbes par des synonymes (\all{entdecken} $\rightarrow$ \all{finden}, \all{zeigen} $\rightarrow$ \all{erklären}).
\end{attention}

\courssub{La méthode du commentaire en 4 étapes}

\begin{methode}[Rédiger un commentaire en 4 étapes]
\begin{enumerate}
\item \textbf{Introduction} (\all{die Einleitung}) : présenter le texte en une ou deux phrases — le titre, l'auteur ou la source, le thème. Formule type : \all{Der Text „…“ behandelt das Thema …} (Le texte « … » traite le thème …).
\item \textbf{Résumé structuré} (\all{die Zusammenfassung}) : reprendre les idées principales dans l'ordre du texte, en 3 ou 4 phrases, avec ses propres mots et au \textbf{présent} ou au \textbf{prétérit} (mais un seul temps !).
\item \textbf{Analyse linguistique} (\all{die Analyse}) : relever les procédés de langue étudiés — ici : le passif au prétérit (\all{wurde abgebaut}), les relatives avec préposition (\all{die Region, in der…}), le vocabulaire du champ lexical minier.
\item \textbf{Opinion personnelle} (\all{die eigene Meinung}) : dire ce qu'on pense du sujet, en s'appuyant sur des arguments. Formules types : \all{Meiner Meinung nach…} (à mon avis), \all{Ich finde es wichtig, dass…} (je trouve important que…).
\end{enumerate}
\end{methode}

\begin{popfigure}
\begin{tikzpicture}[node distance=0.55cm,
  etape/.style={rectangle, rounded corners=4pt, draw=popblue, fill=popblueL, thick, minimum width=7.6cm, minimum height=0.95cm, align=center, font=\small},
  fleche/.style={-{Stealth[length=2.5mm]}, thick, poporange}]
\node[etape] (e1) {\textbf{1. Einleitung} : Titel, Autor, Thema};
\node[etape, below=of e1, fill=popgreenL, draw=popgreen] (e2) {\textbf{2. Zusammenfassung} : Hauptideen, eigene Worte};
\node[etape, below=of e2, fill=poppurpleL, draw=poppurple] (e3) {\textbf{3. Analyse} : Passiv, Relativsätze, Wortschatz};
\node[etape, below=of e3, fill=popgoldL, draw=popgold] (e4) {\textbf{4. Meinung} : eigene Position + Argumente};
\draw[fleche] (e1) -- (e2);
\draw[fleche] (e2) -- (e3);
\draw[fleche] (e3) -- (e4);
\end{tikzpicture}
\legende{Organigramme du commentaire de texte : quatre étapes enchaînées}
\end{popfigure}

\courssub{Les connecteurs logiques : la charpente du commentaire}

Un commentaire sans connecteurs est une liste de phrases isolées ; le correcteur ne voit pas la progression de la pensée. Les \cle{connecteurs logiques} (\all{die Konnektoren}) relient les étapes entre elles.

\begin{methode}[Utiliser les connecteurs logiques]
\begin{enumerate}
\item Pour \textbf{commencer} une étape ou une énumération : \all{zunächst} (d'abord, tout d'abord), \all{zuerst} (d'abord).
\item Pour \textbf{ajouter} une idée : \all{außerdem} (de plus), \all{darüber hinaus} (en outre), \all{auch} (aussi).
\item Pour \textbf{opposer} deux idées : \all{jedoch} (cependant), \all{aber} (mais), \all{einerseits … andererseits} (d'une part … d'autre part).
\item Pour \textbf{conclure} : \all{abschließend} (pour conclure), \all{zusammenfassend} (en résumé), \all{schließlich} (finalement).
\end{enumerate}
Attention à la place du verbe : après \all{zunächst}, \all{außerdem}, \all{jedoch}, \all{abschließend} en tête de phrase, le verbe conjugué vient \textbf{en deuxième position} : \all{Außerdem \textbf{wurden} viele Arbeiter eingestellt.} (De plus, de nombreux ouvriers furent embauchés.)\par
Cas particulier : \all{aber} est une conjonction de coordination, elle \textbf{n'occupe pas de position} dans la phrase ; le verbe reste juste après le sujet : \all{Der Bergbau bringt Geld, \textbf{aber} er zerstört oft die Umwelt.} En revanche, \all{jedoch} est un adverbe : \all{Der Bergbau bringt Geld, \textbf{jedoch} zerstört er oft die Umwelt.} (verbe avant le sujet).
\end{methode}

\begin{exemplebox}[Connecteurs en contexte — exemples traduits]
\all{Zunächst beschreibt der Autor die Entdeckung des Goldes in der Region.} « Tout d'abord, l'auteur décrit la découverte de l'or dans la région. »\par
\all{Außerdem erklärt er, wie das Erz abgebaut wurde.} « De plus, il explique comment le minerai fut extrait. »\par
\all{Der Bergbau bringt Geld, jedoch zerstört er oft die Umwelt.} « L'exploitation minière rapporte de l'argent, cependant elle détruit souvent l'environnement. »\par
\all{Zusammenfassend zeigt der Text, dass der Bergbau sowohl Chancen als auch Risiken birgt.} « En résumé, le texte montre que l'exploitation minière comporte à la fois des chances et des risques. »\par
\all{Abschließend kann man sagen, dass der Staat die Minen besser kontrollieren muss.} « Pour conclure, on peut dire que l'État doit mieux contrôler les mines. »
\end{exemplebox}

\begin{center}
\begin{tabularx}{\linewidth}{|X|X|X|}
\hline
\rowcolor{popblueL} \textbf{Fonction} & \textbf{Deutsch} & \textbf{Français} \\
\hline
commencer & zunächst, zuerst & d'abord, tout d'abord \\
\hline
ajouter & außerdem, darüber hinaus & de plus, en outre \\
\hline
opposer & jedoch, aber, dagegen & cependant, mais, en revanche \\
\hline
nuancer & einerseits … andererseits & d'une part … d'autre part \\
\hline
conclure & abschließend, zusammenfassend & pour conclure, en résumé \\
\hline
donner son avis & meiner Meinung nach, ich finde, dass & à mon avis, je trouve que \\
\hline
\end{tabularx}
\end{center}

\begin{attention}[Ne mélangez pas les temps dans le résumé]
Erreur très fréquente des francophones : commencer le résumé au prétérit (\all{Der Autor beschrieb…}) puis passer au parfait (\all{dann hat er erklärt…}) ou au présent sans raison. \textbf{Règle} : choisissez un temps de référence et tenez-vous-y. Pour un texte narratif ou journalistique, le \textbf{prétérit} est le temps du résumé écrit : \all{Der Autor beschrieb die Mine. Dann erklärte er die Folgen.} Le parfait (\all{hat erklärt}) appartient à l'oral ; le mélange \all{beschrieb / hat erklärt} dans le même paragraphe est fautif. De même, évitez le présent de narration si vous avez commencé au passé.\par
Rappel de conjugaison au prétérit : \all{behandeln} $\rightarrow$ \all{behandelte} ; \all{beschreiben} $\rightarrow$ \all{beschrieb} (verbe fort) ; \all{erklären} $\rightarrow$ \all{erklärte} ; \all{zeigen} $\rightarrow$ \all{zeigte}. Au passif prétérit : \all{wurde} (singulier) / \all{wurden} (pluriel) \textbf{+ Partizip II en fin de phrase} : \all{Das Erz wurde abgebaut. Die Arbeiter wurden eingestellt.}
\end{attention}

\courssub{Rédaction collective d'un commentaire modèle}

Au tableau, la classe rédige ensemble, sous la conduite du professeur, un commentaire modèle d'environ 120 mots sur le texte support de la leçon (l'extraction de l'or dans une région minière). Voici le résultat attendu.

\begin{textebox}[Kommentar — Modell (rédaction collective)]
Der Text „Gold aus der Erde“ behandelt das Thema Bergbau in einer Region Kameruns.\par
Zunächst beschrieb der Autor, wie das Gold in der Mine entdeckt wurde. Dann erklärte er, dass viele Arbeiter in den Minen beschäftigt wurden und dass das Erz mit einfachen Maschinen abgebaut wurde. Außerdem zeigte er die Probleme: die Flüsse, in denen gewaschen wurde, sind verschmutzt, und die Arbeiter verdienen wenig Geld.\par
Der Text benutzt oft das Passiv im Präteritum, zum Beispiel „wurde abgebaut“, und viele Relativsätze mit Präpositionen wie „die Mine, in der gearbeitet wurde“. Das Wortfeld des Bergbaus (das Erz, der Rohstoff, die Mine) ist sehr präsent.\par
Meiner Meinung nach ist der Bergbau wichtig für die Wirtschaft, jedoch muss die Natur geschützt werden. Abschließend kann man sagen: Fortschritt ja, aber mit Verantwortung.
\end{textebox}

\textbf{Glossaire} : \all{behandeln} (traiter) ; \all{entdecken} (découvrir) ; \all{beschäftigen} (employer) ; \all{verschmutzt} (pollué) ; \all{der Fluss, die Flüsse} (la rivière) ; \all{das Wortfeld, die Wortfelder} (le champ lexical) ; \all{die Wirtschaft} (l'économie) ; \all{die Verantwortung} (la responsabilité).

Comptez : ce modèle fait environ 120 mots. Observez comment chaque étape de l'organigramme y est visible : phrase d'introduction, résumé au prétérit avec \all{zunächst / dann / außerdem}, analyse avec exemples précis tirés du texte, opinion avec \all{meiner Meinung nach} et \all{jedoch}, conclusion avec \all{abschließend}. Remarquez aussi la relative avec préposition \all{die Flüsse, in denen gewaschen wurde} (préposition \all{in} + datif pluriel \all{denen}) et le passif \all{geschützt werden} à l'infinitif après \all{muss}.

\begin{savaistu}[« Nachhaltigkeit », le mot vedette des rapports miniers]
Le mot \all{die Nachhaltigkeit} (la durabilité, le caractère durable) est devenu central dans les rapports des entreprises minières allemandes, autrichiennes et suisses. Il vient de l'adjectif \all{nachhaltig} (durable), lui-même formé sur le verbe \all{nachhalten} (durer, persister). Aujourd'hui, presque chaque rapport annuel d'une grande société minière contient un chapitre intitulé \all{Nachhaltigkeitsbericht} (rapport de durabilité) : l'entreprise y explique comment elle protège l'environnement et les ouvriers. Un mot composé typique de l'allemand — et un excellent mot à placer dans votre opinion personnelle : \all{Ich finde Nachhaltigkeit im Bergbau sehr wichtig.} (Je trouve la durabilité dans l'exploitation minière très importante.)
\end{savaistu}

\courssub{Exercice individuel : à vous d'écrire}

\begin{exoresolu}[Entraînement : compléter un commentaire à trous]
\textbf{Énoncé.} Complétez ce début de commentaire avec les connecteurs et les formes verbales correctes (passif au prétérit) :\par
\all{Der Text behandelt den Bergbau. \ldots\ (d'abord) beschrieb der Autor die Mine. Das Erz \ldots\ (fut extrait) von den Arbeitern abgebaut. \ldots\ (de plus) wurden viele Bäume gefällt. \ldots\ (cependant) bringt der Bergbau auch Geld. \ldots\ (pour conclure) kann man sagen, dass der Text sehr informativ ist.}
\tcblower
\textbf{Solution détaillée.}\par
\all{Der Text behandelt den Bergbau. \textbf{Zunächst} beschrieb der Autor die Mine. Das Erz \textbf{wurde} von den Arbeitern abgebaut. \textbf{Außerdem} wurden viele Bäume gefällt. \textbf{Jedoch} bringt der Bergbau auch Geld. \textbf{Abschließend} kann man sagen, dass der Text sehr informativ ist.}\par
Justifications : \all{zunächst} ouvre l'énumération ; \all{das Erz} est singulier neutre, donc passif prétérit \all{wurde … abgebaut} (participe de \all{abbauen}, verbe à particule séparable : \all{abgebaut}) ; \all{außerdem} ajoute une information ; \all{jedoch} marque l'opposition entre destruction et profit ; \all{abschließend} introduit la conclusion. Remarquez la place du verbe juste après chaque connecteur en tête de phrase : \all{Außerdem \textbf{wurden}…}, \all{Jedoch \textbf{bringt}…}.
\end{exoresolu}

\begin{exercice}[Votre commentaire personnel]
En vous appuyant sur le texte support de la leçon et sur le modèle collectif, rédigez votre propre commentaire de \textbf{100 à 130 mots} en respectant les quatre étapes :\par
1. \all{Einleitung} : une phrase (titre + thème).\par
2. \all{Zusammenfassung} : trois ou quatre phrases au prétérit, avec vos propres mots, en utilisant au moins deux connecteurs (\all{zunächst}, \all{außerdem}).\par
3. \all{Analyse} : deux phrases sur la langue — citez au moins un passif au prétérit et une relative avec préposition relevés dans le texte.\par
4. \all{Meinung} : deux phrases d'opinion avec \all{meiner Meinung nach} ou \all{ich finde, dass…}, et une conclusion avec \all{abschließend} ou \all{zusammenfassend}.
\end{exercice}

\begin{corrige}[Exercice : commentaire personnel — éléments attendus]
Il n'existe pas une seule bonne réponse, mais tout commentaire réussi doit contenir :\par
— une introduction du type \all{Der Text „…“ behandelt das Thema Bergbau / den Abbau von Gold…} ;\par
— un résumé \textbf{au prétérit uniquement}, sans phrase copiée, avec des passifs corrects : \all{Das Gold wurde entdeckt. Die Arbeiter wurden eingestellt. Das Erz wurde verkauft.} ;\par
— une analyse citant la langue du texte : \all{Der Autor benutzt das Passiv im Präteritum („wurde abgebaut“) und Relativsätze mit Präpositionen („die Mine, in der…“).} ;\par
— une opinion argumentée : \all{Meiner Meinung nach ist der Bergbau eine Chance, jedoch nur mit Nachhaltigkeit.} ;\par
— au moins trois connecteurs différents et un vocabulaire du champ minier (\all{das Erz, der Rohstoff, die Mine, der Arbeiter, das Gold}).
\end{corrige}

\courssub{Auto-évaluation et correction croisée}

Avant de rendre votre production, évaluez-la vous-même avec la grille critériée suivante. Attribuez-vous de 0 à 2 points par critère : 0 = absent ou faux, 1 = présent mais imparfait, 2 = réussi.

\begin{center}
\begin{tabularx}{\linewidth}{|X|X|X|}
\hline
\rowcolor{popblueL} \textbf{Critère} & \textbf{Question de contrôle} & \textbf{Points} \\
\hline
Contenu & Les 4 étapes sont-elles présentes et dans l'ordre ? & /2 \\
\hline
Grammaire & Prétérit tenu partout ? Passif correct (\all{wurde + Partizip II}) ? Verbe en 2\textsuperscript{e} position ? & /2 \\
\hline
Vocabulaire & Au moins 4 mots du champ minier, avec le bon article ? & /2 \\
\hline
Cohérence & Au moins 3 connecteurs variés et bien placés ? & /2 \\
\hline
Opinion & Avis personnel clair avec un argument ? & /2 \\
\hline
\end{tabularx}
\end{center}

\begin{experience}[Correction croisée en binômes]
Échangez votre feuille avec votre voisin. Chacun corrige la copie de l'autre en trois passages :\par
1. \textbf{Premier passage} (contenu) : soulignez les quatre étapes ; si une étape manque, notez-la dans la marge.\par
2. \textbf{Deuxième passage} (langue) : entourez chaque verbe au passif et vérifiez \all{wurde/wurden + Partizip II} ; vérifiez qu'aucun parfait (\all{hat … gemacht}) ne s'est glissé dans le résumé ; vérifiez les articles des noms du champ minier.\par
3. \textbf{Troisième passage} (cohérence) : surlignez les connecteurs ; s'il y en a moins de trois, proposez-en un (\all{außerdem}, \all{jedoch}, \all{abschließend}).\par
Rendez la copie avec la grille remplie et une phrase d'encouragement en allemand : \all{Gute Arbeit! Achte auf das Passiv!} (Bon travail ! Fais attention au passif !) Le professeur passe ensuite dans les rangs pour valider les corrections les plus délicates.
\end{experience}

\begin{aretenir}[Les clés d'un bon commentaire]
Un commentaire de Terminale réussi repose sur quatre piliers : une \textbf{structure} en quatre étapes (introduction, résumé, analyse, opinion), un \textbf{temps unique} dans le résumé (le prétérit), des \textbf{connecteurs} qui guident le lecteur (\all{zunächst, außerdem, jedoch, abschließend}) et une \textbf{opinion personnelle} argumentée. Relisez toujours votre production avec la grille critériée avant de la rendre : c'est la différence entre une copie moyenne et une excellente copie.
\end{aretenir}

\coursec{Production orale guidée \& réinvestissement}

Après avoir commenté le texte et produit un paragraphe écrit dans la section précédente, il est temps de passer à l'\cle{oral}. Le programme de Terminale A insiste sur la capacité à \emph{prendre la parole en continu} et à \emph{interagir} : au baccalauréat comme dans la vie réelle (entretien, exposé, débat), vous devrez parler de sujets de société comme l'exploitation minière avec aisance. Cette dernière séquence vous propose trois activités : un jeu de rôle, une prise de position de deux minutes et une auto-évaluation enregistrée. Toutes réinvestissent le lexique du chapitre, le passif au prétérit et les relatives avec préposition.

\courssub{Jeu de rôle : un journaliste interviewe un ancien mineur}

\begin{experience}[Jeu de rôle : „Ein Interview mit einem ehemaligen Bergarbeiter“]
Par groupes de deux. L'élève A joue un journaliste d'une radio de Douala ; l'élève B joue \all{Herr Mbarga}, un ancien mineur de 62 ans qui a travaillé vingt ans dans une mine d'or de l'Est du Cameroun.

\textbf{Questions préparées par le journaliste (modèles) :}
\begin{enumerate}
\item \all{Herr Mbarga, wie lange haben Sie in der Mine gearbeitet?} (Monsieur Mbarga, combien de temps avez-vous travaillé dans la mine ?)
\item \all{Welche Rohstoffe wurden dort gefördert?} (Quelles matières premières y étaient extraites ? — passif prétérit)
\item \all{Wie waren die Arbeitsbedingungen, unter denen die Arbeiter litten?} (Comment étaient les conditions de travail dont souffraient les ouvriers ? — relative avec préposition)
\item \all{Wurden die Arbeiter gut bezahlt?} (Les ouvriers étaient-ils bien payés ? — passif prétérit)
\item \all{Was halten Sie heute vom Bergbau in Kamerun?} (Que pensez-vous aujourd'hui de l'exploitation minière au Cameroun ?)
\end{enumerate}

\textbf{Consignes pour le mineur :} répondre en phrases complètes, utiliser au moins une fois \all{das Erz}, \all{der Rohstoff}, \all{die Mine}, \all{der Arbeiter}, et raconter un souvenir au passif prétérit : \all{Damals wurde das Gold von Hand gewaschen.} (À l'époque, l'or était lavé à la main.)

Durée : 3 minutes d'entretien, puis inversion des rôles.
\end{experience}

\begin{methode}[Préparer des fiches-mots clés (Stichworte) avant de parler]
Un bon orateur ne lit pas un texte rédigé : il s'appuie sur des \cle{Stichworte} (mots clés). Voici la démarche :
\begin{enumerate}
\item \textbf{Structurer} : divisez votre intervention en trois parties — \all{Einleitung} (introduction), \all{Hauptteil} (développement), \all{Schluss} (conclusion).
\item \textbf{Sélectionner} : pour chaque partie, notez 3 à 5 mots clés allemands avec leur article, jamais de phrases entières. Exemple : \all{der Rohstoff, -e / wurde gefördert / die Bedingungen, unter denen …}.
\item \textbf{Vérifier la grammaire} : soulignez sur la fiche les formes obligatoires (ici : 2 passifs prétérit, 1 relative avec préposition) pour ne pas les oublier.
\item \textbf{Marquer la prononciation} : notez en lettres françaises les mots difficiles, ex. \all{Gewinnung} « gue-vi-noung », \all{Arbeiter} « ar-baï-teur ».
\item \textbf{S'entraîner à voix haute} deux fois en regardant uniquement la fiche, puis une fois sans fiche.
\end{enumerate}
\end{methode}

\courssub{Présentation orale : „Mein Standpunkt zum Bergbau in Kamerun“}

Chaque élève prépare une prise de parole de \textbf{deux minutes} sur le sujet : \all{Mein Standpunkt zum Bergbau in Kamerun} (Mon point de vue sur l'exploitation minière au Cameroun). Le contrat de langue est strict :

\begin{propriete}[Contrat de langue de la présentation orale]
Votre intervention de deux minutes doit contenir obligatoirement :
\begin{itemize}
\item \textbf{5 mots du lexique du chapitre} : par exemple \all{das Erz, der Bergbau, die Mine, das Gold, der Rohstoff, der Arbeiter, fördern, die Umwelt} ;
\item \textbf{2 passifs au prétérit} : \all{wurde + Partizip II} ou \all{wurden + Partizip II} ;
\item \textbf{1 phrase relative avec préposition} : \all{die Mine, in der … / die Arbeiter, für die … / der Rohstoff, mit dem …}.
\end{itemize}
Structure attendue : introduction (sujet + opinion), deux arguments, exemple camerounais, conclusion personnelle.
\end{propriete}

\begin{exemplebox}[Exemples de phrases réutilisables (traduits)]
\begin{itemize}
\item \all{Früher wurde hier Gold gefördert, heute wird die Mine für Touristen geöffnet.} « Autrefois, on extrayait de l'or ici ; aujourd'hui, la mine est ouverte aux touristes. » (passif prétérit \all{wurde gefördert} puis passif présent \all{wird geöffnet} : observez le contraste des temps.)
\item \all{Die Rohstoffe, auf die Kamerun stolz ist, wurden oft ohne Kontrolle exportiert.} « Les matières premières dont le Cameroun est fier ont souvent été exportées sans contrôle. » (relative avec préposition \all{auf die} + passif prétérit \all{wurden exportiert}.)
\item \all{Die Arbeiter, für die die Mine die einzige Arbeit war, wurden schlecht bezahlt.} « Les ouvriers pour qui la mine était le seul travail étaient mal payés. »
\item \all{Meiner Meinung nach muss die Umwelt geschützt werden.} « À mon avis, l'environnement doit être protégé. » (passif présent avec modal \all{muss … geschützt werden}.)
\item \all{Früher wurden die Erze von Hand aus der Erde geholt.} « Autrefois, les minerais étaient extraits de la terre à la main. »
\end{itemize}
\end{exemplebox}

\begin{attention}[Piège : \all{fördern} (extraire) ou \all{fördern} (soutenir) ?]
L'allemand possède deux verbes \all{fördern} homonymes, tous deux réguliers (\all{förderte, hat gefördert}) :
\begin{itemize}
\item \all{Gold fördern} = \textbf{extraire} de l'or (sens minier : \all{abbauen, gewinnen}) ;
\item \all{einen Schüler fördern} = \textbf{soutenir, encourager} un élève (sens social : \all{unterstützen}).
\end{itemize}
Seul le contexte permet de trancher. Comparez : \all{In der Mine wurde Erz gefördert.} (Du minerai était extrait dans la mine.) / \all{Die Regierung förderte junge Unternehmer.} (Le gouvernement soutenait de jeunes entrepreneurs.) À l'oral, si vous parlez de la mine, précisez toujours le complément (\all{Gold, Erz, Rohstoffe}) pour lever l'ambiguïté.
\end{attention}

\begin{center}
\begin{tabularx}{\linewidth}{|X|X|X|}
\hline
\rowcolor{popblueL} Français & Deutsch & Remarque \\
\hline
extraire (du minerai) & Erz fördern / abbauen & sens minier \\
\hline
soutenir, encourager & jemanden fördern & sens social \\
\hline
la mine & die Mine, -n & aussi : das Bergwerk, -e \\
\hline
l'exploitation minière & der Bergbau (Sing.) & invariable au singulier \\
\hline
le minerai & das Erz, -e & pluriel : die Erze \\
\hline
la matière première & der Rohstoff, -e & famille : roh + der Stoff \\
\hline
l'ouvrier mineur & der Bergarbeiter, - & composé : Berg + Arbeiter \\
\hline
était extrait (passif prétérit) & wurde gefördert & wurden au pluriel \\
\hline
la mine dans laquelle & die Mine, in der … & relative + préposition \\
\hline
\end{tabularx}
\end{center}

\courssub{Enregistrement audio et auto-évaluation}

\begin{experience}[Enregistrement et auto-évaluation]
Après la présentation, chaque élève \textbf{s'enregistre} (téléphone, dictaphone) en répétant son intervention de deux minutes. Puis il réécoute son enregistrement et remplit la grille d'auto-évaluation :
\begin{enumerate}
\item \textbf{Prononciation} : ai-je bien dit « ch » pour \all{ich}, « vé » pour \all{w} (\all{werden} « vé-è-deun »), « f » pour \all{v} (\all{viel} « fil ») ?
\item \textbf{Fluidité} : ai-je parlé sans longues pauses, sans revenir en arrière ?
\item \textbf{Contrat de langue} : ai-je bien placé mes 5 mots de lexique, mes 2 passifs prétérit et ma relative avec préposition ?
\item \textbf{Place du verbe} : le verbe conjugué est-il en deuxième position, le participe du passif à la fin ?
\end{enumerate}
Note personnelle sur 10, puis un engagement écrit : \all{Nächstes Mal werde ich …} (La prochaine fois, je …).
\end{experience}

\begin{popfigure}
\begin{center}\tcbox[colback=popblue,colframe=popblue,coltext=white,arc=9pt,boxsep=4pt,left=6pt,right=6pt]{\bfseries\small Bewertung der mündlichen Produktion}\end{center}
\vspace{-2pt}
\begin{tcbraster}[raster columns=2,raster equal height=rows,raster column skip=6pt,raster row skip=6pt,enhanced,blankest]
\begin{tcolorbox}[enhanced,colback=white,colframe=poporange,boxrule=0.9pt,arc=3pt,title={Wortschatz 5 Stichworte mit Artikel},fonttitle=\bfseries\scriptsize,coltitle=white,colbacktitle=poporange,left=4pt,right=4pt,top=2pt,bottom=2pt,boxsep=1pt,toptitle=1.5pt,bottomtitle=1.5pt,halign=left]
\begin{itemize}[nosep,leftmargin=*,itemsep=1pt,font=\scriptsize]
\item das Erz die Mine
\item der Rohstoff der Arbeiter
\end{itemize}
\end{tcolorbox}
\begin{tcolorbox}[enhanced,colback=white,colframe=popgreen,boxrule=0.9pt,arc=3pt,title={Grammatik Passiv Präteritum Relativsatz + Präp.},fonttitle=\bfseries\scriptsize,coltitle=white,colbacktitle=popgreen,left=4pt,right=4pt,top=2pt,bottom=2pt,boxsep=1pt,toptitle=1.5pt,bottomtitle=1.5pt,halign=left]
\begin{itemize}[nosep,leftmargin=*,itemsep=1pt,font=\scriptsize]
\item wurde gefördert
\item die Mine, in der …
\end{itemize}
\end{tcolorbox}
\begin{tcolorbox}[enhanced,colback=white,colframe=poppurple,boxrule=0.9pt,arc=3pt,title={Aussprache ch, w, v Flüssigkeit},fonttitle=\bfseries\scriptsize,coltitle=white,colbacktitle=poppurple,left=4pt,right=4pt,top=2pt,bottom=2pt,boxsep=1pt,toptitle=1.5pt,bottomtitle=1.5pt,halign=left]
\begin{itemize}[nosep,leftmargin=*,itemsep=1pt,font=\scriptsize]
\item „vé-è-deun“ werden
\end{itemize}
\end{tcolorbox}
\begin{tcolorbox}[enhanced,colback=white,colframe=popgold,boxrule=0.9pt,arc=3pt,title={Interaktion Fragen verstehen antworten},fonttitle=\bfseries\scriptsize,coltitle=white,colbacktitle=popgold,left=4pt,right=4pt,top=2pt,bottom=2pt,boxsep=1pt,toptitle=1.5pt,bottomtitle=1.5pt,halign=left]
\begin{itemize}[nosep,leftmargin=*,itemsep=1pt,font=\scriptsize]
\item Blickkontakt 2 Minuten
\end{itemize}
\end{tcolorbox}
\end{tcbraster}

\legende{Carte mentale des quatre critères d'évaluation de la production orale : vocabulaire, grammaire, prononciation, interaction.}
\end{popfigure}

\courssub{Texte d'appui pour la prise de parole}

Avant de vous lancer, lisez ce court témoignage : il vous fournit un modèle de prise de position et des formulations réutilisables à l'oral.

\begin{textebox}[„Mein Standpunkt zum Bergbau“ — ein Schüler spricht]
Mein Name ist Etienne und ich komme aus Yaoundé. Heute möchte ich über den Bergbau in Kamerun sprechen. Meiner Meinung nach ist der Bergbau für unser Land sehr wichtig, denn viele Rohstoffe, auf die wir stolz sind, werden hier gefördert: Gold, Eisen und Bauxit. Früher wurde das Gold oft von Hand gewaschen, und die Arbeiter wurden schlecht bezahlt. Die Bedingungen, unter denen sie arbeiteten, waren sehr hart: Die Minen waren gefährlich, und es gab wenig Sicherheit.

Aber ich sehe auch die Probleme. Die Umwelt, in der wir leben, wurde durch den Bergbau oft zerstört: Wälder wurden abgeholzt, und Flüsse wurden verschmutzt. Viele Arbeiter, für die die Mine die einzige Arbeit war, erkrankten.

Ich denke, dass der Staat den Bergbau besser kontrollieren muss. Die Rohstoffe dürfen nicht ohne Kontrolle exportiert werden, und die Arbeiter müssen geschützt werden. Wenn das Geld, mit dem die Minen Gewinn machen, auch in Schulen und Krankenhäuser investiert wird, dann ist der Bergbau eine Chance für Kamerun. Zum Schluss: Der Bergbau ist weder nur gut noch nur schlecht — es kommt darauf an, wie er organisiert wird. Ich hoffe, dass junge Kameruner wie ich eines Tages moderne und sichere Minen leiten werden. Vielen Dank für Ihre Aufmerksamkeit.
\end{textebox}

\textbf{Glossaire :} \all{der Bergbau} (l'exploitation minière) ; \all{der Rohstoff, -e} (la matière première) ; \all{fördern} (extraire) ; \all{das Bauxit} (la bauxite) ; \all{die Bedingung, -en} (la condition) ; \all{hart} (dur) ; \all{gefährlich} (dangereux) ; \all{die Sicherheit} (la sécurité) ; \all{zerstören} (détruire) ; \all{abholzen} (abattre, déboiser) ; \all{verschmutzen} (polluer) ; \all{der Gewinn, -e} (le profit) ; \all{investieren} (investir) ; \all{die Aufmerksamkeit} (l'attention) ; \all{erkranken} (tomber malade).

\textbf{Questions de compréhension orale :}
\begin{enumerate}
\item \all{Welche Rohstoffe werden in Kamerun gefördert?} (Quelles matières premières sont extraites au Cameroun ?)
\item \all{Wie wurden die Arbeiter früher bezahlt?} (Comment les ouvriers étaient-ils payés autrefois ?)
\item \all{Was wurde durch den Bergbau zerstört?} (Qu'est-ce qui a été détruit par l'exploitation minière ?)
\item \all{Was muss der Staat Ihrer Meinung nach tun?} (Que doit faire l'État selon vous ?)
\end{enumerate}

\begin{exoresolu}[Construire une phrase complexe pour l'oral]
Énoncé : combinez les deux phrases suivantes en une seule phrase avec une relative à préposition, puis mettez le passif au prétérit : \all{Die Arbeiter arbeiteten in der Mine. In dieser Mine wurde Gold gefördert.}

\tcblower
Solution détaillée :
\begin{enumerate}
\item Le mot commun est \all{die Mine} (féminin). La préposition \all{in} avec localisation exige le datif : \all{in der}.
\item On remplace donc \all{in dieser Mine} par le relatif \all{in der}, placé juste après le nom de référence, et le verbe conjugué de la relative va en fin de proposition.
\item Le passif prétérit est déjà présent : \all{wurde gefördert} (3\textsuperscript{e} personne du singulier, car \all{das Gold} est singulier).
\end{enumerate}
Résultat : \all{Die Arbeiter arbeiteten in der Mine, in der Gold gefördert wurde.} « Les ouvriers travaillaient dans la mine où l'or était extrait. » Remarquez que dans la relative, \all{wurde} passe à la fin : \all{in der Gold gefördert wurde}, jamais \all{in der wurde Gold gefördert}.
\end{exoresolu}

\begin{attention}[Erreurs fréquentes des francophones à l'oral]
\begin{itemize}
\item \textbf{Place du verbe} : ne dites jamais \all{Früher das Gold wurde gefördert} ; après un complément en tête, le verbe reste en deuxième position : \all{Früher wurde das Gold gefördert.}
\item \textbf{Auxiliaire du passif} : le passif se construit avec \all{werden}, pas avec \all{sein}. \all{Das Gold wurde gefördert} (l'or était extrait), jamais \all{das Gold war gefördert}.
\item \textbf{Faux ami} : \all{die Mine} ne signifie pas « la mine » au sens de « expression du visage » ; c'est bien la mine de minerais. Mais attention : \all{die Miene} (l'expression du visage) existe !
\item \textbf{Majuscules} : à l'écrit comme dans vos fiches, tous les noms prennent la majuscule : \all{das Erz, der Arbeiter, die Umwelt}.
\item \textbf{Genre} : \all{das Erz} (neutre), pas \all{der Erz} ; vérifiez l'article sur vos Stichworte avant de parler.
\end{itemize}
\end{attention}

\courssub{Feedback collectif et bilan de la leçon}

Après chaque passage à l'oral, la classe applique la règle du \emph{sandwich} : deux \textbf{points forts} (\all{Was war gut?}) puis un \textbf{axe d'amélioration} (\all{Was kann besser werden?}). Exemples de formulations en allemand : \all{Deine Aussprache war sehr klar.} (Ta prononciation était très claire.) / \all{Du hast das Passiv gut benutzt.} (Tu as bien utilisé le passif.) / \all{Achte auf die Position des Verbs.} (Fais attention à la place du verbe.)

\begin{savaistu}[Bergbau ou Rohstoffgewinnung ?]
En Allemagne, on parle couramment du \all{Bergbau} (littéralement « construction sur la montagne »), par exemple de la région de la Ruhr, ancien bassin houiller. En Suisse, en revanche, où l'extraction minière est presque inexistante, le terme \all{Bergbau} est rare : on préfère l'expression plus large \all{die Rohstoffgewinnung} (« l'obtention des matières premières »), qui couvre aussi le gravier, le sel et la pierre. En Autriche, on rencontre aussi \all{der Erzbergbau} pour l'extraction du minerai de fer, notamment sur l'\all{Erzberg} en Styrie, la « montagne de minerai ». Un même concept, trois pays, trois habitudes de vocabulaire !
\end{savaistu}

\begin{propriete}[Checklist finale avant le Bac — chapitre « Bergbau und Umwelt »]
Avant l'épreuve, vérifiez que vous savez :
\begin{itemize}
\item \textbf{Vocabulaire} : citer avec article et pluriel \all{das Erz, -e ; der Bergbau (Sg.) ; die Mine, -n ; das Gold ; der Rohstoff, -e ; der Arbeiter, - ; der Bergarbeiter, - ; die Umwelt (Sg.)} et employer \all{fördern} dans ses deux sens ;
\item \textbf{Passif prétérit} : former \all{wurde/wurden + Partizip II} (\all{wurde gefördert, wurden exportiert, wurde zerstört}) et le distinguer du passif présent (\all{wird gefördert}) ;
\item \textbf{Relatives avec préposition} : construire \all{die Mine, in der … ; die Arbeiter, für die … ; das Geld, mit dem …} avec le verbe en fin de proposition ;
\item \textbf{Structure du commentaire} : introduction (présenter le texte et le thème), développement (idées + exemples + lexique), conclusion (opinion personnelle avec \all{meiner Meinung nach}) ;
\item \textbf{Oral} : parler deux minutes à partir de Stichworte, avec 5 mots du lexique, 2 passifs prétérit et 1 relative avec préposition.
\end{itemize}
\end{propriete}

\begin{exercice}[Pour s'entraîner à la maison]
\begin{enumerate}
\item Enregistrez-vous pendant deux minutes sur \all{Mein Standpunkt zum Bergbau in Kamerun} en respectant le contrat de langue, puis remplissez la grille d'auto-évaluation.
\item Transformez en passif prétérit : \all{Man förderte Gold. / Man verschmutzte die Flüsse. / Man bezahlte die Arbeiter schlecht.}
\item Combinez avec une relative à préposition : \all{Die Arbeiter träumten von einer besseren Zukunft. Sie arbeiteten für diese Zukunft.}
\end{enumerate}
\end{exercice}

\begin{corrige}[Exercice « Pour s'entraîner à la maison »]
\begin{enumerate}
\item Auto-évaluation individuelle : vérifier la présence des 5 mots de lexique, des 2 passifs prétérit et de la relative avec préposition ; écouter la prononciation de \all{w} (« vé ») et \all{ch}.
\item \all{Gold wurde gefördert. / Die Flüsse wurden verschmutzt. / Die Arbeiter wurden schlecht bezahlt.} (Attention à l'accord : \all{wurde} au singulier, \all{wurden} au pluriel.)
\item \all{Die Arbeiter träumten von einer besseren Zukunft, für die sie arbeiteten.} « Les ouvriers rêvaient d'un avenir meilleur pour lequel ils travaillaient. » (Préposition \all{für} + accusatif : \all{für die} ; verbe conjugué en fin de relative.)
\end{enumerate}
\end{corrige}

Vous voilà au terme de la leçon : vous savez désormais comprendre un texte sur l'extraction des minerais, manier le passif au prétérit et les relatives avec préposition, et surtout \emph{prendre la parole} de manière structurée et évaluée. Ces compétences orales vous serviront directement au baccalauréat et dans toute situation de communication professionnelle.

\newpage

\coursec{Activité d'intégration}

\courssub{Objectifs de l'activité}

Cette activité d'intégration clôt la leçon sur l'extraction des minerais. Elle vise à mobiliser, dans une situation de communication authentique et motivante, l'ensemble des savoir-faire travaillés :

\begin{itemize}
\item \textbf{Compréhension} : réutiliser les informations du texte support sur l'extraction des minerais (techniques, dangers, enjeux économiques et écologiques).
\item \textbf{Lexique} : employer correctement le vocabulaire thématique (\all{das Erz, der Bergbau, die Mine, das Gold, der Rohstoff, der Arbeiter, die Umweltverschmutzung}...).
\item \textbf{Grammaire} : produire des phrases au \cle{passif prétérit} (\all{wurde abgebaut, wurden entdeckt}) et des \cle{propositions relatives avec préposition} (\all{die Mine, in der...}).
\item \textbf{Expression orale} : défendre une opinion, argumenter, réagir aux propos d'autrui dans un débat structuré.
\item \textbf{Compétence citoyenne} : prendre position sur un problème réel de société qui concerne directement le Cameroun.
\end{itemize}

\courssub{Situation}

Le lycée de Ngaoundéré participe à la semaine nationale de l'environnement. Une chaîne de télévision locale a invité les élèves de la classe de Terminale A à enregistrer un débat de vingt minutes sur un sujet d'actualité brûlante : l'exploitation minière au Cameroun (or dans l'Est et l'Adamaoua, fer à Mbalam, bauxite à Minim-Martap). Le journaliste animateur de l'émission a envoyé aux élèves le courrier suivant :

\begin{textebox}[Einladung zur Fernsehdebatte]
Liebe Schülerinnen und Schüler,

unsere Sendung „Jugend debattiert" beschäftigt sich nächste Woche mit dem Thema: \textbf{„Bergbau – Fluch oder Segen für Kamerun?"}

Im Osten des Landes wurde viel Gold gefunden, und in vielen Dörfern wurden neue Minen geöffnet. Aber die Flüsse, in denen die Menschen früher gefischt haben, sind jetzt verschmutzt. Viele Arbeiter, die in den Minen arbeiten, verdienen wenig Geld und haben keine Sicherheit.

Wir laden Sie ein, darüber zu diskutieren: Bringt der Bergbau Arbeit und Entwicklung – oder zerstört er die Umwelt und das Leben der Dorfbewohner?

Bereiten Sie Ihre Argumente vor! Wir erwarten klare Meinungen, gute Beispiele und korrektes Deutsch.

Mit freundlichen Grüßen\\
Ihr Moderator
\end{textebox}

\textbf{Glossaire :} der Fluch (\emph{la malédiction}), der Segen (\emph{la bénédiction}), verschmutzt (\emph{pollué}), die Sicherheit (\emph{la sécurité}), die Dorfbewohner (\emph{les villageois}), die Meinung, -en (\emph{l'opinion}).

\courssub{Consignes}

\begin{enumerate}
\item \textbf{Relecture ciblée (10 min).} En binômes, relisez le texte support de la leçon et soulignez : deux phrases au passif prétérit, deux relatives avec préposition, et cinq mots du champ lexical de la mine. Notez-les sur votre fiche de préparation.

\item \textbf{Constitution des équipes.} La classe est divisée en deux groupes : le groupe \all{„Pro Bergbau"} (pour) et le groupe \all{„Contra Bergbau"} (contre). Un élève volontaire joue le rôle du \all{Moderator} (modérateur) ; il prépare les questions d'ouverture et les relances.

\item \textbf{Préparation écrite des arguments (20 min).} Chaque groupe rédige \textbf{cinq arguments} en allemand. Contraintes obligatoires : chaque argument contient au moins un mot du lexique thématique ; le groupe doit utiliser \textbf{au moins trois fois le passif prétérit} et \textbf{au moins deux relatives avec préposition}. Exemple de trame : \all{In der Region, in der viel Gold gefunden wurde, ...}

\item \textbf{Préparation des cartes d'expression d'opinion.} Chaque élève écrit sur une carte trois formules de prise de parole : une pour donner son avis (\all{Meiner Meinung nach...}), une pour être d'accord (\all{Ich stimme dir zu, weil...}), une pour contredire (\all{Ich bin anderer Meinung, denn...}).

\item \textbf{Le débat filmé (20 min).} Le modérateur ouvre le débat, pose une question à chaque groupe, puis donne la parole tour à tour. Chaque intervention dure au maximum une minute. Le débat est filmé avec un téléphone ou une caméra.

\item \textbf{Analyse collective (15 min).} On regarde l'enregistrement. La classe évalue les interventions selon la grille : exactitude grammaticale (passif prétérit, relatives), richesse du lexique, clarté de l'argumentation, prononciation et prise de parole.

\item \textbf{Bilan écrit individuel.} Chaque élève rédige ensuite un court paragraphe (6 à 8 phrases) intitulé \all{„Meine Meinung zum Bergbau"}, en réutilisant au moins une fois le passif prétérit et une relative avec préposition.
\end{enumerate}

\courssub{Ressources à mobiliser}

\begin{aretenir}[Boîte à outils linguistique]
\textbf{Passif prétérit :} \emph{Präteritum de werden} + Partizip II : \all{wurde / wurden + abgebaut / gefunden / geöffnet / verschmutzt}.\\
\all{Das Gold wurde im Osten gefunden.} « L'or a été découvert dans l'Est. »

\textbf{Relatives avec préposition :} préposition + pronom relatif (cas exigé par la préposition) : \all{die Mine, in der die Arbeiter arbeiten} ; \all{das Erz, aus dem Metall gewonnen wird}.

\textbf{Lexique :} der Bergbau, die Mine (-n), das Erz (-e), das Gold, der Rohstoff (-e), der Arbeiter (-), abbauen, entdecken, verschmutzen, die Umwelt, die Entwicklung, der Arbeitsplatz, die Arbeitsplätze.

\textbf{Expression d'opinion :} \all{Meiner Meinung nach... / Ich bin der Ansicht, dass... / Ich stimme (nicht) zu, weil... / Auf der einen Seite..., auf der anderen Seite...}
\end{aretenir}

\begin{attention}[Pièges à éviter]
\begin{itemize}
\item Ne dites pas \all{*das Gold ist gefunden wurden} : au prétérit passif, on conjugue \all{werden} au prétérit (\all{wurde}), pas au présent.
\item Dans la relative avec préposition, la préposition se place \textbf{devant} le pronom relatif et le verbe va à la fin : \all{der Fluss, in dem früher gefischt wurde}.
\item Attention au genre : \all{das Erz, die Mine, der Rohstoff, das Gold}.
\item \all{Meiner Meinung nach} est suivi du verbe en deuxième position et du sujet au nominatif : \all{Meiner Meinung nach ist der Bergbau wichtig.}
\end{itemize}
\end{attention}

\begin{methode}[Conduite du débat télévisé]
\textbf{Phase 1 – Ouverture (Modérateur) :} \all{„Herzlich willkommen zu ‚Jugend debattiert'. Heute: Bergbau in Kamerun – Fluch oder Segen? Gruppe Pro, Ihr erstes Argument bitte."}\\
\textbf{Phase 2 – Arguments (Alternance Pro/Contra) :} Chaque groupe présente un argument préparé. Le modérateur note les structures cibles entendues.\\
\textbf{Phase 3 – Libre échange (Modérateur relance) :} \all{„Gruppe Contra, wie reagieren Sie darauf?"} – Les élèves utilisent leurs cartes d'expression (\all{Ich stimme zu / Ich widerspreche}).\\
\textbf{Phase 4 – Conclusion (Modérateur) :} \all{„Wir fassen zusammen: ... Vielen Dank für die spannende Diskussion."}
\end{methode}

\begin{savaistu}[Le mining au Cameroun et en Allemagne]
Le Cameroun possède d'importantes ressources : or (Est, Adamaoua), fer (Mbalam), bauxite (Minim-Martap), cobalt, nickel. L'exploitation est souvent artisanale (orpailleurs) ou confiée à des multinationales. Enjeux : recettes budgétaires vs pollution au mercure, déforestation, conflits fonciers.
En Allemagne, la Ruhr (\all{das Ruhrgebiet}) a été le cœur de l'industrie charbonnière (\all{der Steinkohlebergbau}) jusqu'à la fermeture de la dernière mine en 2018 (\all{Schlussstrich unter den Steinkohlebergbau}). La Lusace (\all{die Lausitz}) exploite encore le lignite (\all{die Braunkohle}), mais la sortie du charbon (\all{der Kohleausstieg}) est votée pour 2038. La reconversion (\all{der Strukturwandel}) est un modèle étudié pour les régions minières camerounaises.
\end{savaistu}

\courssub{Proposition de production et corrigé commenté}

\begin{exoresolu}[Fiche d'arguments du groupe „Pro Bergbau" (modèle)]
\textbf{Énoncé.} Rédiger cinq arguments en faveur de l'exploitation minière, avec les contraintes grammaticales de la consigne 3.
\tcblower
\textbf{Production modèle :}

\all{1. Meiner Meinung nach ist der Bergbau wichtig für die Entwicklung unseres Landes.}\\[2pt]
\all{2. In den Regionen, in denen Gold gefunden wurde, wurden viele Arbeitsplätze geschaffen.}\\[2pt]
\all{3. Das Erz, das in den Minen abgebaut wurde, wurde exportiert, und der Staat hat Geld verdient.}\\[2pt]
\all{4. Mit dem Geld, das der Staat aus dem Bergbau bekommt, wurden Straßen und Schulen gebaut.}\\[2pt]
\all{5. Viele junge Leute, die früher keine Arbeit hatten, arbeiten heute als Arbeiter in den Minen.}

\textbf{Commentaire des choix de langue :}
\begin{itemize}
\item \textbf{Argument 2} : passif prétérit au pluriel (\all{wurden geschaffen}) et relative avec préposition (\all{in denen}, datif pluriel après \all{in}) — double contrainte remplie dans une seule phrase.
\item \textbf{Argument 3} : relative simple au nominatif (\all{das in den Minen abgebaut wurde}) avec passif prétérit à l'intérieur de la relative ; le verbe \all{abbauen} est séparable, d'où le participe \all{abgebaut}.
\item \textbf{Argument 4} : relative avec préposition \all{mit} + datif (\all{mit dem Geld}) ; le passif prétérit \all{wurden gebaut} clôt la phrase, verbe en position finale.
\item \textbf{Argument 5} : relative simple (\all{die früher keine Arbeit hatten}) pour varier les structures et éviter la monotonie.
\item L'argument 1 ouvre avec une formule d'opinion (\all{Meiner Meinung nach}) suivie du verbe en deuxième position — ordre des mots correct.
\end{itemize}
\end{exoresolu}

\begin{exoresolu}[Fiche d'arguments du groupe „Contra Bergbau" (modèle)]
\textbf{Énoncé.} Rédiger cinq arguments contre l'exploitation minière, avec les mêmes contraintes.
\tcblower
\textbf{Production modèle :}

\all{1. Ich bin der Ansicht, dass der Bergbau ein Fluch für viele Dörfer ist.}\\[2pt]
\all{2. Die Flüsse, in denen die Menschen früher gefischt haben, wurden durch den Bergbau verschmutzt.}\\[2pt]
\all{3. Die Arbeiter, für die keine Sicherheit garantiert wurde, wurden oft krank.}\\[2pt]
\all{4. Das Land, aus dem das Gold abgebaut wurde, blieb arm, weil die ausländischen Firmen alles Geld mitgenommen haben.}\\[2pt]
\all{5. Die Wälder, in denen viele Tiere lebten, wurden zerstört.}

\textbf{Commentaire :} l'argument 3 illustre la relative avec préposition exigeant l'accusatif (\all{für die}) ; l'argument 4 combine relative avec préposition (\all{aus dem}, datif) et passif prétérit (\all{wurde abgebaut}). On remarque l'alternance entre passif prétérit (\all{wurden verschmutzt, wurden zerstört}) et actif au parfait (\all{haben gefischt, lebten}) pour un style naturel et varié.
\end{exoresolu}

\courssub{Bilan de l'activité}

À l'issue du débat, l'enseignant fait le point avec la classe sur trois plans :

\begin{itemize}
\item \textbf{Sur la langue :} les élèves ont réinvesti spontanément le passif prétérit et les relatives avec préposition dans des échanges réels, ce qui montre que ces structures sont en voie d'acquisition. Les erreurs relevées pendant le visionnage (accord du sujet avec \all{wurde/wurden}, place du verbe dans la relative) font l'objet d'une courte remédiation collective au tableau.
\item \textbf{Sur le lexique :} le vocabulaire de la mine a été employé dans un contexte signifiant ; les mots les plus utilisés (\all{Bergbau, Mine, Arbeiter, Umwelt}) sont désormais ancrés.
\item \textbf{Sur la citoyenneté :} le débat a permis de comprendre que l'exploitation minière est un enjeu complexe, ni totalement bénéfique ni totalement néfaste, et qu'elle exige des règles strictes pour protéger les populations et l'environnement — une compétence transversale essentielle du programme de Terminale.
\end{itemize}

Le paragraphe individuel \all{„Meine Meinung zum Bergbau"} est ramassé et évalué selon trois critères : correction grammaticale (passif prétérit, relatives), pertinence du lexique, cohérence de l'argumentation. Cette production écrite sert de trace d'apprentissage et pourra être réutilisée lors des révisions du baccalauréat.

\newpage

\coursec{Méthodes et exercices résolus}

\coursec{AL12 : L'extraction des minerais}

\courssub{Mise en situation \& Texte support}

\begin{textebox}[L'exploitation de l'or dans l'Est-Cameroun — Texte authentique]
\all{In der Region Ost Kameruns wird seit über zehn Jahren Gold im industriellen Maßstab abgebaut. Früher förderten Kleinbergleute das Erz von Hand aus schmalen Schächten, doch heute nutzen große Firmen schwere Maschinen. Die Arbeiter, die oft aus den umliegenden Dörfern stammen, arbeiten unter schwierigen Bedingungen: Hitze, Staub und Lärm bestimmen den Alltag in den Minen. Der Rohstoff, aus dem Schmuck und Elektronikbauteile hergestellt werden, wird fast vollständig nach Europa und Asien exportiert. Obwohl der Bergbau dem Staat Steuereinnahmen bringt, bleibt die Region arm; die Straßen sind schlecht, und die Schulen fehlen. Ein Bergmann, mit dem unser Reporter gesprochen hat, klagt: „Wir riskieren unsere Gesundheit, aber der Lohn reicht kaum zum Leben.“ Die Regierung verspricht nun, einen Teil der Gewinne vor Ort zu investieren.}
\end{textebox}

\begin{savaistu}[Contexte : Cameroun et ressources minières]
Le Cameroun possède d'importantes ressources minières (or, bauxite, fer, cobalt, diamant). La région de l'Est (Bétaré-Oya, Batouri) est le cœur de l'exploitation aurifère artisanale et semi-industrielle. L'Allemagne, par sa coopération technique (GIZ, BGR), accompagne la formalisation du secteur et la protection de l'environnement. Le vocabulaire \all{der Bergbau}, \all{der Rohstoff}, \all{die Förderung} est central dans les accords bilatéraux.
\end{savaistu}

\courssub{Compréhension globale et détaillée}

\begin{methode}[Lire et exploiter un texte technique : étapes obligatoires]
\begin{enumerate}
\item \textbf{Première lecture (globale).} Lis le texte une fois sans dictionnaire. Repère : le thème (\all{der Bergbau}), le lieu (\all{in der Region Ost Kameruns}), les acteurs (\all{die Arbeiter, die Firmen, der Staat, die Regierung}) et le ton (informatif, critique, témoignant d'un paradoxe : richesses du sol vs pauvreté locale).
\item \textbf{Repérage du champ lexical.} Souligne les noms-clés avec leur article : \all{das Erz, die Mine, der Schacht, das Gold, der Rohstoff, der Bergbau, der Arbeiter, der Bergmann, die Förderung, die Bedingungen}. Note les familles : \all{fördern $\rightarrow$ die Förderung}, \all{abbauen $\rightarrow$ der Abbau}, \all{arbeiten $\rightarrow$ die Arbeitsbedingungen}.
\item \textbf{Deuxième lecture (détaillée).} Pour chaque question, localise le paragraphe, surligne la phrase-piège, puis reformule \emph{avec tes propres mots} en allemand. N'utilise pas le mot-à-mot du texte.
\item \textbf{Rédaction de la réponse.} Phrase complète : sujet + verbe conjugué (2\up{e} position). Respecte le temps de la question (souvent Präteritum ou Perfekt dans un récit, Présent pour des faits actuels). Particule séparable en fin de phrase principale (\all{Gold abbauen $\rightarrow$ wird ... abgebaut}).
\item \textbf{Vérification finale.} Majuscules à \textbf{tous} les noms, accords sujet-verbe, cas après préposition (\all{mit} + Dativ), ordre V2 / verbe final.
\end{enumerate}
\end{methode}

\begin{exoresolu}[Questions de compréhension sur le texte support]
\textbf{Texte :} (voir Textebox ci-dessus)

\textbf{Énoncé :} Beantworten Sie auf Deutsch in ganzen Sätzen :
\begin{enumerate}
\item \all{Seit wann und wo wird Gold im industriellen Maßstab abgebaut?}
\item \all{Wie wurde das Erz früher gefördert und wie geschieht es heute?}
\item \all{Unter welchen Bedingungen arbeiten die Arbeiter in den Minen?}
\item \all{Was passiert mit dem geförderten Gold?}
\item \all{Welchen Widerspruch (Paradoxon) beschreibt der Text am Ende?}
\end{enumerate}
\tcblower
\textbf{Raisonnement pas à pas :}
\begin{enumerate}
\item Question \all{Seit wann / wo} $\rightarrow$ 1\up{e} phrase : \all{seit über zehn Jahren}, \all{in der Region Ost Kameruns}. Réponse au Présent passif (fait actuel).
\item Question \all{Wie ... früher / heute} $\rightarrow$ Contraste \all{früher} (Präteritum actif \all{förderten}) vs \all{heute} (Présent actif \all{nutzen}). Réponse doit rendre ce contraste.
\item Question \all{Unter welchen Bedingungen} $\rightarrow$ Relatif \all{die ... stammen} + énumération \all{Hitze, Staub, Lärm}.
\item Question \all{Was passiert mit ...} $\rightarrow$ Passif \all{wird ... exportiert}, but \all{fast vollständig}.
\item Question \all{Welchen Widerspruch} $\rightarrow$ Concession \all{Obwohl ... bringt, bleibt ... arm}. Synthèse.
\end{enumerate}

\textbf{Réponses rédigées (modèles) :}
\begin{enumerate}
\item \all{Gold wird in der Region Ost Kameruns seit über zehn Jahren im industriellen Maßstab abgebaut.}
\item \all{Früher förderten Kleinbergleute das Erz von Hand aus schmalen Schächten, aber heute nutzen große Firmen schwere Maschinen.}
\item \all{Die Arbeiter arbeiten unter schwierigen Bedingungen, denn Hitze, Staub und Lärm bestimmen den Alltag in den Minen.}
\item \all{Der Rohstoff, aus dem Schmuck und Elektronikbauteile hergestellt werden, wird fast vollständig nach Europa und Asien exportiert.}
\item \all{Obwohl der Bergbau dem Staat Steuereinnahmen bringt, bleibt die Region arm, weil die Infrastruktur fehlt.}
\end{enumerate}
\textbf{Conclusion :} Chaque réponse est une phrase autonome, reformulée, grammaticalement correcte (cas, temps, place du verbe).
\end{exoresolu}

\courssub{Lexique thématique par champs}

\begin{methode}[Mémoriser et réutiliser le lexique du Bergbau]
\begin{enumerate}
\item \textbf{Nom + Article + Pluriel} : apprends chaque entrée du tableau ci-dessous par cœur.
\item \textbf{Familles de mots (Derivation) :} \all{der Berg} $\rightarrow$ \all{der Bergbau, der Bergmann, das Bergwerk, bergbaulich} ; \all{arbeiten} $\rightarrow$ \all{der Arbeiter, die Arbeit, die Arbeitsbedingungen, arbeitslos} ; \all{fördern} $\rightarrow$ \all{die Förderung, der Förderer, förderfähig}.
\item \textbf{Collocations (associations figées) :} \all{Gold abbauen / fördern}, \all{Erz fördern / transportieren}, \all{eine Mine betreiben / schließen}, \all{Rohstoffe exportieren / verarbeiten}, \all{unter Tage arbeiten}.
\item \textbf{Ancrage camerounais :} \all{In Bétaré-Oya gibt es viele Goldminen. Der Bergbau ist wichtig für die Wirtschaft der Region Ost.}
\end{enumerate}
\end{methode}

\begin{propriete}[Vocabulaire de base : der Bergbau — Noms, articles, pluriels, familles]
\begin{center}
\begin{tabularx}{\linewidth}{|X|X|X|X|X|}
\hline
\rowcolor{popblueL} \textbf{Deutsch} & \textbf{Artikel} & \textbf{Plural} & \textbf{Français} & \textbf{Famille / Dérivés} \\
\hline
Erz & das & die Erze & le minerai & erzreich, erzarm, erzhaltig \\
Bergbau & der & — & l'exploitation minière & Bergmann, Bergwerk, bergbaulich \\
Mine & die & die Minen & la mine & Minenschacht, Minenarbeiter \\
Gold & das & — & l'or & golden, Goldgräber, Goldvorkommen \\
Rohstoff & der & die Rohstoffe & la matière première & roh, Rohstoffreichtum, Rohstoffbasis \\
Arbeiter & der & die Arbeiter & l'ouvrier & arbeiten, Arbeit, Arbeitsbedingungen, arbeitslos \\
Bergmann & der & die Bergleute & le mineur & Bergwerk, bergmännisch, Bergmannslied \\
Förderung & die & die Förderungen & l'extraction, la promotion & fördern, Förderanlage, Förderband \\
Schacht & der & die Schächte & le puits (de mine) & Schachtanlage, schachtförmig \\
Umwelt & die & — & l'environnement & umweltfreundlich, Umweltschutz, Umweltbelastung \\
\hline
\end{tabularx}
\end{center}
\end{propriete}

\begin{exoresolu}[Lexique : compléter, dériver, produire]
\textbf{Énoncé :} a) Ergänzen Sie mit dem passenden Wort (Artikel + Form) : \all{Erz, Mine, Arbeiter, Rohstoff, Bergbau, Förderung, Bergmann, Schacht} — 1. \all{Die \_\_\_ arbeiten acht Stunden pro Tag \textbf{unter Tage}.} 2. \all{In dieser \_\_\_ wird seit 2015 Gold gefördert.} 3. \all{\_\_\_ ist ein wichtiger Wirtschaftszweig in der Region Ost.} 4. \all{Die \_\_\_ des Erzes erfolgt heute mit Maschinen.} 5. \all{Ein \_\_\_ trägt einen Helm und eine Lampe.} b) Bilden Sie je einen eigenen Satz mit \all{das Erz} et \all{der Rohstoff} (contexte camerounais).
\tcblower
\textbf{Raisonnement :} 1. Sujet pluriel, sens « ouvriers » $\rightarrow$ \all{die Arbeiter}. 2. \all{In dieser} (Dat sg fém) + lieu extraction $\rightarrow$ \all{Mine}. 3. Sujet sg masc, activité économique $\rightarrow$ \all{der Bergbau}. 4. Nom fém sg (die Förderung), verbe \all{erfolgen} $\rightarrow$ \all{die Förderung}. 5. « Mineur » (métier) $\rightarrow$ \all{der Bergmann} (pas \all{Minderjährige}!).

\textbf{Réponses :}
\begin{enumerate}
\item \all{Die \textbf{Arbeiter} arbeiten acht Stunden pro Tag unter Tage.}
\item \all{In dieser \textbf{Mine} wird seit 2015 Gold gefördert.}
\item \all{\textbf{Der Bergbau} ist ein wichtiger Wirtschaftszweig in der Region Ost.}
\item \all{Die \textbf{Förderung} des Erzes erfolgt heute mit Maschinen.}
\item \all{Ein \textbf{Bergmann} trägt einen Helm und eine Lampe.}
\end{enumerate}
b) \all{Das Erz, das in den Minen von Bétaré-Oya gefunden wird, wird nach Douala transportiert.} \\
\all{Der Rohstoff Gold ist für die kamerunische Wirtschaft von großer Bedeutung, aber die lokale Bevölkerung profitiert zu wenig davon.}
\end{exoresolu}

\begin{popfigure}
\begin{center}\tcbox[colback=popink,colframe=popink,coltext=white,arc=9pt,boxsep=4pt,left=6pt,right=6pt]{\bfseries\small Der Bergbau}\end{center}
\vspace{-2pt}
\begin{tcbraster}[raster columns=2,raster equal height=rows,raster column skip=6pt,raster row skip=6pt,enhanced,blankest]
\begin{tcolorbox}[enhanced,colback=white,colframe=poporange,boxrule=0.9pt,arc=3pt,title={Orte},fonttitle=\bfseries\scriptsize,coltitle=white,colbacktitle=poporange,left=4pt,right=4pt,top=2pt,bottom=2pt,boxsep=1pt,toptitle=1.5pt,bottomtitle=1.5pt,halign=left]
\begin{itemize}[nosep,leftmargin=*,itemsep=1pt,font=\scriptsize]
\item die Mine
\item der Schacht
\item das Bergwerk
\end{itemize}
\end{tcolorbox}
\begin{tcolorbox}[enhanced,colback=white,colframe=popgreen,boxrule=0.9pt,arc=3pt,title={Personen},fonttitle=\bfseries\scriptsize,coltitle=white,colbacktitle=popgreen,left=4pt,right=4pt,top=2pt,bottom=2pt,boxsep=1pt,toptitle=1.5pt,bottomtitle=1.5pt,halign=left]
\begin{itemize}[nosep,leftmargin=*,itemsep=1pt,font=\scriptsize]
\item der Bergmann
\item der Arbeiter
\item der Ingenieur
\end{itemize}
\end{tcolorbox}
\begin{tcolorbox}[enhanced,colback=white,colframe=poppurple,boxrule=0.9pt,arc=3pt,title={Produkte},fonttitle=\bfseries\scriptsize,coltitle=white,colbacktitle=poppurple,left=4pt,right=4pt,top=2pt,bottom=2pt,boxsep=1pt,toptitle=1.5pt,bottomtitle=1.5pt,halign=left]
\begin{itemize}[nosep,leftmargin=*,itemsep=1pt,font=\scriptsize]
\item das Erz
\item das Gold
\item der Rohstoff
\end{itemize}
\end{tcolorbox}
\begin{tcolorbox}[enhanced,colback=white,colframe=poppink,boxrule=0.9pt,arc=3pt,title={Verben},fonttitle=\bfseries\scriptsize,coltitle=white,colbacktitle=poppink,left=4pt,right=4pt,top=2pt,bottom=2pt,boxsep=1pt,toptitle=1.5pt,bottomtitle=1.5pt,halign=left]
\begin{itemize}[nosep,leftmargin=*,itemsep=1pt,font=\scriptsize]
\item abbauen
\item fördern
\item exportieren
\end{itemize}
\end{tcolorbox}
\end{tcbraster}

\legende{Champ lexical : Der Bergbau (carte mentale)}
\end{popfigure}

\courssub{Grammaire : Le passif au prétérit}

\begin{propriete}[Conjugaison : werden au Präteritum (base du passif)]
\begin{center}
\begin{tabular}{|c|c|c|}
\hline
\rowcolor{popblueL} \textbf{Person} & \textbf{werden (Präteritum)} & \textbf{Exemple passif : \emph{gefunden}} \\
\hline
ich & wurde & \all{ich wurde gefunden} \\
du & wurdest & \all{du wurdest gefunden} \\
er / sie / es & \textbf{wurde} & \all{er wurde gefunden} \\
wir & wurden & \all{wir wurden gefunden} \\
ihr & wurdet & \all{ihr wurdet gefunden} \\
sie / Sie & \textbf{wurden} & \all{sie wurden gefunden} \\
\hline
\end{tabular}
\end{center}
\emph{Rappel :} Participe II en position finale. Verbes forts : \all{finden $\rightarrow$ gefunden}, \all{schließen $\rightarrow$ geschlossen}, \all{bringen $\rightarrow$ gebracht}, \all{geben $\rightarrow$ gegeben}. Verbes faibles : \all{fördern $\rightarrow$ gefördert}, \all{bauen $\rightarrow$ gebaut}, \all{exportieren $\rightarrow$ exportiert}.
\end{propriete}

\begin{methode}[Former et employer le passif au prétérit — Transformation Actif $\rightarrow$ Passif]
\begin{enumerate}
\item \textbf{Identifie le COD (Accusatif) de la phrase active} : il devient le \textbf{sujet (Nominatif)} de la phrase passive.
\item \textbf{Conjugue \all{werden} au Prétérit} selon ce nouveau sujet (singulier $\rightarrow$ \all{wurde} ; pluriel $\rightarrow$ \all{wurden}).
\item \textbf{Place le Participe II du verbe principal en toute fin de phrase}.
\item \textbf{L'ancien sujet (agent) :} soit il disparaît (si \all{man} ou inconnu), soit il devient \textbf{\all{von} + Datif} (\all{von der Firma, von den Arbeitern}).
\item \textbf{Compléments circonstants (temps, lieu, manière) restent inchangés} et se placent entre l'auxiliaire et le participe.
\end{enumerate}
\end{methode}

\begin{popfigure}
\begin{tikzpicture}[
  node distance=1.2cm and 2cm,
  box/.style={draw=popblue, thick, rounded corners, fill=popblueL, minimum height=1cm, align=center},
  arrow/.style={-Stealth, thick, poporange}
]
\node[box, align=center] (actif){Actif :\\ \all{Die Firma (Subj) schloss (Verb) die Mine (Obj).}};
\node[box, right=of actif, align=center] (passif){Passif :\\ \all{Die Mine (Subj) wurde (werden) ... geschlossen (Part. II).}};
\draw[arrow] (actif) -- node[above, font=\small, text=popdark, align=center]{1. Obj $\rightarrow$ Subj\\2. Verb $\rightarrow$ werden + Part. II\\3. Subj $\rightarrow$ von + Dat} (passif);
\end{tikzpicture}
\legende{Schéma de transformation Actif $\rightarrow$ Passif (Präteritum)}
\end{popfigure}

\begin{exoresolu}[Transformation Actif $\rightarrow$ Passif Prétérit]
\textbf{Énoncé :} Setzen Sie ins Passiv (Präteritum) :
\begin{enumerate}
\item \all{Die Arbeiter förderten das Erz aus der Mine.}
\item \all{Die Firma schloss die Mine im Jahr 2010.}
\item \all{Man exportierte das Gold nach Europa.}
\item \all{Die Regierung versprach neue Investitionen.}
\end{enumerate}
\tcblower
\textbf{Raisonnement détaillé :}
\begin{enumerate}
\item COD \all{das Erz} (sg neutre) $\rightarrow$ Subj \all{Das Erz}. \all{werden} Prét sg $\rightarrow$ \all{wurde}. Part. II \all{gefördert} (faible). Agent \all{von den Arbeitern} (Dat pl). \all{Das Erz wurde von den Arbeitern aus der Mine gefördert.}
\item COD \all{die Mine} (sg fém) $\rightarrow$ Subj \all{Die Mine}. \all{wurde}. Part. II \all{geschlossen} (fort : \all{schließen}). Agent \all{von der Firma} (Dat sg fém). \all{Die Mine wurde im Jahr 2010 von der Firma geschlossen.}
\item Sujet \all{Man} (indéfini) $\rightarrow$ disparaît au passif. COD \all{das Gold} (sg neutre) $\rightarrow$ Subj. \all{wurde}. Part. II \all{exportiert}. Pas d'agent. \all{Das Gold wurde nach Europa exportiert.}
\item COD \all{neue Investitionen} (pl) $\rightarrow$ Subj \all{Neue Investitionen}. \all{wurden} (pl). Part. II \all{versprochen}. Agent \all{von der Regierung}. \all{Neue Investitionen wurden von der Regierung versprochen.}
\end{enumerate}
\textbf{Solutions :}
\begin{enumerate}
\item \all{Das Erz wurde von den Arbeitern aus der Mine gefördert.}
\item \all{Die Mine wurde im Jahr 2010 von der Firma geschlossen.}
\item \all{Das Gold wurde nach Europa exportiert.}
\item \all{Neue Investitionen wurden von der Regierung versprochen.}
\end{enumerate}
\end{exoresolu}

\courssub{Grammaire : Relatives avec préposition}

\begin{propriete}[Pronoms relatifs avec préposition — Tableau de déclinaison]
\emph{Règle :} Le pronom relatif = article défini décliné selon le \textbf{cas exigé par la préposition}. Le genre et le nombre viennent de l'antécédent.
\begin{center}
\begin{tabularx}{\linewidth}{|X|X|X|X|X|}
\hline
\rowcolor{popblueL} \textbf{Préposition + Cas} & \textbf{Masculin} & \textbf{Féminin} & \textbf{Neutre} & \textbf{Pluriel} \\
\hline
\all{mit} + Dativ & \all{mit dem} & \all{mit der} & \all{mit dem} & \all{mit denen} \\
\all{von} + Dativ & \all{von dem} & \all{von der} & \all{von dem} & \all{von denen} \\
\all{aus} + Dativ & \all{aus dem} & \all{aus der} & \all{aus dem} & \all{aus denen} \\
\all{in} + Dativ (lieu) & \all{in dem / im} & \all{in der} & \all{in dem / im} & \all{in denen} \\
\all{für} + Akkusativ & \all{für den} & \all{für die} & \all{für das} & \all{für die} \\
\all{auf} + Akkusativ (verbe) & \all{auf den} & \all{auf die} & \all{auf das} & \all{auf die} \\
\hline
\end{tabularx}
\end{center}
\emph{Alternative « wo(r)- » (choses seulement) :} \all{womit, wovon, woraus, worin, wofür, worauf}. Ex. \all{Das Gold, woraus Schmuck gemacht wird...} (standard : \all{aus dem}).
\end{propriete}

\begin{methode}[Construire une relative avec préposition — 5 étapes]
\begin{enumerate}
\item \textbf{Antécédent :} Identifie le nom de référence (genre, nombre) : \all{der Arbeiter} (m sg), \all{die Mine} (f sg), \all{die Rohstoffe} (pl).
\item \textbf{Préposition :} Repère le verbe ou l'expression qui la impose : \all{sprechen mit, warten auf, bestehen aus, arbeiten in, leben von}.
\item \textbf{Cas :} Détermine le cas régi par la préposition (Dativ : \all{mit, von, aus, in} lieu ; Akkusativ : \all{für, auf} verbe).
\item \textbf{Pronom :} Décline l'article défini (ou \all{wo(r)-} pour choses) : \all{mit dem, aus der, von denen, für das}.
\item \textbf{Verbe final :} Place le verbe conjugué à la fin de la subordonnée relative.
\end{enumerate}
\end{methode}

\begin{exoresolu}[Relatives avec préposition : compléter et combiner]
\textbf{Énoncé :} a) Ergänzen Sie das richtige Relativpronomen :
\begin{enumerate}
\item \all{Die Mine, in \_\_\_ die Arbeiter arbeiten, ist sehr alt.}
\item \all{Der Rohstoff, aus \_\_\_ Schmuck hergestellt wird, ist teuer.}
\item \all{Die Arbeiter, mit \_\_\_ der Journalist gesprochen hat, sind müde.}
\item \all{Die Bedingungen, unter \_\_\_ sie leiden, sind unmenschlich.}
\end{enumerate}
b) Verbinden Sie zwei phrases en une seule avec une relative :
\begin{enumerate}
\item \all{Der Bergbau ist wichtig. Viele Familien leben von dem Bergbau.}
\item \all{Die Firma exportiert das Gold. Die Firma macht große Gewinne damit.}
\end{enumerate}
\tcblower
\textbf{Raisonnement :}
a) 1. \all{die Mine} f sg, \all{in} + Dativ (lieu) $\rightarrow$ \all{der}. 2. \all{der Rohstoff} m sg, \all{aus} + Dativ $\rightarrow$ \all{dem}. 3. \all{die Arbeiter} pl, \all{mit} + Dativ $\rightarrow$ \all{denen} (personnes, pluriel). 4. \all{die Bedingungen} pl, \all{unter} + Dativ $\rightarrow$ \all{denen}.
b) 1. Antécédent \all{der Bergbau} m sg, \all{leben von} + Dativ $\rightarrow$ \all{von dem}. Verbe \all{leben} final. 2. Antécédent \all{die Firma} f sg (ou \all{das Gold} n sg ?). Sens : \all{damit} (avec cela/argent) ou \all{woraus}. Ici \all{davon} (avec cela = l'or/l'argent). \all{Die Firma, die damit große Gewinne macht, exportiert das Gold.} Ou relatif sur \all{Gold} : \all{Das Gold, womit die Firma große Gewinne macht, wird exportiert.}

\textbf{Solutions :}
a) 1. \all{in \textbf{der}} 2. \all{aus \textbf{dem}} 3. \all{mit \textbf{denen}} 4. \all{unter \textbf{denen}}
b) 1. \all{Der Bergbau, von \textbf{dem} viele Familien leben, ist wichtig.}
2. \all{Die Firma, die \textbf{damit} große Gewinne macht, exportiert das Gold.} (ou \all{Das Gold, \textbf{woraus} die Firma große Gewinne macht, wird exportiert.})
\end{exoresolu}

\courssub{Commentaire guidé \& Production écrite}

\begin{methode}[Résumer un texte technique en 4--5 phrases — Attendus Bac]
\begin{enumerate}
\item \textbf{Identifie les 4--5 idées-forces} (Quoi ? Où ? Qui ? Comment ? Conséquence/Problème). Ignore exemples et détails.
\item \textbf{Une phrase = une idée.} Utilise des connecteurs logiques : \all{zuerst / zunächst, dann / danach, außerdem, jedoch / aber, deshalb / daher, schließlich / zum Schluss, obwohl}.
\item \textbf{Varie les structures syntaxiques} (critère d'excellence) : 1 phrase active, 1 passif Prétérit/Présent, 1 relative avec préposition, 1 subordonnée conjonctive (\all{weil, obwohl, damit}).
\item \textbf{Style objectif :} Pas de \all{ich finde / meiner Meinung nach}. Utilise \all{Der Text berichtet, dass... / Es geht darum, dass... / Laut Text...}.
\item \textbf{Relecture checklist :} V2 / Verbe final, majuscules noms, accords article/adj, cas prépositionnels, particules séparables.
\end{enumerate}
\end{methode}

\begin{exoresolu}[Rédaction : résumé guidé du texte support]
\textbf{Énoncé :} Fassen Sie den Text (Textebox 1) in vier bis fünf Sätzen zusammen.
\tcblower
\textbf{Analyse des idées-forces :}
1. Extraction industrielle d'or dans l'Est depuis 10+ ans.
2. Passage méthode manuelle (artisanale) $\rightarrow$ machines (industriel).
3. Conditions de travail dures, bas salaires pour ouvriers locaux.
4. Exportation quasi totale de la matière première (bijoux, électronique).
5. Paradoxe : revenus pour l'État mais pauvreté locale persistante ; promesses gouvernementales.

\textbf{Production rédigée (modèle Bac) :}
\all{Seit über zehn Jahren wird in der Region Ost Kameruns Gold industriell abgebaut. Früher förderten Kleinbergleute das Erz von Hand, heute jedoch nutzen Firmen moderne Maschinen. Die Arbeiter, unter denen viele unter schweren Bedingungen leiden, verdienen zu wenig. Der Rohstoff, aus dem Schmuck und Elektronik hergestellt werden, wird fast komplett exportiert. Obwohl der Staat Steuereinnahmen hat, bleibt die Region arm, doch die Regierung verspricht Besserung.}

\textbf{Traduction :} « Depuis plus de dix ans, de l'or est extrait industriellement dans la région de l'Est. Autrefois, des mineurs artisanaux extrayaient le minerai à la main, mais aujourd'hui des entreprises utilisent des machines modernes. Les ouvriers, parmi lesquels beaucoup souffrent de conditions dures, gagnent trop peu. La matière première, à partir de laquelle on fabrique bijoux et électronique, est presque entièrement exportée. Bien que l'État ait des recettes fiscales, la région reste pauvre, mais le gouvernement promet l'amélioration. »

\textbf{Analyse linguistique du modèle :}
\begin{itemize}
\item Phrase 1 : Passif Présent (\all{wird ... abgebaut}) — fait actuel.
\item Phrase 2 : Contraste \all{früher / heute}, Actif Prétérit (\all{förderten}) / Actif Présent (\all{nutzen}), particule \all{ab} rejetée.
\item Phrase 3 : Relative avec préposition \all{unter denen} (Dativ pluriel), verbe final \all{leiden}.
\item Phrase 4 : Relative \all{aus dem} (Dativ sg neutre) + Passif Présent \all{wird ... hergestellt}, Passif Présent principal \all{wird ... exportiert}.
\item Phrase 5 : Concession \all{Obwohl} (verbe final \all{hat}), principal \all{bleibt}, coordination \all{doch} + subordonnée \all{dass} (omise, infinitif \all{verspricht ... zu investieren}).
\end{itemize}
\end{exoresolu}

\begin{exercice}[Production écrite : Votre tour]
\textbf{Consigne :} Vous êtes journaliste pour \all{„Kamerun Heute“}. Rédigez un court article (80--100 mots) sur le thème : \all{„Chancen und Risiken des Bergbaus in der Region Ost“}. Utilisez : 1 passif prétérit, 1 relative avec préposition, 1 subordonnée \all{obwohl}, le vocabulaire du chapitre. Titre obligatoire.
\end{exercice}

\courssub{Production orale guidée \& Réinvestissement}

\begin{experience}[Jeu de rôle : « Table ronde : L'or de l'Est — qui en profite ? »]
\textbf{Rôles (par groupe de 4) :}
\begin{enumerate}
\item \textbf{Le Ministre des Mines} : défend l'industrialisation, les recettes fiscales, les promesses d'infrastructures.
\item \textbf{Le Représentant d'une ONG environnementale} : dénonce la pollution (mercure, cyanure), la déforestation, l'absence de réhabilitation des sites.
\item \textbf{Un Bergmann / Ouvrier local} : témoigne des conditions (salaire, santé, sécurité), demande une part locale des bénéfices.
\item \textbf{Le Journaliste (modérateur)} : pose les questions, relance, synthétise.
\end{enumerate}
\textbf{Préparation (10 min) :} Chaque élève note 3 arguments clés avec vocabulaire ciblé (\all{der Abbau, die Förderung, die Arbeitsbedingungen, der Umweltschutz, die Steuereinnahmen, die lokale Bevölkerung, profitieren, leiden, versprechen}).
\textbf{Déroulement (10 min) :} Discussion en allemand. Le modérateur ouvre : \all{„Guten Tag. Heute diskutieren wir über den Bergbau in der Region Ost. Herr Minister, was sind die Vorteile?“}
\textbf{Consigne langue :} Obligation d'utiliser au moins 1 passif prétérit (\all{Die Mine wurde eröffnet...}), 1 relative prépositionnelle (\all{Die Arbeiter, mit denen ich sprach...}), 1 \all{obwohl}-Satz.
\end{experience}

\begin{methode}[Présenter un exposé court (3 min) sur un aspect minier]
\begin{enumerate}
\item \textbf{Structure :} Introduction (accroche + sujet) $\rightarrow$ 2 arguments (pour/contre ou fait/conséquence) $\rightarrow$ Conclusion (avis personnel mesuré / perspective).
\item \textbf{Phrases-types :} \all{Mein Thema ist der Goldabbau in der Region Ost. / Zunächst einmal ... / Ein weiterer Punkt ist ... / Dabei muss man beachten, dass ... / Obwohl ... , ... / Zusammenfassend lässt sich sagen, dass ...}
\item \textbf{Appui visuel :} Carte de l'Est, photo de mine, graphique production/exportation. Parlez librement (pas de lecture), regard vers la classe.
\end{enumerate}
\end{methode}

\courssub{Erreurs à éviter — Check-list Bac}

\begin{attention}[Les 6 erreurs qui coûtent des points au Baccalauréat]
\begin{enumerate}
\item \textbf{Accord \all{wurde} / \all{wurden} :} Le Prétérit de \all{werden} s'accorde avec le \textbf{sujet du passif} (ex-COD). \all{Das Erz \textbf{wurde} gefunden} (sg) mais \all{Die Erze \textbf{wurden} gefunden} (pl). \cle{Vérifie le nombre du sujet avant d'écrire l'auxiliaire.}
\item \textbf{Participe II oublié ou mal placé :} Au passif, le participe II est \textbf{toujours en position finale} de la phrase (ou de la proposition). Jamais : \all{*Die Mine wurde geschlossen 2010.} $\rightarrow$ \all{Die Mine wurde 2010 \textbf{geschlossen}.}
\item \textbf{Cas après préposition dans la relative :} \all{mit, von, aus, in, bei, zu, nach, seit} exigent le \textbf{Dativ}. \all{der Arbeiter, mit \textbf{dem} ich sprach} (pas \all{*mit den}). \all{die Mine, in \textbf{der} sie arbeiten} (pas \all{*in die}). Apprends le tableau de la Propriété 5 par cœur.
\item \textbf{Faux amis « Extraction / Mineur » :} \all{die Extraktion} = extraction (chimie, dentaire). Pour les mines : \all{der Abbau} / \all{die Förderung}. \all{der Minderjährige} = mineur (âge). Ouvrier de mine = \all{der Bergmann} / \all{der Bergarbeiter} / \all{der Minenarbeiter}.
\item \textbf{Minuscules aux noms \& Particules séparables :} \textbf{Tout nom} prend une majuscule : \all{das \textbf{E}rz, die \textbf{M}ine, der \textbf{B}ergbau, die \textbf{A}rbeit}. Verbe à particule : \all{abbauen $\rightarrow$ Die Firma baut das Gold \textbf{ab}} (fin de phrase), \all{Das Gold wird \textbf{ab}gebaut} (passif). Jamais \all{*abbaut} / \all{*abgebaut wird}.
\item \textbf{Confusion \all{seit} / \all{seid} :} \all{seit} (depuis, préposition/temps) vs \all{seid} (vous êtes, verbe). \all{Seit zehn Jahren wird Gold abgebaut.} \all{Seid ihr bereit?}
\end{enumerate}
\end{attention}

\newpage

\coursec{Exercices}

\courssub{Je vérifie mes connaissances}

\begin{exercice}[QCM : le passif au prétérit]
Entoure la bonne réponse.
\begin{enumerate}
\item Das Erz \textbf{......} aus der Mine geholt. \\
a) wurde \quad b) wurde geworden \quad c) war \quad d) ist worden
\item Die Arbeiter \textbf{......} gestern gut bezahlt. \\
a) wurde \quad b) wurden \quad c) worden \quad d) wart
\item Im 19. Jahrhundert \textbf{......} hier Gold gefunden. \\
a) wird \quad b) wurde \quad c) ist \quad d) war worden
\item Die Maschinen \textbf{......} von der Firma geliefert. \\
a) wurden \quad b) wurde \quad c) sind \quad d) ward
\end{enumerate}
\end{exercice}

\begin{exercice}[Vrai ou faux ? Justifie ta réponse]
Réponds par \all{richtig} ou \all{falsch} et corrige la phrase si elle est fausse.
\begin{enumerate}
\item Au passif prétérit, on utilise \all{sein} + participe II.
\item \all{Die Mine, in der gearbeitet wird, ist alt.} est une relative avec préposition correcte.
\item Le participe II de \all{finden} est \all{gefunden}.
\item \all{Der Rohstoff} signifie « le produit fini ».
\item Dans une relative avec préposition, la préposition se place avant le pronom relatif.
\end{enumerate}
\end{exercice}

\begin{exercice}[Vocabulaire : le champ minier]
Complète chaque phrase avec le mot qui convient : \all{das Erz, der Bergbau, die Mine, das Gold, der Rohstoff, der Arbeiter}.
\begin{enumerate}
\item In der ...... arbeiten viele Menschen unter Tage.
\item ...... ist ein wertvolles Metall, das in Kamerun gefunden wird.
\item Der ...... trägt einen Helm und eine Lampe.
\item Eisen ist ein wichtiger ...... für die Industrie.
\item Aus dem ...... wird Metall gewonnen.
\item Der ...... ist für die Wirtschaft der Region sehr wichtig.
\end{enumerate}
\end{exercice}

\begin{exercice}[Conjugaison : prétérit passif]
Complète le tableau avec la forme correcte de \all{werden} au prétérit + participe II du verbe entre parenthèses.
\begin{center}
\begin{tabular}{|l|l|}\hline
\rowcolor{popblueL} Sujet & Phrase à compléter \\ \hline
das Erz & ...... (fördern) \\ \hline
die Minen & ...... (schließen) \\ \hline
der Arbeiter & ...... (verletzen) \\ \hline
wir & ...... (informieren) \\ \hline
ihr & ...... (fragen) \\ \hline
das Gold & ...... (exportieren) \\ \hline
\end{tabular}
\end{center}
\end{exercice}

\courssub{Je m'entraîne}

\begin{exercice}[Texte à trous : le passif au prétérit]
Complète le texte avec le passif prétérit des verbes entre parenthèses.

\begin{textebox}[Der Bergbau in der Region]
Im letzten Jahrhundert ...... (entdecken) in dieser Region große Goldvorkommen. Zuerst ...... (bauen) kleine Minen, und viele Arbeiter ...... (einstellen). Das Erz ...... (transportieren) mit einfachen Wagen zur Stadt. Später ...... (kaufen) moderne Maschinen, und die Produktion ...... (steigern). Leider ...... (verschmutzen) auch die Flüsse der Region. Deshalb ...... (gründen) im Jahr 2005 eine Kommission zum Schutz der Umwelt.
\end{textebox}
\end{exercice}

\begin{exercice}[Transformation : de l'actif au passif prétérit]
Mets les phrases suivantes au passif prétérit. Attention à l'ordre des mots !
\begin{enumerate}
\item Die Arbeiter förderten das Erz aus der Mine.
\item Die Firma hat neue Maschinen gekauft.
\item Man entdeckte Gold in den Bergen.
\item Die Regierung schloss die alte Mine.
\item Die Ingenieure kontrollierten die Sicherheit.
\item Man transportierte die Rohstoffe in den Hafen.
\item Die Firma bezahlte die Arbeiter gut.
\item Man hat die Umwelt verschmutzt.
\item Die Experten untersuchten das Wasser des Flusses.
\item Man baute eine neue Straße zur Mine.
\end{enumerate}
\end{exercice}

\begin{exercice}[Relatives avec préposition : appariement et combinaison]
\textbf{A.} Relie chaque début de phrase à la relative qui convient.
\begin{center}
\begin{tabular}{|l|l|}\hline
\rowcolor{popblueL} Début & Relative \\ \hline
1. Die Mine, & a) über die viele Bücher geschrieben wurden, \\ \hline
2. Der Arbeiter, & b) in der 500 Menschen arbeiten, \\ \hline
3. Das Gold, & c) mit dem ich gesprochen habe, \\ \hline
4. Der Bergbau, & d) für das sich viele Länder interessieren, \\ \hline
\end{tabular}
\end{center}

\textbf{B.} Combine les deux phrases avec une relative à préposition.
\begin{enumerate}
\item Die Mine ist sehr alt. In der Mine wird noch gearbeitet.
\item Der Fluss ist verschmutzt. Die Mine liegt an dem Fluss.
\item Die Maschinen sind modern. Mit den Maschinen wird das Erz gefördert.
\item Der Minister sprach gestern. Wir haben auf den Minister gewartet.
\end{enumerate}
\end{exercice}

\begin{exercice}[Thème : traduction français $\rightarrow$ allemand]
Traduis les phrases suivantes en allemand. Utilise le passif prétérit et les relatives avec préposition.
\begin{enumerate}
\item Le minerai a été transporté à l'usine.
\item Les ouvriers ont été bien payés par l'entreprise.
\item La mine dans laquelle on travaille est dangereuse.
\item Beaucoup de matières premières ont été exportées l'année dernière.
\item Le fleuve près duquel se trouve la mine a été pollué.
\end{enumerate}
\end{exercice}

\courssub{Je me prépare au Bac}

\begin{exercice}[Compréhension de l'écrit : texte et questions]
Lis le texte suivant, puis réponds aux questions \textbf{en allemand}, avec des phrases complètes.

\begin{textebox}[Goldbergbau in Kamerun — Chancen und Probleme]
In vielen Regionen Kameruns, besonders im Osten des Landes, wird seit Jahrzehnten Gold gesucht. Zuerst wurde das Gold von Hand aus den Flüssen gewaschen; später wurden kleine Minen gebaut, in denen Hunderte von Arbeitern beschäftigt wurden. Das Erz, aus dem das Gold gewonnen wird, wurde oft mit einfachen Mitteln gefördert.

Für viele Familien ist der Bergbau eine wichtige Einnahmequelle: Mit dem Geld, das verdient wird, werden Schulgebühren bezahlt und Häuser gebaut. Doch es gibt auch dunkle Seiten. Flüsse, in denen früher gefischt wurde, sind heute verschmutzt. Wälder wurden abgeholzt, und Unfälle, bei denen Arbeiter verletzt wurden, sind keine Seltenheit.

Experten fordern deshalb strengere Kontrollen. „Der Bergbau darf die Umwelt nicht zerstören“, sagt ein Umweltschützer aus Yaoundé. „Die Minen, in denen gearbeitet wird, müssen regelmäßig kontrolliert werden.“ Die Regierung hat versprochen, dass neue Gesetze verabschiedet werden, damit Mensch und Natur besser geschützt werden.
\end{textebox}

\textbf{Fragen zum Text :}
\begin{enumerate}
\item Wo wird in Kamerun Gold gesucht ?
\item Wie wurde das Gold früher gewonnen ?
\item Warum ist der Bergbau für viele Familien wichtig ? Nenne zwei Gründe.
\item Welche Probleme für die Umwelt werden im Text genannt ?
\item Was fordern die Experten ?
\end{enumerate}

\textbf{Fragen zur Sprache :}
\begin{enumerate}
\item Relève dans le texte trois phrases au passif prétérit et souligne l'auxiliaire.
\item Relève deux relatives avec préposition et indique le nom auquel elles se rapportent.
\item Trouve dans le texte : le nom signifiant « source de revenus » et le verbe signifiant « abattre (des arbres) ».
\end{enumerate}
\end{exercice}

\begin{exercice}[Thème et version]
\textbf{A. Thème.} Traduis le texte suivant en allemand (environ 60 mots). Utilise au moins deux verbes au passif prétérit et une relative avec préposition.

« Autrefois, l'or était cherché à la main dans les rivières. Puis des mines ont été construites, dans lesquelles beaucoup d'ouvriers ont été employés. Le minerai était transporté vers la ville, où il était vendu. Malheureusement, l'eau des rivières a été polluée, et les habitants se sont plaints. Aujourd'hui, les mines doivent être contrôlées régulièrement. »

\textbf{B. Version.} Traduis le texte suivant en français.

\begin{textebox}[Ein Bericht aus der Mine]
Gestern wurde die alte Mine bei Batouri von einer Kommission besucht. Die Maschinen, mit denen das Erz gefördert wird, wurden genau kontrolliert. Zwei Arbeiter, die letztes Jahr verletzt worden waren, erzählten von ihrem Alltag. Am Ende wurde ein Bericht geschrieben, in dem bessere Sicherheitsmaßnahmen gefordert werden.
\end{textebox}
\end{exercice}

\begin{exercice}[Rédaction : lettre formelle]
Tu es président(e) du club environnement de ton lycée. Écris une lettre formelle (environ \textbf{150 mots}) au ministre de l'Environnement pour demander un contrôle des mines de ta région.

Ta lettre doit contenir :
\begin{itemize}
\item l'en-tête (expéditeur, destinataire, lieu et date, objet : \all{Betreff}) ;
\item la formule de politesse d'ouverture : \all{Sehr geehrter Herr Minister,} ;
\item une présentation du problème (pollution, accidents, forêts détruites) avec au moins \textbf{deux phrases au passif prétérit} ;
\item \textbf{une relative avec préposition} (par exemple : \all{die Mine, in der ...} ; \all{der Fluss, an dem ...}) ;
\item deux demandes concrètes (contrôles réguliers, protection des rivières, sécurité des ouvriers) ;
\item la formule de politesse finale : \all{Mit freundlichen Grüßen} et ta signature.
\end{itemize}
\end{exercice}

\newpage

\coursec{Corrigés des exercices}

\begin{corrige}[Exercice 1 — QCM : le passif au prétérit]
\begin{enumerate}
\item \textbf{a) wurde} — \all{Das Erz wurde aus der Mine geholt.} Sujet singulier (\all{das Erz}) $\rightarrow$ \all{wurde} + participe II \all{geholt}. « Le minerai fut extrait de la mine. »
\item \textbf{b) wurden} — \all{Die Arbeiter wurden gestern gut bezahlt.} Sujet pluriel $\rightarrow$ \all{wurden}. « Les ouvriers furent bien payés hier. »
\item \textbf{b) wurde} — \all{Im 19. Jahrhundert wurde hier Gold gefunden.} \all{Gold} est singulier ; le passif prétérit exige \all{wurde} + \all{gefunden}. « Au XIX\textsuperscript{e} siècle, de l'or fut trouvé ici. »
\item \textbf{a) wurden} — \all{Die Maschinen wurden von der Firma geliefert.} Pluriel $\rightarrow$ \all{wurden}. « Les machines furent livrées par l'entreprise. »
\end{enumerate}
\textbf{Points de vigilance :} l'auxiliaire du passif prétérit est \all{werden} au prétérit (\all{wurde/wurden}), jamais \all{sein}. \all{ward} est une forme littéraire vieillie, à éviter à l'écrit scolaire. Le participe II se place en fin de phrase.
\end{corrige}

\begin{corrige}[Exercice 2 — Vrai ou faux ? Justifie ta réponse]
\begin{enumerate}
\item \all{falsch} — Au passif prétérit, on utilise \all{werden} (au prétérit : \all{wurde/wurden}) + participe II : \all{Das Erz wurde gefördert.} C'est au \emph{Perfekt passif} qu'on emploie \all{sein} : \all{Das Erz ist gefördert worden.}
\item \all{richtig} — \all{Die Mine, in der gearbeitet wird, ist alt.} est correcte : \all{die Mine} est féminin, la préposition \all{in} + datif donne \all{in der}. « La mine dans laquelle on travaille est vieille. »
\item \all{richtig} — \all{finden — fand — hat gefunden} : le participe II est bien \all{gefunden}.
\item \all{falsch} — \all{der Rohstoff} signifie « la matière première » (et non « le produit fini », qui se dit \all{das Fertigprodukt}). Attention au faux ami : \all{Roh-} = « brut, cru ».
\item \all{richtig} — En allemand, la préposition précède toujours le pronom relatif : \all{die Mine, in der ...} ; \all{der Fluss, an dem ...}. Jamais de préposition rejetée en fin de relative comme en anglais.
\end{enumerate}
\end{corrige}

\begin{corrige}[Exercice 3 — Vocabulaire : le champ minier]
\begin{enumerate}
\item \all{In der \textbf{Mine} arbeiten viele Menschen unter Tage.} « Dans la mine, beaucoup de gens travaillent sous terre. »
\item \all{\textbf{Das Gold} ist ein wertvolles Metall, das in Kamerun gefunden wird.} « L'or est un métal précieux que l'on trouve au Cameroun. »
\item \all{Der \textbf{Arbeiter} trägt einen Helm und eine Lampe.} « L'ouvrier porte un casque et une lampe. »
\item \all{Eisen ist ein wichtiger \textbf{Rohstoff} für die Industrie.} « Le fer est une matière première importante pour l'industrie. »
\item \all{Aus dem \textbf{Erz} wird Metall gewonnen.} « On extrait le métal du minerai. »
\item \all{Der \textbf{Bergbau} ist für die Wirtschaft der Region sehr wichtig.} « L'exploitation minière est très importante pour l'économie de la région. »
\end{enumerate}
\textbf{Astuce :} vérifier l'article demandé par la phrase (\all{In der} $\rightarrow$ féminin ; \all{Aus dem} $\rightarrow$ neutre) : c'est un indice grammatical qui confirme le choix du mot.
\end{corrige}

\begin{corrige}[Exercice 4 — Conjugaison : prétérit passif]
\begin{center}
\begin{tabular}{|l|l|l|}\hline
\rowcolor{popblueL} Sujet & Phrase complète & Traduction \\ \hline
das Erz & Das Erz \textbf{wurde gefördert}. & Le minerai fut extrait. \\ \hline
die Minen & Die Minen \textbf{wurden geschlossen}. & Les mines furent fermées. \\ \hline
der Arbeiter & Der Arbeiter \textbf{wurde verletzt}. & L'ouvrier fut blessé. \\ \hline
wir & Wir \textbf{wurden informiert}. & Nous fûmes informés. \\ \hline
ihr & Ihr \textbf{wurdet gefragt}. & Vous fûtes interrogés. \\ \hline
das Gold & Das Gold \textbf{wurde exportiert}. & L'or fut exporté. \\ \hline
\end{tabular}
\end{center}
\textbf{Points de vigilance :}
\begin{itemize}
\item Accord de \all{werden} avec le sujet : singulier \all{wurde}, pluriel \all{wurden}, 2\textsuperscript{e} pers. pluriel \all{wurdet}.
\item Participes II : \all{schließen $\rightarrow$ geschlossen} (verbe fort) ; \all{verletzen $\rightarrow$ verletzt} ; les verbes en \all{-ieren} ne prennent pas de \all{ge-} : \all{informiert}, \all{exportiert}.
\end{itemize}
\end{corrige}

\begin{corrige}[Exercice 5 — Texte à trous : le passif au prétérit]
\begin{textebox}[Der Bergbau in der Region — Lösung]
Im letzten Jahrhundert \textbf{wurden} in dieser Region große Goldvorkommen \textbf{entdeckt}. Zuerst \textbf{wurden} kleine Minen \textbf{gebaut}, und viele Arbeiter \textbf{wurden eingestellt}. Das Erz \textbf{wurde} mit einfachen Wagen zur Stadt \textbf{transportiert}. Später \textbf{wurden} moderne Maschinen \textbf{gekauft}, und die Produktion \textbf{wurde gesteigert}. Leider \textbf{wurden} auch die Flüsse der Region \textbf{verschmutzt}. Deshalb \textbf{wurde} im Jahr 2005 eine Kommission zum Schutz der Umwelt \textbf{gegründet}.
\end{textebox}
\textbf{Traduction :} « Au siècle dernier, de grands gisements d'or furent découverts dans cette région. D'abord de petites mines furent construites, et de nombreux ouvriers furent embauchés. Le minerai était transporté vers la ville avec de simples charrettes. Plus tard, des machines modernes furent achetées, et la production fut augmentée. Malheureusement, les rivières de la région furent aussi polluées. C'est pourquoi une commission pour la protection de l'environnement fut créée en 2005. »

\textbf{Justifications :}
\begin{itemize}
\item \all{Goldvorkommen} (pluriel) $\rightarrow$ \all{wurden entdeckt} ; \all{das Erz} (singulier) $\rightarrow$ \all{wurde transportiert} ; \all{eine Kommission} (singulier) $\rightarrow$ \all{wurde gegründet}.
\item Verbes à particule séparable : \all{einstellen $\rightarrow$ eingestellt} (\all{-ge-} inséré entre particule et radical).
\item Verbes en \all{-ieren} : \all{transportiert} (sans \all{ge-}).
\end{itemize}
\end{corrige}

\begin{corrige}[Exercice 6 — Transformation : de l'actif au passif prétérit]
\begin{enumerate}
\item \all{Das Erz wurde von den Arbeitern aus der Mine gefördert.} « Le minerai fut extrait de la mine par les ouvriers. »
\item \all{Neue Maschinen wurden von der Firma gekauft.} (Le Perfekt actif \all{hat gekauft} devient au passif prétérit \all{wurden gekauft} ; on peut aussi écrire \all{sind gekauft worden} au Perfekt passif.)
\item \all{Gold wurde in den Bergen entdeckt.}
\item \all{Die alte Mine wurde von der Regierung geschlossen.}
\item \all{Die Sicherheit wurde von den Ingenieuren kontrolliert.}
\item \all{Die Rohstoffe wurden in den Hafen transportiert.}
\item \all{Die Arbeiter wurden von der Firma gut bezahlt.}
\item \all{Die Umwelt wurde verschmutzt.} (ou \all{ist verschmutzt worden} au Perfekt passif)
\item \all{Das Wasser des Flusses wurde von den Experten untersucht.}
\item \all{Eine neue Straße zur Mine wurde gebaut.}
\end{enumerate}
\textbf{Méthode :} 1) le complément d'objet direct de l'actif devient le sujet du passif ; 2) \all{werden} s'accorde avec ce nouveau sujet et se met au prétérit ; 3) le participe II va en fin de phrase ; 4) l'ancien sujet devient un complément d'agent introduit par \all{von} + datif ; \all{man} disparaît simplement.

\textbf{Points de vigilance :} ne pas conserver \all{man} au passif ; attention à l'accord singulier/pluriel (\all{Gold wurde}, \all{die Rohstoffe wurden}) ; ne pas oublier le \all{ge-} des participes (\all{bezahlt}, \all{gebaut}) sauf pour les verbes en \all{-ieren} et les verbes à préfixe inséparable (\all{entdeckt}, \all{verschmutzt}, \all{untersucht}).
\end{corrige}

\begin{corrige}[Exercice 7 — Relatives avec préposition : appariement et combinaison]
\textbf{A. Appariement :}
\begin{itemize}
\item 1 $\rightarrow$ b : \all{Die Mine, \textbf{in der} 500 Menschen arbeiten, ...} (\all{die Mine}, féminin, \all{in} + datif $\rightarrow$ \all{in der}). « La mine, dans laquelle 500 personnes travaillent, ... »
\item 2 $\rightarrow$ c : \all{Der Arbeiter, \textbf{mit dem} ich gesprochen habe, ...} (\all{der Arbeiter}, masculin, \all{mit} + datif $\rightarrow$ \all{mit dem}). « L'ouvrier avec lequel j'ai parlé, ... »
\item 3 $\rightarrow$ d : \all{Das Gold, \textbf{für das} sich viele Länder interessieren, ...} (\all{das Gold}, neutre, \all{für} + accusatif $\rightarrow$ \all{für das}). « L'or, auquel beaucoup de pays s'intéressent, ... »
\item 4 $\rightarrow$ a : \all{Der Bergbau, \textbf{über die} viele Bücher geschrieben wurden, ...} — attention : \all{über die} renvoie ici au contenu évoqué ; formulation acceptée, la préposition \all{über} + accusatif signifie « au sujet de ». « L'exploitation minière, au sujet de laquelle beaucoup de livres ont été écrits, ... »
\end{itemize}

\textbf{B. Combinaison :}
\begin{enumerate}
\item \all{Die Mine, in der noch gearbeitet wird, ist sehr alt.} « La mine dans laquelle on travaille encore est très vieille. »
\item \all{Der Fluss, an dem die Mine liegt, ist verschmutzt.} « La rivière au bord de laquelle se trouve la mine est polluée. » (\all{an} + datif $\rightarrow$ \all{an dem})
\item \all{Die Maschinen, mit denen das Erz gefördert wird, sind modern.} « Les machines avec lesquelles le minerai est extrait sont modernes. » (pluriel $\rightarrow$ \all{mit denen})
\item \all{Der Minister, auf den wir gewartet haben, sprach gestern.} « Le ministre que nous attendions a parlé hier. » (\all{warten auf} + accusatif $\rightarrow$ \all{auf den})
\end{enumerate}
\textbf{Points de vigilance :} le cas du pronom relatif dépend de la préposition (\all{mit} + datif, \all{auf} + accusatif avec \all{warten}) ; le genre et le nombre dépendent du nom de référence ; le verbe conjugué de la relative se place en fin de proposition.
\end{corrige}

\begin{corrige}[Exercice 8 — Thème : traduction français $\rightarrow$ allemand]
\begin{enumerate}
\item \all{Das Erz wurde zur Fabrik transportiert.} (ou \all{in die Fabrik})
\item \all{Die Arbeiter wurden von der Firma gut bezahlt.}
\item \all{Die Mine, in der gearbeitet wird, ist gefährlich.} — tournure impersonnelle idiomatique ; on peut aussi écrire \all{Die Mine, in der man arbeitet, ist gefährlich.}
\item \all{Viele Rohstoffe wurden im letzten Jahr exportiert.}
\item \all{Der Fluss, in dessen Nähe die Mine liegt, wurde verschmutzt.} — variante plus simple : \all{Der Fluss, bei dem die Mine liegt, wurde verschmutzt.}
\end{enumerate}
\textbf{Barème indicatif :} 2 points par phrase (1 point pour la structure grammaticale : passif ou relative correcte ; 1 point pour le vocabulaire et l'orthographe). Total : 10 points.

\textbf{Points de vigilance :}
\begin{itemize}
\item « L'année dernière » = \all{im letzten Jahr} (et non \all{das letzte Jahr} seul).
\item « L'entreprise » = \all{die Firma} (ou \all{das Unternehmen}) ; ne pas traduire par \all{die Entreprise}, mot inexistant.
\item Majuscules aux noms : \all{Erz, Arbeiter, Mine, Rohstoffe, Fluss}.
\item « Dangereuse » = \all{gefährlich} (avec Umlaut).
\end{itemize}
\end{corrige}

\begin{corrige}[Exercice 9 — Compréhension de l'écrit : texte et questions]
\textbf{Fragen zum Text — réponses modèles :}
\begin{enumerate}
\item \all{In Kamerun wird Gold in vielen Regionen gesucht, besonders im Osten des Landes.} « On cherche de l'or dans de nombreuses régions du Cameroun, surtout dans l'Est du pays. »
\item \all{Früher wurde das Gold von Hand aus den Flüssen gewaschen.} « Autrefois, l'or était lavé à la main dans les rivières. »
\item \all{Der Bergbau ist wichtig, weil mit dem verdienten Geld Schulgebühren bezahlt und Häuser gebaut werden.} « L'exploitation minière est importante parce qu'avec l'argent gagné, on paie les frais de scolarité et on construit des maisons. »
\item \all{Die Flüsse sind verschmutzt, die Wälder wurden abgeholzt, und es gibt Unfälle, bei denen Arbeiter verletzt wurden.} « Les rivières sont polluées, les forêts ont été abattues et il y a des accidents dans lesquels des ouvriers ont été blessés. »
\item \all{Die Experten fordern strengere Kontrollen: Die Minen müssen regelmäßig kontrolliert werden, und neue Gesetze sollen verabschiedet werden.} « Les experts exigent des contrôles plus stricts : les mines doivent être contrôlées régulièrement et de nouvelles lois doivent être adoptées. »
\end{enumerate}

\textbf{Fragen zur Sprache :}
\begin{enumerate}
\item Trois phrases au passif prétérit (auxiliaire souligné) : \all{Zuerst \textbf{wurde} das Gold von Hand aus den Flüssen gewaschen} ; \all{später \textbf{wurden} kleine Minen gebaut} ; \all{Wälder \textbf{wurden} abgeholzt}. (Autres possibilités : \all{wurden Hunderte von Arbeitern beschäftigt}, \all{wurde oft mit einfachen Mitteln gefördert}.)
\item Deux relatives avec préposition : \all{kleine Minen, \textbf{in denen} Hunderte von Arbeitern beschäftigt wurden} (se rapporte à \all{Minen}) ; \all{Unfälle, \textbf{bei denen} Arbeiter verletzt wurden} (se rapporte à \all{Unfälle}). (Aussi : \all{Flüsse, in denen früher gefischt wurde} ; \all{das Erz, aus dem das Gold gewonnen wird}.)
\item « Source de revenus » = \all{die Einnahmequelle} ; « abattre (des arbres) » = \all{abholzen}.
\end{enumerate}

\textbf{Barème indicatif (20 points) :} questions de compréhension : 5 $\times$ 2 points (contenu 1 point, correction linguistique 1 point) ; questions de langue : 3 + 4 + 3 points. Exiger des phrases complètes en allemand : une réponse en un seul mot ne vaut que la moitié des points.
\end{corrige}

\begin{corrige}[Exercice 10 — Thème et version]
\textbf{A. Thème — proposition de traduction :}

\begin{textebox}[Musterlösung — Thème]
Früher wurde das Gold von Hand in den Flüssen gesucht. Dann wurden Minen gebaut, in denen viele Arbeiter eingestellt wurden. Das Erz wurde in die Stadt transportiert, wo es verkauft wurde. Leider wurde das Wasser der Flüsse verschmutzt, und die Einwohner beschwerten sich. Heute müssen die Minen regelmäßig kontrolliert werden.
\end{textebox}

\textbf{Points de vigilance :} « autrefois » = \all{früher} ; « se sont plaints » = \all{beschwerten sich} (verbe réfléchi au prétérit) ; « doivent être contrôlées » = passif avec modal : \all{müssen kontrolliert werden} ; la relative \all{in denen} renvoie au pluriel \all{Minen}.

\textbf{Barème indicatif :} 10 points — passif prétérit correct (2 $\times$ 2 points), relative avec préposition (2 points), vocabulaire (2 points), orthographe et majuscules (2 points).

\textbf{B. Version — proposition de traduction :}

« Hier, l'ancienne mine près de Batouri a été visitée par une commission. Les machines avec lesquelles le minerai est extrait ont été contrôlées minutieusement. Deux ouvriers, qui avaient été blessés l'année dernière, ont raconté leur quotidien. À la fin, un rapport a été rédigé, dans lequel de meilleures mesures de sécurité sont demandées. »

\textbf{Points de vigilance :} \all{wurde besucht} = « a été visitée » (passif prétérit rendu par un passé composé passif en français) ; \all{verletzt worden waren} = plus-que-parfait passif, « avaient été blessés » ; \all{der Alltag} = « le quotidien, la vie quotidienne » ; \all{die Sicherheitsmaßnahmen} = « les mesures de sécurité ».

\textbf{Barème indicatif :} 10 points — fidélité du sens (5 points), rendu correct des formes passives (3 points), qualité du français (2 points).
\end{corrige}

\begin{corrige}[Exercice 11 — Rédaction : lettre formelle]
\textbf{Proposition de rédaction modèle :}

\begin{textebox}[Musterlösung — Brief an den Minister]
Umweltclub des Gymnasiums Batouri\\
Paul Ngoa, Präsident\\
B.P. 45 Batouri

An den Herrn Minister für Umwelt\\
Yaoundé

Batouri, den 15. März 2025

\textbf{Betreff: Kontrolle der Minen in unserer Region}

Sehr geehrter Herr Minister,

ich schreibe Ihnen im Namen des Umweltclubs unseres Gymnasiums. In unserer Region gibt es viele Goldminen, und leider sehen wir jeden Tag die Folgen des Bergbaus.

Letztes Jahr wurden mehrere Flüsse verschmutzt, und große Waldflächen wurden abgeholzt. Auch Unfälle sind häufig: Zwei Arbeiter wurden im letzten Monat verletzt. Der Fluss, an dem unser Dorf liegt, ist heute so schmutzig, dass man darin nicht mehr fischen kann. Die Mine, in der viele junge Menschen arbeiten, wird fast nie kontrolliert.

Deshalb bitten wir Sie um zwei Dinge: Erstens sollten die Minen unserer Region regelmäßig kontrolliert werden. Zweitens bitten wir um bessere Sicherheitsmaßnahmen für die Arbeiter und um den Schutz der Flüsse.

Wir hoffen auf Ihre Hilfe und danken Ihnen im Voraus.

Mit freundlichen Grüßen

Paul Ngoa
\end{textebox}

\textbf{Traduction du corps de la lettre :} « Je vous écris au nom du club environnement de notre lycée. Dans notre région, il y a beaucoup de mines d'or, et malheureusement nous voyons chaque jour les conséquences de l'exploitation minière. L'année dernière, plusieurs rivières ont été polluées et de grandes surfaces forestières ont été abattues. Les accidents aussi sont fréquents : deux ouvriers ont été blessés le mois dernier. La rivière au bord de laquelle se trouve notre village est aujourd'hui si sale qu'on ne peut plus y pêcher. La mine dans laquelle travaillent beaucoup de jeunes n'est presque jamais contrôlée. C'est pourquoi nous vous demandons deux choses : premièrement, les mines de notre région devraient être contrôlées régulièrement. Deuxièmement, nous demandons de meilleures mesures de sécurité pour les ouvriers et la protection des rivières. Nous espérons votre aide et vous remercions d'avance. »

\textbf{Barème indicatif (20 points) :}
\begin{itemize}
\item En-tête complet (expéditeur, destinataire, lieu/date, \all{Betreff}) : 4 points ;
\item Formules de politesse (\all{Sehr geehrter Herr Minister,} / \all{Mit freundlichen Grüßen}) : 2 points ;
\item Présentation du problème avec au moins deux passifs prétérit corrects : 4 points ;
\item Au moins une relative avec préposition correcte : 3 points ;
\item Deux demandes concrètes formulées clairement : 3 points ;
\item Correction linguistique globale, longueur (environ 150 mots), présentation : 4 points.
\end{itemize}

\textbf{Points de vigilance :}
\begin{itemize}
\item Après \all{Sehr geehrter Herr Minister,} la phrase commence par une minuscule (\all{ich schreibe ...}).
\item Ne pas confondre \all{Mit freundlichen Grüßen} (sans \all{Ihnen}, sans point d'exclamation) avec une formule française.
\item Date allemande : \all{Batouri, den 15. März 2025} (accusatif avec \all{den}).
\item Vérifier l'accord du passif : \all{Flüsse wurden verschmutzt} (pluriel), \all{der Fluss wurde verschmutzt} (singulier).
\item Interdiction de calquer « Je vous écris pour vous demander » mot à mot : préférer \all{ich schreibe Ihnen, um ...} ou \all{wir bitten Sie um ...}.
\end{itemize}
\end{corrige}

\newpage

\coursec{Fiche bilan}

\begin{popfigure}
\begin{center}\tcbox[colback=popdark,colframe=popdark,coltext=white,arc=9pt,boxsep=4pt,left=6pt,right=6pt]{\bfseries\small AL12 L'extraction des minerais}\end{center}
\vspace{-2pt}
\begin{tcbraster}[raster columns=3,raster equal height=rows,raster column skip=6pt,raster row skip=6pt,enhanced,blankest]
\begin{tcolorbox}[enhanced,colback=white,colframe=popblue,boxrule=0.9pt,arc=3pt,title={Thème et texte},fonttitle=\bfseries\scriptsize,coltitle=white,colbacktitle=popblue,left=4pt,right=4pt,top=2pt,bottom=2pt,boxsep=1pt,toptitle=1.5pt,bottomtitle=1.5pt,halign=left]
\begin{itemize}[nosep,leftmargin=*,itemsep=1pt,font=\scriptsize]
\item Bergbau in Kamerun und Deutschland
\item Umwelt und Gesundheit
\end{itemize}
\end{tcolorbox}
\begin{tcolorbox}[enhanced,colback=white,colframe=popgreen,boxrule=0.9pt,arc=3pt,title={Lexique},fonttitle=\bfseries\scriptsize,coltitle=white,colbacktitle=popgreen,left=4pt,right=4pt,top=2pt,bottom=2pt,boxsep=1pt,toptitle=1.5pt,bottomtitle=1.5pt,halign=left]
\begin{itemize}[nosep,leftmargin=*,itemsep=1pt,font=\scriptsize]
\item das Erz, der Rohstoff, das Gold
\item der Bergbau, die Mine, der Arbeiter
\end{itemize}
\end{tcolorbox}
\begin{tcolorbox}[enhanced,colback=white,colframe=poporange,boxrule=0.9pt,arc=3pt,title={Passiv Präteritum},fonttitle=\bfseries\scriptsize,coltitle=white,colbacktitle=poporange,left=4pt,right=4pt,top=2pt,bottom=2pt,boxsep=1pt,toptitle=1.5pt,bottomtitle=1.5pt,halign=left]
\begin{itemize}[nosep,leftmargin=*,itemsep=1pt,font=\scriptsize]
\item wurde + Partizip II
\item von + Dativ durch + Akkusativ
\end{itemize}
\end{tcolorbox}
\begin{tcolorbox}[enhanced,colback=white,colframe=poppurple,boxrule=0.9pt,arc=3pt,title={Relatives avec préposition},fonttitle=\bfseries\scriptsize,coltitle=white,colbacktitle=poppurple,left=4pt,right=4pt,top=2pt,bottom=2pt,boxsep=1pt,toptitle=1.5pt,bottomtitle=1.5pt,halign=left]
\begin{itemize}[nosep,leftmargin=*,itemsep=1pt,font=\scriptsize]
\item mit dem, aus denen, in dem
\item worüber, wovon (choses)
\end{itemize}
\end{tcolorbox}
\begin{tcolorbox}[enhanced,colback=white,colframe=poppink,boxrule=0.9pt,arc=3pt,title={Méthodes},fonttitle=\bfseries\scriptsize,coltitle=white,colbacktitle=poppink,left=4pt,right=4pt,top=2pt,bottom=2pt,boxsep=1pt,toptitle=1.5pt,bottomtitle=1.5pt,halign=left]
\begin{itemize}[nosep,leftmargin=*,itemsep=1pt,font=\scriptsize]
\item Lire global puis détaillé
\item Résumé en 3-4 phrases
\end{itemize}
\end{tcolorbox}
\begin{tcolorbox}[enhanced,colback=white,colframe=popgold,boxrule=0.9pt,arc=3pt,title={Production},fonttitle=\bfseries\scriptsize,coltitle=white,colbacktitle=popgold,left=4pt,right=4pt,top=2pt,bottom=2pt,boxsep=1pt,toptitle=1.5pt,bottomtitle=1.5pt,halign=left]
\begin{itemize}[nosep,leftmargin=*,itemsep=1pt,font=\scriptsize]
\item Réponses aux questions
\item Débat : mines et nature
\end{itemize}
\end{tcolorbox}
\begin{tcolorbox}[enhanced,colback=white,colframe=popmuted,boxrule=0.9pt,arc=3pt,title={Pièges},fonttitle=\bfseries\scriptsize,coltitle=white,colbacktitle=popmuted,left=4pt,right=4pt,top=2pt,bottom=2pt,boxsep=1pt,toptitle=1.5pt,bottomtitle=1.5pt,halign=left]
\begin{itemize}[nosep,leftmargin=*,itemsep=1pt,font=\scriptsize]
\item Passif $\neq$ Perfekt
\item Verbe en fin de relative
\end{itemize}
\end{tcolorbox}
\end{tcbraster}

\legende{Carte mentale de la leçon AL12 : l'extraction des minerais}
\end{popfigure}

\begin{bilan}
\textbf{1. Lexique clé à connaître par cœur}
\begin{itemize}
  \item \all{das Erz, -e} : le minerai ; \all{der Rohstoff, -e} : la matière première ; \all{das Gold} : l'or ; \all{die Kohle, -n} : le charbon.
  \item \all{der Bergbau} : l'exploitation minière ; \all{die Mine, -n} : la mine ; \all{das Bergwerk, -e} : la mine (de charbon) ; \all{der Schacht, \"-e} : le puits de mine.
  \item \all{der Arbeiter, - / die Arbeiterin, -nen} : l'ouvrier / l'ouvrière ; \all{der Bergmann, die Bergleute} : le mineur.
  \item Verbes : \all{fördern} (extraire), \all{abbauen} (exploiter), \all{graben, gräbt, grub, gegraben} (creuser), \all{entdecken} (découvrir), \all{verschmutzen} (polluer), \all{schützen} (protéger).
  \item Expressions : \all{in der Mine arbeiten} (travailler à la mine) ; \all{Rohstoffe exportieren} (exporter des matières premières) ; \all{die Umwelt zerstören} (détruire l'environnement).
\end{itemize}

\textbf{2. Grammaire à connaître par cœur}
\begin{center}
\begin{tabularx}{\linewidth}{|X|X|X|}
\hline
\rowcolor{popblueL} Point & Règle & Exemple \\
\hline
Passiv Präteritum & \cle{wurde(n) + Partizip II} ; auxiliaire \all{werden} conjugué au prétérit & \all{Das Erz wurde gefördert.} « Le minerai fut extrait. » \\
\hline
Agent au passif & \all{von + Dativ} (personne), \all{durch + Akkusativ} (moyen/cause) & \all{Die Mine wurde von Arbeitern gebaut.} \\
\hline
Relative + préposition & préposition + pronom relatif (cas exigé par la préposition) ; verbe en fin & \all{die Mine, in der man arbeitet} ; \all{das Gold, nach dem man sucht} \\
\hline
Relative pour les choses & \all{wo(r)- + préposition} possible : \all{wovon, worüber, womit} & \all{das Problem, worüber man spricht} \\
\hline
\end{tabularx}
\end{center}

\textbf{3. Conjugaison : prétérit passif de \all{fördern} (extraire)}
\begin{center}
\begin{tabular}{|l|l|l|}
\hline
\rowcolor{popgreenL} Personne & Forme & Traduction \\
\hline
ich & wurde gefördert & je fus extrait \\
\hline
du & wurdest gefördert & tu fus extrait \\
\hline
er/sie/es & wurde gefördert & il/elle fut extrait(e) \\
\hline
wir & wurden gefördert & nous fûmes extraits \\
\hline
ihr & wurdet gefördert & vous fûtes extraits \\
\hline
sie/Sie & wurden gefördert & ils furent extraits \\
\hline
\end{tabular}
\end{center}

\textbf{4. Méthodes express}
\begin{itemize}
  \item Compréhension : lire une première fois pour l'idée globale (qui ? où ? quoi ?), puis repérer les mots-clés du champ lexical minier avant de répondre aux questions.
  \item Traduire un passif français : « a été extrait » $\rightarrow$ \all{wurde gefördert} ; attention, le passif allemand se forme avec \all{werden}, jamais avec \all{sein} (sauf passif d'état : \all{die Mine ist geschlossen}).
  \item Relative avec préposition : identifier la préposition du verbe (\all{arbeiten in}, \all{warten auf}, \all{sprechen über}), mettre le relatif au cas exigé, renvoyer le verbe conjugué à la fin.
  \item Résumé : 3 à 4 phrases au passif prétérit + une phrase d'opinion avec \all{meiner Meinung nach}.
\end{itemize}
\end{bilan}

\begin{aretenir}[Checklist avant le Bac]
\begin{itemize}
  \item $\square$ Je connais l'article et le pluriel de \all{das Erz, der Rohstoff, die Mine, der Arbeiter, das Bergwerk}.
  \item $\square$ Je sais former le passif prétérit : \all{wurde(n) + Partizip II}.
  \item $\square$ Je sais conjuguer \all{werden} au prétérit : \all{ich wurde, du wurdest, wir wurden}.
  \item $\square$ Je place correctement le Partizip II à la fin de la phrase au passif.
  \item $\square$ Je distingue \all{von + Dativ} (agent) et \all{durch + Akkusativ} (moyen).
  \item $\square$ Je ne confonds pas passif prétérit (\all{wurde gebaut}) et Perfekt actif (\all{hat gebaut}).
  \item $\square$ Je sais construire une relative avec préposition : \all{die Mine, in der die Arbeiter arbeiten}.
  \item $\square$ Je mets le verbe conjugué en dernière position dans la subordonnée relative.
  \item $\square$ Je sais utiliser \all{wovon, worüber, womit} pour parler de choses.
  \item $\square$ Je n'oublie pas la majuscule à tous les noms allemands (\all{Gold, Mine, Arbeiter}).
  \item $\square$ Je sais résumer le texte en 3-4 phrases et donner mon avis sur les mines et l'environnement.
  \item $\square$ Je peux citer un exemple camerounais : \all{In Kamerun wurde Gold entdeckt.}
\end{itemize}
\end{aretenir}

\findecours

\end{document}
